%% ============================================================================
%%  N O T H I N G   O R I G I N A L   N E E D   S U R V I V E
%%
%%  Origin, Continuation, and the Geometry of Descent
%%
%%  Flyxion
%%  Independent Researcher
%%
%%  Compile with LuaLaTeX or XeLaTeX.
%%  Two passes are recommended. The document concerns histories of states;
%%  it would be discourteous for the document itself to pretend that one
%%  encounter with its own source is sufficient.
%% ============================================================================

\RequirePackage{iftex}

\ifPDFTeX
  \PackageError{originist-preamble}
    {This document requires LuaLaTeX or XeLaTeX}
    {pdfLaTeX cannot report the Unicode history of this document with
     sufficient candor. Compile with LuaLaTeX or XeLaTeX instead.}
\fi

\documentclass[
  11pt,
  a4paper,
  DIV=12,
  BCOR=0mm,
  parskip=half-
]{scrartcl}

% ============================================================================
%  ENGINE
% ============================================================================

\RequirePackage{ifluatex}
\RequirePackage{ifxetex}

\newcommand{\enginename}{unknown engine}

\ifluatex
  \renewcommand{\enginename}{LuaLaTeX}
\fi

\ifxetex
  \renewcommand{\enginename}{XeLaTeX}
\fi

% ============================================================================
%  TYPOGRAPHY
% ============================================================================

\usepackage{fontspec}
\usepackage{unicode-math}

\setmainfont{TeX Gyre Pagella}
\setsansfont{TeX Gyre Heros}[Scale=MatchLowercase]
\setmonofont{TeX Gyre Cursor}[Scale=MatchLowercase]
\setmathfont{TeX Gyre Pagella Math}

\usepackage{microtype}

\setlength{\emergencystretch}{2em}

% ============================================================================
%  MATHEMATICS
% ============================================================================

\usepackage{mathtools}
\usepackage{amsthm}
\usepackage{thmtools}
\usepackage{centernot}

\allowdisplaybreaks[2]

% ============================================================================
%  GENERAL DOCUMENT INFRASTRUCTURE
% ============================================================================

\usepackage{xcolor}
\usepackage{csquotes}
\usepackage{epigraph}
\usepackage{marginnote}
\usepackage{tikz-cd}
\usepackage[most]{tcolorbox}
\tcbuselibrary{skins,breakable}
\usepackage{scrlayer-scrpage}

% enumitem is deliberately available for technical appendices and unusual
% footnotes. Its presence is not permission to turn the running argument
% into an inventory.
\usepackage[shortlabels]{enumitem}

% ============================================================================
%  PROVENANCE PALETTE
% ============================================================================

\definecolor{verbatimink}{HTML}{14213D}
\definecolor{reconink}{HTML}{8A5A2E}
\definecolor{uncertaingray}{HTML}{6B6B6B}
\definecolor{woundred}{HTML}{7A1F2B}
\definecolor{invariantgreen}{HTML}{2E5339}

% ============================================================================
%  HYPERLINKS AND CROSS-REFERENCES
% ============================================================================

\usepackage[
  colorlinks=true,
  unicode=true
]{hyperref}

\hypersetup{
  linkcolor  = verbatimink,
  citecolor  = verbatimink,
  urlcolor   = reconink,
  pdftitle   = {Nothing Original Need Survive},
  pdfauthor  = {Flyxion},
  pdfsubject = {Origin, Continuation, and the Geometry of Descent},
  pdfcreator = {\enginename}
}

\usepackage[nameinlink,noabbrev]{cleveref}

% ============================================================================
%  PROVENANCE CLASSES
%
%  These commands describe epistemic status in the source. Their visual
%  appearance is secondary and may be altered without rewriting the
%  manuscript's provenance declarations.
% ============================================================================

\newcommand{\directstyle}[1]{%
  {\color{verbatimink}#1}%
}

\newcommand{\reconstyle}[1]{%
  {\itshape\color{reconink}#1}%
}

\newcommand{\uncertainstyle}[1]{%
  {\sffamily\color{uncertaingray}#1}%
}

\newcommand{\verbzero}[1]{\directstyle{#1}}
\newcommand{\reconstructed}[1]{\reconstyle{#1}}
\newcommand{\uncertain}[1]{\uncertainstyle{#1}}

% ============================================================================
%  PROVENANCE INDICES
%
%       \provn{0}
%
%  produces
%
%       [0]
%
%  as a superscript.
%
%       \provn[R\circ C\circ R]{3}
%
%  records both generation count and the known transformation chain.
%
%  The full form is deliberately cumbersome. If provenance becomes
%  cumbersome, the typography should not conceal that fact.
% ============================================================================

\newcommand{\provn}[2][]{%
  \if\relax\detokenize{#1}\relax
    \textsuperscript{$[#2]$}%
  \else
    \textsuperscript{$[#2:#1]$}%
  \fi
}

\newcommand{\provfull}[4]{%
  \textsuperscript{%
    $[#1:#2;\,#3;\,p=#4]$%
  }%
}

\newcommand{\provq}{%
  \textsuperscript{$[?]$}%
}

% ============================================================================
%  THE UNIQUE UNKNOWN-DEPTH WOUND
%
%  The manuscript permits exactly one occurrence of [\omega?].
%  If a second genuinely becomes necessary, the formal architecture has
%  changed and this constraint must be reconsidered explicitly.
% ============================================================================

\newcounter{woundtagcount}

\newcommand{\provomega}{%
  \stepcounter{woundtagcount}%
  \ifnum\value{woundtagcount}>1
    \PackageError{originist-preamble}
      {\string\provomega\space used more than once}
      {The [omega?] provenance class is reserved for one structural wound.
       A recurring wound is a convention. Reconsider the manuscript before
       weakening this constraint.}
  \fi
  \textsuperscript{$[\omega?]$}%
}

% Non-counting display form. For the appendix's own discussion OF the
% [\omega?] tag, and for any later prose that needs to point back at the
% wound without re-inflicting it. Does not touch \c{woundtagcount} and so
% never trips the scarcity guard above.
\newcommand{\provomegadisplay}{\textsuperscript{$[\omega?]$}}

% ============================================================================
%  RECONSTRUCTION LEDGER
%
%  This is an occurrence ledger, not a measure of genealogical depth.
%  Genealogical depth belongs to the provenance index. The ledger answers
%  only how often the manuscript has explicitly presented material as
%  reconstructed.
% ============================================================================

\newcounter{reconspanledger}

\newcommand{\recon}[1]{%
  \stepcounter{reconspanledger}%
  \reconstyle{#1}%
  \marginnote{%
    \tiny\color{reconink}%
    recon.\ \#\thereconspanledger
  }%
}

% ============================================================================
%  CORE RELATIONAL VOCABULARY
%
%  This vocabulary is intentionally small. New relations should appear only
%  when an example demonstrates that the existing relations are insufficient.
% ============================================================================

\newcommand{\state}[1]{x_{#1}}
\newcommand{\traj}[1]{\gamma_{#1}}
\newcommand{\graphG}{\mathcal G}
\newcommand{\Inv}{\mathcal I}
\newcommand{\Adm}{\mathcal A}
\newcommand{\World}{\mathcal W}

\newcommand{\eqbyte}{=_{\mathrm{byte}}}
\newcommand{\eqmaterial}{=_{\mathrm{material}}}
\newcommand{\eqprov}{=_{\mathrm{prov}}}

\newcommand{\simI}{\simeq_{\Inv}}
\newcommand{\simgen}{\sim_{\mathrm{gen}}}
\newcommand{\simconstraint}{\sim_{\mathcal C}}

\newcommand{\leadcausal}{%
  \mathrel{\rightsquigarrow_{\mathrm{causal}}}%
}

\newcommand{\eqworld}{%
  \equiv_{\mathrm{world}}%
}

\newcommand{\lossorig}{%
  \mathscr L_{\mathrm{orig}}%
}

\newcommand{\Truth}{\operatorname{Truth}}
\newcommand{\Prov}{\operatorname{Prov}}
\newcommand{\Depth}{\operatorname{Depth}}
\newcommand{\Hist}{\operatorname{Hist}}
\newcommand{\End}{\operatorname{End}}
\newcommand{\Orig}{\operatorname{Orig}}
\newcommand{\Cont}{\operatorname{Cont}}
\newcommand{\Hom}{\operatorname{Hom}}
\newcommand{\Mor}{\operatorname{Mor}}
\newcommand{\State}{\operatorname{State}}
\newcommand{\Trajectory}{\operatorname{Trajectory}}
\newcommand{\Provenance}{\operatorname{Provenance}}
\newcommand{\Admissibility}{\operatorname{Admissibility}}
\newcommand{\Worldhood}{\operatorname{Worldhood}}
\newcommand{\History}{\operatorname{History}}
\newcommand{\Content}{\operatorname{Content}}
\newcommand{\Pop}{\operatorname{Pop}}
\newcommand{\Refuse}{\operatorname{Refuse}}
\newcommand{\Bind}{\operatorname{Bind}}
\newcommand{\Collapse}{\operatorname{Collapse}}

\newcommand{\nimplies}{%
  \centernot\Rightarrow
}

% ============================================================================
%  TRAJECTORIES
%
%  The first form is deliberately elementary. More elaborate provenance
%  graphs are introduced only after linear trajectories have failed.
% ============================================================================

\newcommand{\trajectory}[4]{%
  #1
  \xrightarrow{#2}
  #3
  \longrightarrow
  #4%
}

\newcommand{\pathclass}[2]{%
  \left[#1\right]_{#2}%
}

\newcommand{\quotient}[2]{%
  #1/{#2}%
}

% ============================================================================
%  WOUNDS
%
%  A wound is not a red box. A wound is an irreversible loss of a distinction.
%  The box merely makes certain wounds visible.
% ============================================================================

\newcommand{\Wound}[2]{%
  \lossorig\!\left(#1\,;\,#2\right)%
}

\newtcolorbox{woundbox}{
  enhanced,
  breakable,
  colback=white,
  colframe=woundred,
  boxrule=0.6pt,
  arc=0pt,
  left=7pt,
  right=7pt,
  top=6pt,
  bottom=6pt,
  borderline={0.6pt}{0pt}{woundred,dashed}
}

% ============================================================================
%  HELD PROPOSITIONS
% ============================================================================

\newtcolorbox{held}[1][]{
  enhanced,
  breakable,
  colback=white,
  colframe=verbatimink,
  boxrule=0.45pt,
  arc=0pt,
  left=7pt,
  right=7pt,
  top=6pt,
  bottom=6pt,
  fonttitle=\normalfont\scshape\small,
  title={#1}
}

\newcommand{\statehistory}{%
  \begin{held}
    \centering
    \Large
    $\textsc{State}\neq\textsc{History}$
  \end{held}
}

% ============================================================================
%  THEOREM-LIKE ENVIRONMENTS
%
%  Conventional mathematical taxonomy is retained where useful, but the
%  manuscript also names the forms of reasoning peculiar to its argument.
% ============================================================================

\theoremstyle{definition}

\newtheorem{definition}{Definition}[section]
\newtheorem{pathology}[definition]{Pathology}
\newtheorem{objection}[definition]{Originist Objection}
\newtheorem{crossing}[definition]{Boundary Crossing}

\theoremstyle{plain}

\newtheorem{proposition}[definition]{Proposition}
\newtheorem{theorem}[definition]{Theorem}
\newtheorem{lemma}[definition]{Lemma}
\newtheorem{corollary}[definition]{Corollary}

\theoremstyle{remark}

\newtheorem{remark}[definition]{Remark}
\newtheorem{example}[definition]{Example}

\newtheorem*{wounded}{The Wound}

\crefname{pathology}{pathology}{pathologies}
\Crefname{pathology}{Pathology}{Pathologies}

\crefname{objection}{Originist objection}{Originist objections}
\Crefname{objection}{Originist Objection}{Originist Objections}

\crefname{crossing}{boundary crossing}{boundary crossings}
\Crefname{crossing}{Boundary Crossing}{Boundary Crossings}

% ============================================================================
%  TRAJECTORY DIAGRAMS
% ============================================================================

\newenvironment{trajectorydiagram}
  {%
    \begin{center}
    \begin{tikzcd}[
      row sep=1.6em,
      column sep=2.8em
    ]
  }
  {%
    \end{tikzcd}
    \end{center}
  }

% ============================================================================
%  MOVEMENTS
%
%  The document calls its principal divisions Movements rather than Sections.
%  This is not intended to imply that their sequence is reversible.
% ============================================================================

\renewcommand{\sectionformat}{%
  \textsc{Movement \Roman{section}}\autodot\enskip
}

\setkomafont{section}{%
  \normalfont\Large\scshape
}

\setkomafont{subsection}{%
  \normalfont\large\itshape
}

\setkomafont{subsubsection}{%
  \normalfont\normalsize\scshape
}

\newcommand{\movementquestion}[1]{%
  \begin{center}
    \small\itshape
    #1
  \end{center}
  \medskip
}

% ============================================================================
%  PREFATORY REGISTER
%
%  Every Movement begins in ordinary language. Formal notation follows only
%  after the question has been made intelligible without it.
% ============================================================================

\newenvironment{orientation}
  {%
    \begin{quote}
    \small
    \setlength{\parindent}{1em}
  }
  {%
    \end{quote}
  }

% ============================================================================
%  EPIGRAPHS
% ============================================================================

\setlength{\epigraphwidth}{0.72\textwidth}
\setlength{\epigraphrule}{0pt}

\newcommand{\movementepigraph}[2]{%
  \epigraph{\emph{#1}}{#2}%
}

% ============================================================================
%  RUNNING HEADS
% ============================================================================

\pagestyle{scrheadings}
\clearpairofpagestyles

\ihead{%
  \small\scshape
  Nothing Original Need Survive
}

\ohead{%
  \small
  \pagemark
}

\ifoot{%
  \scriptsize
  reconstructed spans encountered:
  \thereconspanledger
}

\ofoot{%
  \scriptsize
  \enginename
}

% ============================================================================
%  FOOTNOTES
%
%  Ordinary footnotes remain ordinary. Provenance-bearing notes are explicit.
%  The manuscript does not redefine \footnote globally because editorial
%  annotation and provenance annotation are not the same operation.
% ============================================================================

\newcommand{\pfootnote}[2][?]{%
  \footnote{%
    \textsuperscript{$[#1]$}\,
    #2%
  }%
}

% ============================================================================
%  APPENDIX NAMES
% ============================================================================

\renewcommand{\appendixname}{Technical Appendix}

% ============================================================================
%  TITLE
% ============================================================================

\title{%
  \bfseries
  Nothing Original Need Survive
}

\subtitle{%
  Origin, Continuation, and the Geometry of Descent\\[4pt]
  \normalfont\small
  Card's Originism, Zebrowski's \emph{Macrolife},
  Le Guin's \emph{The Dispossessed}, and the Problem of
  Historical Continuity
}

\author{%
  Flyxion\\
  \normalfont\small
  Independent Researcher
}

\date{%
  July 2026\\[2pt]
  \footnotesize
  compiled under \enginename\\[2pt]
  \itshape
  a generation-zero build is not a guarantee
}

% ============================================================================
%  DOCUMENT
% ============================================================================

\begin{document}

\maketitle

\begin{abstract}
What makes a later thing a continuation of an earlier one when the material,
location, representation, and even the recoverable record of the earlier thing
may all have changed? Conversely, what can a present record establish about an
origin that no longer survives? These questions are often separated into
problems of historical knowledge, personal or collective identity, archival
reconstruction, and persistence through change. This essay argues that they
share a common structural difficulty. States do not, in general, contain enough
information to determine the histories by which they arose.

Three works of speculative fiction provide unusually severe probes of this
difficulty. Orson Scott Card's ``The Originist'' approaches it retrospectively:
a present assembled from incomplete records is used to reason toward an absent
origin. George Zebrowski's \emph{Macrolife} approaches it prospectively:
civilization persists by becoming increasingly independent of its original
material and planetary conditions. Ursula K.\ Le Guin's
\emph{The Dispossessed} introduces a third operation, return, in which departure
and recurrence reveal that neither an origin nor the agent returning to it
remains simply identical with its earlier state.

The resulting argument moves from states to trajectories, from trajectories
to provenance structures, and from provenance to constraints on admissible
continuation. These distinctions are then related to the author's existing
work on the Relativistic Scalar--Vector Plenum (RSVP), TARTAN, Chain of Memory,
Spherepop, irreversibility, and the No-World argument. The governing claim is
simple but its consequences are not: continuity is not adequately represented
as a relation between states. It is a property of histories, and histories
themselves may branch, converge, become partially unrecoverable, or cease to
support the continuation of a world even where some descendant structure
survives.

The document's citation and typographic apparatus participates in this
argument. Claims may carry provenance indices recording direct inspection,
known reconstruction depth, or unresolved derivational history. These indices
do not represent truth values. They record a different coordinate of knowledge:
how a statement reached the present document.
\end{abstract}

\clearpage

% ============================================================================
%  PREFACE
% ============================================================================

\section*{Preface: What It Means for Something to Continue}
\addcontentsline{toc}{section}{Preface: What It Means for Something to Continue}

Most questions about whether something has survived appear easy until the
survival is imperfect. A book remains the same book when it is moved from one
shelf to another. A city remains the same city after a building is demolished.
A person remains the same person after replacing the cells from which much of
the body was once composed. A language remains recognizable despite the death
of every speaker who spoke it several centuries ago. In ordinary life these
judgments rarely require a theory. We inherit conventions about which changes
matter and which do not, and those conventions usually work well enough.

Difficulty begins when several kinds of continuity separate.

A copied file may contain exactly the same bytes as another file while having
a completely different history. A damaged manuscript may be reconstructed so
successfully that the reconstructed passage is indistinguishable from the lost
one, even though the passage now before us was produced by an editor rather
than transmitted continuously from the earlier document. A ship may preserve
its name, function, crew, route, and institutional identity while every piece
of its material is gradually replaced. A civilization may survive the
destruction of its home planet while changing its bodies, habitats,
technologies, and forms of social organization. Conversely, a civilization may
retain buildings, names, and institutions while losing so much of the
relational structure that once constituted its world that calling the result
the same civilization begins to conceal more than it explains.

The ordinary word \enquote{same} performs too many jobs in such cases.

This essay begins from the suspicion that many apparent puzzles of identity
are produced by asking states to answer questions that only histories can
answer. If two objects presently look alike, their present similarity does not
tell us whether one descended from the other, whether both descended from a
common source, whether they were produced independently, or whether one was
constructed specifically to imitate the other. Likewise, present difference
does not establish historical discontinuity. A descendant may differ radically
from its ancestor precisely because an unbroken process of adaptation connected
them.

The distinction is elementary in ordinary experience. Finding two identical
handwritten notes on a table does not establish that they were written by the
same person. Discovering that one note was copied from the other changes what
we know about their relationship without changing either note. Discovering
instead that both were independently copied from a third note changes the
relationship again. Nothing about the visible sentences needs to change while
our account of what those sentences are evidence of changes substantially.

The same difficulty appears at much larger scales.

In Card's ``The Originist,'' the problem points backward. Surviving evidence
must support claims about an origin that is no longer directly available.
Originism therefore concerns more than the storage of records. It concerns the
relations among records and the histories through which present evidence became
available.

In Zebrowski's \emph{Macrolife}, the problem points forward. Survival ceases to
mean preserving a fixed material arrangement and increasingly means maintaining
a viable continuation through transformations that may eventually leave little
of the original substrate intact.

Le Guin's \emph{The Dispossessed} complicates both directions. Departure is
followed by return, but return does not restore an earlier state. The traveler
has changed, the place of return has continued to develop, and the origin
itself turns out to have an origin. The apparently simple line connecting
departure and return therefore opens into a history of boundaries,
inheritances, institutional transformations, and branching transmissions.

These literary works will not be treated as authorities from which an ontology
can simply be extracted. They are useful because they produce difficult cases.
Each forces a distinction that a simpler representation would prefer not to
make. Their role in the argument is therefore closer to the role of a
counterexample in mathematics than to the role of evidence in literary
history.

The formal apparatus will be introduced in the same manner. No distinction
will be retained merely because it permits an attractive symbol. A new
relation earns its place when two cases that must remain distinguishable would
otherwise collapse into one. Conversely, when two histories can safely be
treated as equivalent for a particular purpose, the equivalence will be stated
with respect to that purpose rather than elevated silently into identity.

This last qualification is important. Mathematics obtains much of its power
from forgetting differences. Objects are identified up to isomorphism,
homeomorphism, gauge transformation, or some other equivalence precisely
because the discarded differences are irrelevant to the question under study.
Nothing in the present argument opposes abstraction as such. The concern is
narrower: a distinction should not be discarded before one has established
that the operations to be performed later do not require it.

The central danger is therefore not forgetting.

It is forgetting that one has forgotten.

This is the point at which provenance becomes more than bibliographic
bookkeeping. Provenance records which distinctions have survived a sequence of
transformations and which have not. A claim may remain true while its
provenance becomes uncertain. A reconstruction may reproduce an original
perfectly while remaining a reconstruction. A descendant may preserve an
invariant without preserving the path by which that invariant was transmitted.
A world may leave traces from which some of its structure can be reconstructed
without thereby becoming recoverable as the world that produced those traces.

The argument that follows will therefore resist a premature declaration of
what its fundamental object must be. It may prove sufficient to speak of
trajectories. Branching histories may require graphs. Questions of admissible
continuation may require trajectories together with constraints. Questions of
worldhood may force another extension. The stopping condition is not that every
possible distinction has been represented. Such a condition could never be
met. The argument closes when the examples under consideration cease to force
an enlargement of the kind of object required to represent them.

Closure, in other words, will itself be treated historically.

\clearpage

% ============================================================================
%  NOTE ON PROVENANCE
% ============================================================================

\section*{A Note on Provenance}
\addcontentsline{toc}{section}{A Note on Provenance}

The superscripts used in this essay should not be read as confidence scores.

A claim marked \provn{0} indicates that, for the purpose for which the claim is
being used, the relevant wording or evidence has been inspected directly in
the source identified by the accompanying citation. A claim marked
\provq\ indicates that some relevant part of its derivational history has not
yet been established. Further distinctions will be introduced only when the
argument requires them.

Thus the notation

\[
  \varphi^{[0]}
\]

does not mean that \(\varphi\) is certainly true, while

\[
  \psi^{[?]}
\]

does not mean that \(\psi\) is probably false. Directly inspected sources can
contain errors. Claims of uncertain provenance can be true. The superscript
records a relation between the present statement and the evidence through which
it entered the argument.

The distinction can be expressed provisionally as

\[
  \Truth(\varphi)
  \neq
  \Prov(\varphi).
\]

At this stage, however, the notation should be understood in the ordinary
language just given rather than treated as a developed formal system. One of
the purposes of the essay is to determine how much structure such a system
actually requires.

The document also distinguishes ordinary editorial footnotes from
provenance-bearing notes. A footnote may clarify a term, supply historical
context, or acknowledge a complication without making a claim about the
derivational depth of the sentence to which it is attached. Provenance is
therefore not silently identified with citation.

The distinction will matter later.

\clearpage

% ============================================================================
%  CONTENTS
% ============================================================================

\tableofcontents

\clearpage

% ============================================================================
%  MOVEMENT I
%  ORIGIN
% ============================================================================

\section{Origin}
\label{mov:origin}

\movementquestion{Can the past be recovered from its traces?}

\begin{orientation}
There is an ordinary sense in which finding an origin means finding a place.
A family comes from a town, a manuscript comes from an archive, a species
comes from an ancestral population, and a civilization comes from somewhere.
If the place can be identified, the problem appears to be solved.

But origins become more difficult when the thing that began no longer exists
in a recoverable form. Suppose the first settlement has disappeared, its
records have been copied and recopied, its descendants have migrated, and the
stories told about it have been reconstructed from fragments. Finding a
plausible location would then answer only one question among many. We would
still want to know which present traces actually descend from the lost
settlement, which merely resemble its traces, which have been reconstructed,
and which principles of organization persisted after the original material
vanished.

Orson Scott Card's ``The Originist'' is useful because it begins inside this
difficulty. Its subject is ostensibly the origin of humanity, but its deeper
problem is the status of evidence after the origin has receded beyond direct
inspection. The originist cannot simply look backward. The originist must
reason through what survived.
\end{orientation}

% -----------------------------------------------------------------------------
%  1.1  THE ORIGINIST'S PROBLEM
% -----------------------------------------------------------------------------

\subsection{The Originist's Problem}
\label{sec:originist-problem}

Orson Scott Card's ``The Originist,'' first published in the anthology
\emph{Foundation's Friends}, places Leyel Forska within Isaac Asimov's
Foundation universe during the period surrounding Hari Seldon's work on
psychohistory \cite{card1989}. The story inherits the scale of Asimov's
Galactic Empire while directing attention toward a comparatively narrow
scholarly problem. Leyel is concerned with human origins: not merely with the
location from which humanity began its expansion, but with the reconstruction
of the conditions and processes under which recognizably human civilization
emerged.

The distinction matters. A search for a planet and a search for an origin are
not necessarily the same search.

If the problem were merely geographical, originism could in principle
terminate at a coordinate. Let \(E\) denote Earth and let \(p(E)\) denote its
location. The successful recovery of \(p(E)\) would then settle the relevant
question. Yet a coordinate says almost nothing about why the civilization
whose descendants now occupy the Galaxy developed there, which inherited
structures distinguish genuine descent from later resemblance, or which
features of contemporary humanity carry information about the conditions of
its emergence.

The geographical question has the schematic form

\[
  \text{Where was } O?
\]

whereas the originist's question is closer to

\[
  \text{What can the present establish about } O?
\]

These questions coincide only under unusually favorable archival conditions.

Card's originism is therefore valuable here not because the fictional
discipline supplies a ready-made philosophy of provenance. It does not.
Rather, the story dramatizes a circumstance in which an origin has become
epistemically remote while remaining causally indispensable. Humanity must
have had an earlier history. The loss of direct access to that history does
not remove it from the causal ancestry of the present. What disappears is
something more specific: the ability to move without difficulty from present
evidence to a unique account of how the present arose.

That is the first asymmetry on which this essay depends.

The past can determine the present without the present uniquely determining
the past.

% -----------------------------------------------------------------------------
%  1.2  THE LIBRARY AND THE SCHOLAR
% -----------------------------------------------------------------------------

\subsection{The Library and the Scholar}
\label{sec:library-scholar}

The Imperial Library gives this asymmetry an institutional form. A library
preserves records. An originist does something different with them.

This distinction can initially be stated without mathematical machinery.
Imagine that an archive possesses a document copied two thousand years ago
from another document that no longer survives. The archive's task is
satisfied, in one important sense, by preserving the surviving copy. It can
record where the copy is stored, when it entered the collection, what material
it is written on, and perhaps what earlier catalogues said about it.

None of this by itself establishes the entire history of the text.

The scholar asks questions that the shelf cannot answer. Was the surviving
document copied faithfully? Was its source itself a copy? Were passages
silently repaired? Did two apparently independent documents derive from a
common intermediary? Does a sentence preserve ancient wording, or is it a
later reconstruction that happened to become canonical?

The difference is sufficiently important to give the two functions names.

\begin{definition}[Archivist]
\label{def:archivist}
An \emph{Archivist} is an agent or process whose primary operation is the
preservation and retrieval of available states together with whatever
explicit metadata accompanies them.
\end{definition}

\begin{definition}[Originist]
\label{def:originist}
An \emph{Originist} is an agent or process whose primary operation is
reasoning about the derivational histories by which available states may have
arisen, including distinctions that are not recoverable from the states
considered in isolation.
\end{definition}

The capitalization is deliberate. These definitions do not purport to define
Card's fictional profession exhaustively. They name two computational roles
that the literary case makes unusually easy to separate.

The Archivist asks, in effect,

\[
  \text{What survives?}
\]

The Originist asks

\[
  \text{How did what survives get here?}
\]

An ideal institution may perform both operations. Indeed, Card's setting is
interesting precisely because the Library and the scholarly practice of
originism place them in close proximity. But proximity should not be mistaken
for identity.

A perfectly preserved state can have an uncertain history.

Conversely, a badly damaged state can sometimes belong to a well-established
history.

The amount of surviving content and the quality of surviving provenance are
therefore distinct quantities.

% -----------------------------------------------------------------------------
%  1.3  FROM STATES TO TRAJECTORIES
% -----------------------------------------------------------------------------

\subsection{From States to Trajectories}
\label{sec:states-trajectories}

The preceding distinction suggests the first formal enlargement of the object
under study.

Let \(X\) be a space of possible states. For the moment, nothing substantial
depends on what kind of states these are. They may be documents, organisms,
institutional configurations, software artifacts, civilizations, or
representations within a computational system.

A state is written

\[
  x\in X.
\]

If only the present state matters, this notation may be sufficient. The
Originist's problem begins when it is not.

Suppose that \(x_0\) undergoes a sequence of transformations

\[
  T_1,T_2,\ldots,T_n
\]

producing states

\[
  x_0,x_1,\ldots,x_n.
\]

We define the corresponding trajectory as follows.

\begin{definition}[Derivational trajectory]
\label{def:trajectory}
A \emph{derivational trajectory} terminating in \(x_n\) is an ordered sequence

\[
  \traj{x_n}
  =
  \left(
    x_0
    \xrightarrow{T_1}
    x_1
    \xrightarrow{T_2}
    \cdots
    \xrightarrow{T_n}
    x_n
  \right),
\]

where each \(T_i\) records a transformation relevant to the provenance
question under consideration.
\end{definition}

The phrase \enquote{relevant to the provenance question} is necessary. A
trajectory containing every physical event in the causal history of a
manuscript would be useless. The temperature of the room in which a copyist
worked might belong to the total causal history of a document without
belonging to the derivational history of its wording.

A trajectory is therefore already selective.

This fact will become important later, because selecting which transformations
count is itself an operation capable of losing information. For the present,
however, it is enough to notice that the endpoint map

\[
  \End(\traj{x_n})=x_n
\]

forgets almost everything about the trajectory.

Many different histories may possess the same endpoint.

Formally, if \(\Gamma(X)\) denotes the relevant space of trajectories in
\(X\), then

\[
  \End:\Gamma(X)\longrightarrow X
\]

need not be injective.

This elementary observation contains the first formal version of the
Originist's problem.

\begin{proposition}[Endpoint insufficiency]
\label{prop:endpoint-insufficiency}
In any provenance system for which the endpoint map
\(\End:\Gamma(X)\rightarrow X\) is non-injective, the present state is
insufficient to determine its derivational history.
\end{proposition}

\begin{proof}
Non-injectivity means that there exist distinct trajectories
\(\gamma_A,\gamma_B\in\Gamma(X)\) such that

\[
  \gamma_A\neq\gamma_B
\]

while

\[
  \End(\gamma_A)=\End(\gamma_B).
\]

Let their common endpoint be \(x\). Observation of \(x\) alone therefore
cannot determine whether the derivational history was \(\gamma_A\) or
\(\gamma_B\). Hence the endpoint is insufficient to determine derivational
history.
\end{proof}

The proof is trivial. Its consequences are not.

% -----------------------------------------------------------------------------
%  1.4  CONVERGENT RECONSTRUCTION
% -----------------------------------------------------------------------------

\subsection{Convergent Reconstruction}
\label{sec:convergent-reconstruction}

Consider an original state \(O\).

One surviving object, \(x_A\), is copied directly from it:

\[
  \gamma_A:
  \qquad
  O\longrightarrow x_A.
\]

A second object passes through repeated reconstruction and compression:

\[
  \gamma_B:
  \qquad
  O
  \longrightarrow R_1
  \longrightarrow C_1
  \longrightarrow R_2
  \longrightarrow C_2
  \longrightarrow x_B.
\]

Suppose now that the final states are indistinguishable under the strongest
state-level comparison available to us. For a digital artifact, we may take
the extreme case:

\[
  x_A \eqbyte x_B.
\]

Nevertheless,

\[
  \gamma_A\neq\gamma_B.
\]

The equality of endpoints has not collapsed the difference between their
histories.

This is the first pathology around which the rest of the essay will grow.

\begin{pathology}[Convergent reconstruction]
\label{path:convergent-reconstruction}
Two derivationally distinct trajectories may terminate in states that are
identical under a chosen state-level equivalence relation:

\[
  \gamma_A\neq\gamma_B,
  \qquad
  \End(\gamma_A)\sim\End(\gamma_B).
\]

Consequently, endpoint equivalence does not in general imply provenance
equivalence.
\end{pathology}

In the byte-identical case,

\[
  x_A\eqbyte x_B
  \nimplies
  x_A\eqprov x_B.
\]

The notation is intentionally stronger than the example requires. Nothing
about the argument depends upon digital files. Byte identity merely makes the
counterexample difficult to evade because it removes ordinary appeals to
visible difference.

The two artifacts can be identical now.

They can still have different histories.

% -----------------------------------------------------------------------------
%  1.5  GENERATION
% -----------------------------------------------------------------------------

\subsection{Generation Without Difference}
\label{sec:generation}

The distinction becomes particularly important when transformations are
regenerative.

Suppose an artifact is summarized, regenerated from the summary, summarized
again, and regenerated again. The resulting artifact may contain sentences
identical to sentences in the original. A comparison system that evaluates
only present content will naturally assign no penalty to an exact match.

An Originist may nevertheless wish to distinguish

\[
  \varphi^{[0]}
\]

from

\[
  \varphi^{[1]}.
\]

At this stage the superscripts should be read minimally. The first records
directly inspected material. The second records material separated from the
relevant source by one recognized derivational transformation.

The important point is not that

\[
  \varphi^{[0]}
\]

is somehow semantically superior to

\[
  \varphi^{[1]}.
\]

They may express precisely the same proposition. They may even contain
precisely the same characters.

Rather,

\[
  \Prov\!\left(\varphi^{[0]}\right)
  \neq
  \Prov\!\left(\varphi^{[1]}\right).
\]

Generation is therefore not a property of visible content considered alone.
It is a property of content situated in a derivational history.

This gives a more precise interpretation to the computational role called
Originist.

\begin{definition}[Generation depth]
\label{def:generation-depth}
Given a designated source state \(O\) and a derivational trajectory

\[
  \gamma:
  O
  \xrightarrow{T_1}
  x_1
  \xrightarrow{T_2}
  \cdots
  \xrightarrow{T_n}
  x_n,
\]

the \emph{generation depth} of \(x_n\) relative to \(O\), under the chosen
transformation vocabulary, is

\[
  \Depth_O(x_n;\gamma)=n.
\]
\end{definition}

The dependence on \(\gamma\) is essential.

If two histories converge on the same endpoint, then it is possible that

\[
  x_A=x_B
\]

while

\[
  \Depth_O(x_A;\gamma_A)
  \neq
  \Depth_O(x_B;\gamma_B).
\]

Thus there is no well-defined function

\[
  \Depth_O:X\longrightarrow\mathbb N
\]

unless additional assumptions guarantee a unique relevant derivational
history for each state.

Generation belongs first to a trajectory, not to an endpoint.

% -----------------------------------------------------------------------------
%  1.6  APPROXIMATION WITHOUT ANCESTRY
% -----------------------------------------------------------------------------

\subsection{Approximation Without Ancestry: Shannon's Markov Languages}
\label{sec:shannon-markov}

The trajectory-relative reading of generation just established has an
unusually early and clean illustration, from outside the manuscript's usual
sources.

Claude Shannon's 1948 paper on communication theory constructs successive
statistical approximations to English by conditioning each generated symbol
on an increasing span of preceding context \cite{shannon1948}. Schematically,
an order-\(k\) approximation sets

\[
  P(w_{t+1}\mid w_t,w_{t-1},\ldots)
  =
  P(w_{t+1}\mid w_t,\ldots,w_{t-k+1}),
\]

and generates symbols accordingly. As \(k\) increases, or as the conditioning
unit moves from letters to words to word pairs, the generated sequences
acquire progressively stronger surface resemblance to ordinary English,
despite arising from a generative mechanism with no author, no intention, and
no historical relation to any actual English utterance.

\begin{definition}[Order-\(k\) statistical generation]
\label{def:order-k-generation}
A sequence \(m=(w_1,\ldots,w_n)\) is an \emph{order-\(k\) statistical
generation} over a source distribution if each \(w_{t+1}\) is drawn according
to \(P(w_{t+1}\mid w_t,\ldots,w_{t-k+1})\) for some fixed \(k\), with no
further dependence on \(t\) or on any external record.
\end{definition}

Shannon's own examples already anticipate the interesting case. He reports
that as \(k\) grows the generated samples resemble ordinary English more
closely, and that they exhibit apparent structure extending beyond what the
order-\(k\) dependency explicitly supplies \cite{shannon1948}. That second
observation matters more than the first for this manuscript. It is not
merely that Markov generation can look like English. It is that a reader
examining the output can be led to attribute a longer-range organizing
history to it than the generator actually possessed.

\begin{pathology}[Statistical resemblance without genealogy]
\label{path:shannon-resemblance}
Let \(h\) be a state produced by ordinary human linguistic production, and
let \(m\) be an order-\(k\) statistical generation for some fixed \(k\).
Suppose a selected observer or discriminator \(\mathcal S\) finds the two
increasingly similar as \(k\) grows:

\[
  h\sim_{\mathcal S}m.
\]

Then nothing follows about the relation between their generative histories.
In particular,

\[
  h\sim_{\mathcal S}m
  \;\nimplies\;
  h\eqprov m.
\]
\end{pathology}

\begin{proof}
This is an instance of \cref{path:convergent-reconstruction} in which one of
the two trajectories has no originating state at all in the relevant sense:
\(h\) descends from an ordinary history of authorship, while \(m\) descends
from a local statistical generator with no comparable history to disclose.
Endpoint resemblance under \(\mathcal S\) is, by hypothesis, a fact about
\(\End(\gamma_h)\) and \(\End(\gamma_m)\) alone; provenance equivalence is a
fact about \(\gamma_h\) and \(\gamma_m\) as trajectories. The two need not
agree, for the reason already established in
\cref{path:convergent-reconstruction}.
\end{proof}

What makes the Shannon case sharper than an ordinary instance of convergent
reconstruction is the direction of the failure. \Cref{sec:generation} showed
that generation depth cannot be read off an endpoint without specifying a
trajectory. Order-\(k\) generation shows something stronger: an observer can
perceive organization at a range the generator never implemented. The
failure is not merely that history is invisible at the endpoint; it is that
the endpoint can suggest a richer history than the one that actually
produced it.

\begin{held}[What Shannon supplies]
\centering
Statistical structure can outrun the dependency structure explicitly
encoded by its generator.
\end{held}

This is a narrower and more defensible claim than it may first appear.
Shannon's framework brackets meaning from the engineering problem of
communication throughout, and nothing here should be read as claiming that
order-\(k\) generation produces, or fails to produce, meaning
\cite{shannon1948}. The claim is only about genealogy, not semantics:

\[
  \boxed{
    \text{linguistic appearance}
    \;\nimplies\;
    \text{linguistic genealogy.}
  }
\]

An artifact can enter the region of state space occupied by ordinary
linguistic production without traversing any of the histories by which that
region is ordinarily reached. \Cref{sec:prompt-injection-authority} returns
to this once representations are permitted to act, rather than merely to
resemble.

% -----------------------------------------------------------------------------
%  1.7  ORIGINISM AS AN INVERSE PROBLEM
% -----------------------------------------------------------------------------

\subsection{Originism as an Inverse Problem}
\label{sec:inverse-problem}

The preceding formalism also clarifies why originism should not be imagined as
simple reversal.

Suppose a forward process is known:

\[
  O
  \longrightarrow
  x_1
  \longrightarrow
  x_2
  \longrightarrow
  P.
\]

If every transformation were invertible and perfectly recorded, recovering
\(O\) from \(P\) would be straightforward in principle. One could compose the
inverse transformations.

Historical reconstruction is interesting precisely because this condition
usually fails.

Compression may discard distinctions. Copying may omit metadata. Translation
may preserve one structure while altering another. Independent traditions may
merge. Records may disappear. Later agents may reconstruct missing material.
A present state can therefore be compatible with more than one possible past.

Let

\[
  \Gamma^{-}(P)
\]

denote the set of admissible trajectories terminating at the present state
\(P\). Originist inference seeks not necessarily a unique member of this set,
but constraints upon it.

Schematically,

\[
  P
  \quad\rightsquigarrow\quad
  \Gamma^{-}(P).
\]

The arrow here is epistemic rather than historical. The actual causal
trajectory ran toward \(P\). Inference runs over candidate descriptions in
the opposite direction.

This distinction prevents a common mistake. Reconstructing a plausible past
does not make the reconstruction causally ancestral to the evidence from
which it was inferred.

The map and the history point in opposite directions.

% -----------------------------------------------------------------------------
%  1.8  THE ORIGINIST OBJECTION
% -----------------------------------------------------------------------------

\subsection{The First Originist Objection}
\label{sec:first-originist-objection}

Mathematical reasoning routinely simplifies objects by identifying states
that differ only in ways irrelevant to the problem at hand. Such
identifications are indispensable. Without them, abstraction would be
impossible.

Originism does not object to equivalence.

It objects to equivalence without an account of what the equivalence forgets.

Suppose

\[
  q:X\longrightarrow X/{\sim}
\]

identifies states under an equivalence relation \(\sim\). If two states
possess distinguishable provenance histories but are mapped to the same
equivalence class, then the quotient has erased a distinction available
before the operation.

Whether that erasure matters depends upon what one intends to do next.

\begin{objection}[Preliminary form]
\label{obj:preliminary-originist}
An equivalence relation is epistemically dangerous for an Originist task when
it identifies states whose distinctions are required to discriminate among
admissible derivational histories.
\end{objection}

The qualification \enquote{for an Originist task} prevents the objection from
becoming a general hostility toward quotienting. Byte identity may be exactly
the correct equivalence for verifying a file transfer. It may be exactly the
wrong equivalence for determining whether two files possess the same
genealogy.

The question is therefore never merely

\[
  x\sim y?
\]

It is also

\[
  \text{What future inference becomes impossible after treating }
  x\text{ and }y\text{ as the same?}
\]

The formal loss operator associated with this question will not be introduced
yet. The examples required to define it adequately have not occurred.

% -----------------------------------------------------------------------------
%  1.9  THE ABSENT GROUND
% -----------------------------------------------------------------------------

\subsection{The Absent Ground}
\label{sec:absent-ground}

Card's originists work under a condition that can now be stated more
precisely. The origin is causally prior but epistemically unavailable.
Present evidence is not the origin. Nor is a successful reconstruction
identical with the origin merely because it reproduces some of its structure.

This produces three objects that ordinary historical language often slides
between:

\[
  O,
  \qquad
  \gamma,
  \qquad
  \widehat O.
\]

Here \(O\) is the historical origin, \(\gamma\) is a trajectory connecting
that origin to surviving evidence, and \(\widehat O\) is a reconstruction
inferred from that evidence.

Even in the best case,

\[
  \widehat O\simeq O
\]

requires an explicit account of what the equivalence symbol means.

Structural agreement is not historical identity.

Semantic agreement is not material identity.

Successful reconstruction is not temporal reversal.

An origin may therefore become unrecoverable without becoming irrelevant.
Its effects remain distributed through descendants, records, constraints, and
recurrences. Originism begins precisely where possession of the origin ends.

One early account of the fictional discipline describes Leyel's research as
oriented toward general principles of human origins rather than merely toward
the identification of a single planetary point of departure
\cite{card1989}.\pfootnote[0]{The distinction is retained here because
it motivates the separation between locating an origin and recovering the
constraints under which an origin can be inferred.}

This shift from point-location to constraint is the deepest reason Card's
fictional discipline belongs in the present argument. The Originist does not
need to possess the ground in order for the ground to constrain what may
truthfully be said about the present.

% -----------------------------------------------------------------------------
%  1.10  WHAT MOVEMENT I HAS FORCED
% -----------------------------------------------------------------------------

\subsection{What the Origin Has Already Failed to Tell Us}
\label{sec:origin-forced}

We began with a state.

The literary problem forced us to add its trajectory.

That enlargement was not motivated by a preference for process philosophy or
by a prior commitment to a particular metaphysics. It was forced by a simple
counterexample:

\[
  x_A=x_B,
  \qquad
  \gamma_A\neq\gamma_B.
\]

Once this case is admitted, no theory that represents only the endpoint can
recover every distinction relevant to provenance.

The conclusion should nevertheless remain modest. Nothing in this movement
establishes that trajectories are the fundamental objects of reality. It
establishes only that states are insufficient for the particular class of
questions under consideration.

That difference will govern the rest of the essay.

\begin{held}[Endpoint Insufficiency]
\centering
A present state may preserve the result of a history\\
without preserving enough information to identify that history.
\end{held}

Card has therefore supplied the backward-facing problem:

\[
  \boxed{
  \text{surviving traces}
  \;\rightsquigarrow\;
  \{\text{admissible pasts}\}
  }.
\]

But nothing yet tells us what happens when the direction of the question is
reversed.

Suppose the origin is known.

Suppose the trajectory proceeds forward without interruption.

Suppose almost nothing materially original survives at its endpoint.

What, then, has continued?

That is the problem of \emph{Macrolife}.

% ============================================================================
%  MOVEMENT II
%  CONTINUATION
% ============================================================================

\section{Continuation}
\label{mov:continuation}

\movementquestion{What may something become and still count as its continuation?}

\begin{orientation}
The problem of origin begins with something that already exists and asks how
it came to be there. The problem of continuation reverses the direction of
attention. We begin with something whose existence is not in doubt and ask
what may happen to it without bringing its history to an end.

Ordinary cases make this question seem easier than it is. A house remains the
same house after it is painted. A person remains the same person after a
haircut. A city survives the replacement of roads, buildings, governments,
and inhabitants. In such cases we rarely need to specify which changes are
permitted because the changes remain comfortably inside inherited conventions
of identity.

The difficulty appears when replacement ceases to be local. What if every
material component is eventually exchanged? What if the environment that
defined the original object is abandoned? What if continuation requires not
preserving a form but repeatedly changing it? What if the descendant survives
only because almost nothing about the original arrangement is held fixed?

George Zebrowski's \emph{Macrolife} pushes precisely in this direction. Its
large-scale technological imagination treats planetary existence not as the
permanent condition of civilization but as one historical phase among others.
The relevant question is therefore not whether an original state can be kept
unchanged. It is whether a sequence of transformations can remain a
continuation even when its endpoint would no longer be recognized as materially
identical with its beginning.
\end{orientation}

% -----------------------------------------------------------------------------
%  2.1  THE QUESTION POINTS FORWARD
% -----------------------------------------------------------------------------

\subsection{The Question Points Forward}
\label{sec:question-forward}

Movement~\ref{mov:origin} considered an inverse problem. Given a present state
\(P\), the Originist reasons over histories capable of terminating in that
state:

\[
  P
  \;\rightsquigarrow\;
  \Gamma^{-}(P).
\]

The causal history itself runs in the opposite direction. The present is an
effect from which possible antecedents are inferred.

\emph{Macrolife} changes the orientation. Its characteristic problem begins
with an existing civilization and asks what forms its continuation may take
as the conditions of survival change \cite{zebrowski1979}. The question is
prospective:

\[
  O
  \longrightarrow
  \{\text{possible descendants}\}.
\]

This is not simply Card's problem with the arrows reversed.

In the retrospective case, the historical transformations have already
occurred. Their effects constrain reconstruction. In the prospective case,
the relevant transformations have not all occurred, and the space of possible
continuations is correspondingly open.

The asymmetry can be expressed provisionally as

\[
  \text{past inference}
  :
  P\rightsquigarrow\Gamma^{-}(P),
\]

whereas

\[
  \text{future continuation}
  :
  O\rightsquigarrow\Gamma^{+}(O).
\]

Here \(\Gamma^{+}(O)\) denotes candidate trajectories beginning at \(O\).

The two spaces are related by history, but they are not interchangeable.
A realized present may constrain its possible pasts without uniquely
determining them, while a present origin may permit many futures of which
only one, if any, will actually occur.

Originism contracts an already realized history into a set of admissible
reconstructions.

Continuation opens an unrealized history into a set of possible descendants.

% -----------------------------------------------------------------------------
%  2.2  MACROLIFE AND MATERIAL INDEPENDENCE
% -----------------------------------------------------------------------------

\subsection{Macrolife and Material Independence}
\label{sec:macrolife-material}

Zebrowski's central speculative move is to treat large artificial habitats,
spacefaring communities, and eventually much larger forms of organized life
as continuations of terrestrial civilization rather than as merely detached
artifacts produced by it \cite{zebrowski1979}. The conceptual importance of
this move lies less in any particular engineering proposal than in the
relation it establishes between survival and transformation.

A planet-bound civilization can easily mistake persistence of environment for
persistence of civilization. The landscape, atmosphere, gravitational field,
ecology, inherited architecture, and biological substrate appear as stable
background conditions. They are so persistent relative to an individual
human life that they can be treated as though they were not themselves parts
of a historical arrangement.

Macrolife removes that convenience.

Once a civilization becomes capable of constructing and inhabiting its own
environments, a distinction appears between preserving a civilization and
preserving the conditions under which that civilization first developed.
These projects may coincide for a time. They need not coincide indefinitely.

Let \(C_0\) denote a civilization at some initial stage. After a sequence of
transformations,

\[
  C_0
  \xrightarrow{T_1}
  C_1
  \xrightarrow{T_2}
  \cdots
  \xrightarrow{T_n}
  C_n,
\]

it may be the case that

\[
  C_0\neq_{\mathrm{material}}C_n.
\]

Indeed, under sufficiently extensive replacement, there may be no particular
material component whose persistence is necessary for the entire sequence to
be regarded as continuous.

This gives the second major failure of a state-based criterion.

\begin{pathology}[Material discontinuity without historical discontinuity]
\label{path:material-discontinuity}
There exist plausible histories in which

\[
  x_0\neq_{\mathrm{material}}x_n
\]

while the sequence

\[
  x_0
  \longrightarrow
  x_1
  \longrightarrow
  \cdots
  \longrightarrow
  x_n
\]

is nevertheless treated as a continuous history of transformation.
Consequently, material identity is not a necessary condition for every
relevant notion of continuation.
\end{pathology}

The conclusion does not establish that material continuity is irrelevant.
It establishes only that material identity and historical continuation cannot
be silently identified.

% -----------------------------------------------------------------------------
%  2.3  THE SHIP THAT NEVER NEEDED TO BE A SHIP
% -----------------------------------------------------------------------------

\subsection{Replacement Without Resurrection}
\label{sec:replacement}

The familiar Ship of Theseus problem is often posed as a puzzle concerning the
identity of an object after gradual replacement. Its usefulness here is more
limited and more precise.

Suppose an object \(S_0\) undergoes component replacement over time:

\[
  S_0
  \longrightarrow
  S_1
  \longrightarrow
  S_2
  \longrightarrow
  \cdots
  \longrightarrow
  S_n.
\]

At each transition, only part of the object is replaced. No stage requires the
construction of a complete copy from an external description. There is a
continuous process connecting adjacent states.

Now compare this history with a different operation. Destroy \(S_0\), retain
a sufficiently complete description \(D(S_0)\), and later construct \(S^\ast\)
from that description:

\[
  S_0
  \longrightarrow
  D(S_0)
  \dashrightarrow
  S^\ast.
\]

It is possible that

\[
  S_n\eqbyte S^\ast
\]

for an appropriate digital analogue, or that

\[
  S_n\simI S^\ast
\]

under a chosen structural invariant.

Nevertheless, the histories differ.

The first is replacement within an ongoing trajectory.

The second contains reconstruction from a representation.

This difference is particularly important in discussions of copying,
uploading, restoration, and consciousness transfer. If endpoint similarity is
treated as sufficient, gradual expansion and destructive copying become
difficult to distinguish. Once trajectory is retained, they are different
operations before any metaphysical conclusion about personal identity is
drawn.

\begin{proposition}[Replacement--reconstruction distinction]
\label{prop:replacement-reconstruction}
Endpoint equivalence between a continuously transformed descendant and a
reconstructed copy does not imply equivalence of their derivational
trajectories.
\end{proposition}

\begin{proof}
Let

\[
  \gamma_{\mathrm{cont}}
  =
  \left(
  S_0\to S_1\to\cdots\to S_n
  \right)
\]

be a trajectory in which each state arises through local transformation of
the preceding state. Let

\[
  \gamma_{\mathrm{recon}}
  =
  \left(
  S_0\to D(S_0)\dashrightarrow S^\ast
  \right)
\]

contain a stage at which the later object is generated from a representation
rather than by local continuation of the material organization represented by
the preceding object.

Even if

\[
  S_n\sim S^\ast,
\]

the ordered transformations constituting the two trajectories differ.
Therefore

\[
  \gamma_{\mathrm{cont}}
  \neq
  \gamma_{\mathrm{recon}}.
\]

Endpoint equivalence alone cannot remove this distinction.
\end{proof}

Notice what the proposition does \emph{not} establish. It does not prove that
\(S_n\) is the original while \(S^\ast\) is not. Such a conclusion would
require a criterion of identity not yet supplied.

The proposition establishes something weaker and more useful: a theory that
collapses these histories merely because their endpoints agree has discarded
information before establishing that the information is irrelevant.

% -----------------------------------------------------------------------------
%  2.4  CAUSAL CONTINUITY
% -----------------------------------------------------------------------------

\subsection{The Temptation of Causal Continuity}
\label{sec:causal-continuity}

A natural response is to replace material identity with causal continuity.

Instead of requiring

\[
  x\eqmaterial y,
\]

we might require

\[
  x\leadcausal y.
\]

The proposal has considerable force. A continuously modified ship is causally
connected to its earlier stages. A civilization constructing new habitats is
causally connected to the civilization that designed them. A living organism
is causally continuous with earlier stages despite extensive material
turnover.

Causal continuity therefore captures something material identity misses.

It is still too weak.

Consider a file \(F\) corrupted by a destructive process:

\[
  F
  \leadcausal
  F_{\mathrm{corrupt}}.
\]

The corruption is causally downstream from the file. Yet this fact alone does
not establish that the corrupted object is an adequate continuation of the
file for every purpose.

Similarly, a city may causally produce its ruins:

\[
  C
  \leadcausal
  R.
\]

The ruins would not exist in their present configuration without the city.
Yet the claim that the city therefore continues to exist as the ruins is at
best context-dependent.

Biological examples make the problem sharper. Tissue can undergo a continuous
causal process producing pathological growth. The resulting structure may be
causally continuous with the tissue from which it arose while violating the
organizational conditions under which we would call it continuation of the
same functioning tissue.

Thus

\[
  x\leadcausal y
\]

does not by itself tell us which properties, relations, or constraints must
survive for \(y\) to count as an acceptable continuation of \(x\).

\begin{pathology}[Causal succession without admissible continuation]
\label{path:causal-insufficiency}
There exist causally continuous transformations

\[
  x\leadcausal y
\]

for which \(y\) fails a criterion required for the relevant notion of
continuation. Therefore causal continuity is not sufficient, by itself, to
define continuation.
\end{pathology}

This pathology forces the next enlargement of the formalism.

% -----------------------------------------------------------------------------
%  2.5  ADMISSIBILITY
% -----------------------------------------------------------------------------

\subsection{Admissible Transformation}
\label{sec:admissibility}

The question is no longer merely whether a transformation occurred.

The question is whether it belongs to the class of transformations under
which the relevant continuation relation is preserved.

Let \(X\) again denote a state space. We now associate with the problem a
collection

\[
  \Adm
\]

of admissible transformations or admissible trajectories.

The term \enquote{admissible} is intentionally relational. A transformation
is not admissible in the abstract. It is admissible relative to a criterion,
task, invariant, or world-constituting structure.

\begin{definition}[Admissible continuation]
\label{def:admissible-continuation}
Let \(\Adm\) be a specified family of transformations on a state space \(X\).
A trajectory

\[
  \gamma:
  x_0
  \xrightarrow{T_1}
  x_1
  \xrightarrow{T_2}
  \cdots
  \xrightarrow{T_n}
  x_n
\]

is an \emph{\(\Adm\)-admissible continuation} when each transformation, or the
trajectory considered as a whole when admissibility is nonlocal, satisfies
the constraints defining \(\Adm\).
\end{definition}

This definition deliberately permits admissibility to be nonlocal.

It may be tempting to write

\[
  T_i\in\Adm
  \qquad
  \text{for every }i
\]

and conclude that the entire trajectory is admissible. That condition is
useful when admissibility composes locally. It need not always do so.

A sequence of individually harmless transformations may accumulate into a
change that destroys an invariant. Conversely, a locally disruptive
transformation may belong to a globally admissible repair process.

The distinction will later connect directly to TARTAN and to the author's
constraint-based treatment of realization. For now, the important result is
simpler.

Continuation requires more than causal ancestry.

It requires a specification of what kinds of transformation the relevant
continuation relation permits.

% -----------------------------------------------------------------------------
%  2.6  CONTINUATION IS TYPED
% -----------------------------------------------------------------------------

\subsection{There Is No Unqualified Continuation Relation}
\label{sec:typed-continuation}

Once admissibility is explicit, the phrase \enquote{the same thing continued}
becomes suspiciously incomplete.

Continued as what?

A manuscript may continue as a semantic work while ceasing to continue as a
physical artifact. A biological lineage may continue while every organism
that once instantiated it dies. An institution may continue legally while its
practices change radically. A civilization may continue genealogically while
ceasing to preserve a particular political order.

Continuation is therefore typed by the structure whose persistence is being
asserted.

Let \(\Inv\) denote a selected family of invariants. We may write

\[
  x\simI y
\]

when \(x\) and \(y\) agree with respect to those invariants.

But even this does not settle the historical question.

Two independently constructed objects may satisfy

\[
  x\simI y
\]

without either descending from the other.

Thus invariant preservation and genealogy remain distinct.

\begin{proposition}[Invariant agreement does not establish descent]
\label{prop:invariant-not-descent}
For a nontrivial invariant-equivalence relation \(\simI\),

\[
  x\simI y
\]

does not in general imply the existence of a historical trajectory from
\(x\) to \(y\).
\end{proposition}

\begin{proof}
Choose two states independently constructed so as to instantiate the same
selected invariants. By construction,

\[
  x\simI y.
\]

Their independent construction entails that neither need belong to the
derivational ancestry of the other. Hence invariant agreement alone does not
establish descent.
\end{proof}

We have therefore accumulated several relations that ordinary identity
language tends to compress:

\[
  \text{material identity},
  \qquad
  \text{causal continuity},
  \qquad
  \text{invariant preservation},
  \qquad
  \text{admissible continuation}.
\]

They are not interchangeable.

The notation will remain deliberately incomplete until concrete pathologies
force further distinctions.

% -----------------------------------------------------------------------------
%  2.7  COPY AND EXPANSION
% -----------------------------------------------------------------------------

\subsection{Copying and Continuous Expansion}
\label{sec:copy-expansion}

The distinction between copying and continuous expansion provides a useful
stress test.

Suppose a system \(S\) is represented with sufficient accuracy to construct a
second system \(S'\). The copying process has the schematic form

\[
  S
  \longrightarrow
  D(S)
  \longrightarrow
  S'.
\]

Now suppose instead that \(S\) expands by incorporating new components while
maintaining active causal and organizational coupling between old and new
regions:

\[
  S_0
  \hookrightarrow
  S_1
  \hookrightarrow
  S_2
  \hookrightarrow
  \cdots.
\]

At some later stage, none of the material initially present in \(S_0\) need
remain.

If endpoint structure alone is considered, the two processes may eventually
be made arbitrarily similar. If trajectories are retained, however, they
remain distinguishable.

The first process establishes a new realization from a description or
transfer channel.

The second enlarges an already operating realization.

This distinction has consequences for the author's earlier argument concerning
consciousness transfer, where destructive copying and continuous expansion
were treated as importantly different candidate processes of persistence.
The present framework supplies a more general reason for that difference:
they occupy different regions of trajectory space.

No claim about consciousness itself is required at this stage.

The relevant principle is simply

\[
  \text{endpoint resemblance}
  \neq
  \text{trajectory equivalence}.
\]

Macrolife generalizes the same issue from an individual system to a
civilization. A civilization that expands into constructed environments does
not need to preserve its initial matter. But neither does the mere existence
of a later structure resembling the civilization establish that the
civilization continued into it.

The missing variable is the admissible path.

% -----------------------------------------------------------------------------
%  2.8  IRREVERSIBILITY
% -----------------------------------------------------------------------------

\subsection{Why the Two Directions Do Not Cancel}
\label{sec:irreversibility}

It would now be easy to draw a pleasing symmetry.

Card gives

\[
  ?\longrightarrow P,
\]

while Zebrowski gives

\[
  O\longrightarrow ?.
\]

One might therefore imagine Originism and Macrolife as inverse views of a
single reversible relation.

That would be a mistake.

Suppose a transformation

\[
  T:x\longrightarrow y
\]

discards information. Then the existence of \(T\) does not guarantee an
inverse transformation

\[
  T^{-1}:y\longrightarrow x.
\]

Indeed, if distinct states satisfy

\[
  T(x_1)=T(x_2)=y,
\]

then \(y\) alone cannot determine which antecedent was realized.

This is the same non-injectivity encountered in
\cref{prop:endpoint-insufficiency}, now viewed as an irreversible operation.

\begin{definition}[Origin-relative irreversibility]
\label{def:irreversibility}
A transformation \(T:X\rightarrow Y\) is
\emph{origin-relative irreversible} with respect to a distinction
\(\delta\) when states distinguished by \(\delta\) can be mapped by \(T\)
to outputs from which that distinction cannot be recovered without
additional information.
\end{definition}

This definition does not require thermodynamic irreversibility. It describes
irreversibility relative to a provenance distinction.

Compression is the obvious example. If two distinguishable inputs are mapped
to one representation, the representation may preserve everything required
for one task while destroying what is required for another.

The consequence for the temporal symmetry is immediate.

Originist inference begins after transformations have occurred and asks which
pasts remain compatible with their surviving traces.

Macrolife continuation begins before all transformations have occurred and
asks which futures remain admissible.

The first is constrained by realized losses.

The second is constrained by possible transformations.

They face each other across the present, but they do not undo one another.

% -----------------------------------------------------------------------------
%  2.9  THE TWO CONES
% -----------------------------------------------------------------------------

\subsection{Past Contraction and Future Expansion}
\label{sec:two-cones}

Let \(P\) denote a present state.

Define

\[
  \Gamma^{-}_{\Adm}(P)
\]

as the family of admissible histories terminating in \(P\), and

\[
  \Gamma^{+}_{\Adm}(P)
\]

as the family of admissible continuations beginning at \(P\).

The present then occupies an asymmetric position:

\[
  \Gamma^{-}_{\Adm}(P)
  \longrightarrow
  P
  \longrightarrow
  \Gamma^{+}_{\Adm}(P).
\]

The notation should not be mistaken for a claim that every member of the left
family actually occurred or that every member of the right family will occur.
The left side contains candidate histories consistent with surviving
constraints. The right side contains candidate continuations permitted by
current constraints.

Evidence tends to reduce the left family.

Possibility tends to enlarge the right family.

Additional historical evidence may eliminate candidate pasts:

\[
  \Gamma^{-}_{\Adm}(P)
  \supseteq
  \Gamma^{-}_{\Adm}(P\mid E_1)
  \supseteq
  \Gamma^{-}_{\Adm}(P\mid E_1,E_2).
\]

By contrast, newly available technologies or organizational transformations
may enlarge the space of possible futures:

\[
  \Gamma^{+}_{\Adm_1}(P)
  \subseteq
  \Gamma^{+}_{\Adm_2}(P),
\]

when \(\Adm_2\) permits transformations unavailable under \(\Adm_1\).

This is the precise sense in which the two literary probes face opposite
temporal directions without becoming inverses.

% -----------------------------------------------------------------------------
%  2.10  THE MACROLIFE LIMIT
% -----------------------------------------------------------------------------

\subsection{When Almost Nothing Original Remains}
\label{sec:macrolife-limit}

The strongest version of the Macrolife problem occurs when material
replacement approaches completeness.

Suppose

\[
  C_0
  \longrightarrow
  C_1
  \longrightarrow
  \cdots
  \longrightarrow
  C_n
\]

is treated as a continuous civilizational history. Let \(M(C_i)\) denote the
material constituents of \(C_i\).

It may eventually be the case that

\[
  M(C_0)\cap M(C_n)=\varnothing.
\]

Nothing in the initial material inventory survives.

Yet the conclusion

\[
  C_0\not\rightsquigarrow C_n
\]

does not follow merely from that fact.

If continuation is carried by organization, causal dependence, transmitted
constraints, institutional inheritance, memory, reproduction, or some
combination of these, then disappearance of the original matter may be
compatible with historical continuation.

This produces the question from which the title of the essay is drawn.

\begin{held}[Nothing Original Need Survive]
\centering
If continuation is carried by an admissible history rather than by
persistence of an initial material state, then complete replacement of the
original material is compatible, in principle, with continuation.
\end{held}

The proposition is deliberately conditional.

The title does not claim that preservation never matters. Nor does it claim
that every causal descendant counts as continuation.

It claims that original material need not be the invariant that does the
work.

What does the work instead remains to be determined.

% -----------------------------------------------------------------------------
%  2.11  INSTRUCTIONS AS HISTORICAL OBJECTS
% -----------------------------------------------------------------------------

\subsection{Instructions as Historical Objects: Prompt Injection and the
Provenance of Constraints}
\label{sec:prompt-injection-authority}

The admissibility apparatus developed in this movement has a consequence
worth making explicit before the movement closes: a representation's
provenance can determine not only what it means but what it is entitled to
do.

Suppose a system is operating under an admissibility structure \(\mathcal A\)
determined by a constraint family \(\mathcal C_0\) (\cref{sec:admissibility}),
and receives some representation \(p\) through an ordinary data channel. The
transformation of interest is not an ordinary state change

\[
  x\longrightarrow x',
\]

but

\[
  (\mathcal C_0,x)
  \xrightarrow{\;p\;}
  (\mathcal C_p,x'),
\]

in which \(p\) is interpreted as licensed to modify the constraint family
itself, rather than merely to serve as data evaluated under it.

\begin{definition}[Prompt injection]
\label{def:prompt-injection}
Prompt injection occurs when a representation whose provenance licenses it
only as data under \(\mathcal C_0\) is interpreted as authority over the
constraint family governing subsequent transformations, producing
\(\mathcal C_p\neq\mathcal C_0\) without any change in the admissibility
structure having been independently warranted.
\end{definition}

This is deliberately not a definition in terms of malicious content. Malice
is incidental to the structural failure. Consider a string \(p_A\) supplied
by an authorized controller, and a byte-identical string \(p_B\) occurring
inside a document retrieved for an unrelated purpose, such as summarization.
At the state level,

\[
  p_A\eqbyte p_B.
\]

But their provenance differs:

\[
  \Prov(p_A)\neq\Prov(p_B).
\]

\begin{pathology}[Provenance-insensitive authority attribution]
\label{path:prompt-injection}
A system exhibits provenance-insensitive authority attribution when it
licenses

\[
  p_A\eqbyte p_B
  \;\Longrightarrow\;
  p_A=_{\mathrm{authority}}p_B,
\]

that is, when byte-level equivalence is treated as sufficient for equal
instructional authority, regardless of \(\Prov(p_A)\) and \(\Prov(p_B)\).
\end{pathology}

\begin{proof}
This is the strong Originist objection of
\cref{sec:first-originist-objection} instantiated at \(\delta=\Prov\). The
system quotients away exactly the distinction — provenance — that a later
operation (interpreting \(p\) as an instruction) requires. By
\cref{sec:first-originist-objection}, this is precisely the condition under
which a quotient becomes a wound rather than a harmless abstraction.
\end{proof}

\begin{held}[What prompt injection supplies]
\centering
Semantic content does not determine instructional authority.
\end{held}

The computational analogue of the manuscript's central distinction is now
available in operational form:

\[
  \boxed{\State\neq\History}
  \qquad\text{becomes}\qquad
  \boxed{\text{content}\neq\text{authority.}}
\]

Chain of Memory, introduced later, supplies the constructive half of this
pathology: a system does not need to retain everything about an
instruction's history in order to avoid this failure. It needs to preserve
exactly those provenance distinctions on which future authority and
admissibility judgments depend — no more, and no less.
\Cref{sec:shannon-markov} established that content does not reveal its
provenance; this section has shown why failing to track that fact can
change what content is permitted to do.

% -----------------------------------------------------------------------------
%  2.12  FROM CONTINUATION TO WORLD
% -----------------------------------------------------------------------------

\subsection{The First Appearance of the World Problem}
\label{sec:first-world-problem}

At civilizational scale, a further difficulty begins to appear.

Suppose \(C_t\) and \(C_{t+n}\) belong to an admissible civilizational
trajectory:

\[
  C_t
  \leadcausal
  C_{t+n}.
\]

Suppose further that enough organizational structure survives for us to call
the latter a descendant of the former.

A stronger question remains unanswered.

Did the same \emph{world} continue?

A civilization can alter its environment, institutions, bodies, memory
systems, perceptual interfaces, spatial organization, and relations among
agents. At some point it becomes possible that the genealogy of the
civilization remains intelligible while the relational field within which its
earlier forms made sense has changed profoundly.

We must therefore distinguish, at least provisionally,

\[
  \text{persistence of an object},
\]

\[
  \text{persistence of a civilization},
\]

and

\[
  \text{persistence of a world}.
\]

No formal definition of world-continuation will be offered here. Doing so
would outrun the examples.

The distinction has appeared, but it has not yet been earned mathematically.

% -----------------------------------------------------------------------------
%  2.13  WHAT MOVEMENT II HAS FORCED
% -----------------------------------------------------------------------------

\subsection{What Continuation Has Failed to Preserve}
\label{sec:continuation-forced}

Movement~\ref{mov:origin} established that a state does not determine its
history.

The present movement has established a complementary failure: difference of
state does not determine discontinuity of history.

The two failures can be placed side by side:

\[
  x_A=x_B
  \quad\text{while}\quad
  \gamma_A\neq\gamma_B,
\]

and

\[
  x_A\neq_{\mathrm{material}}x_B
  \quad\text{while}\quad
  x_A\leadcausal x_B.
\]

State equality can conceal historical difference.

State difference can coexist with historical continuation.

Neither direction permits identity of states to substitute for analysis of
trajectories.

But causal trajectory alone has also failed. Corruption, destruction, and
pathological transformation can all be causally continuous. We were therefore
forced to introduce admissibility:

\[
  \gamma
  \quad\longrightarrow\quad
  (\gamma,\Adm).
\]

This is the second enlargement of the representational object.

It should again be read minimally. We have not shown that
\((\gamma,\Adm)\) is fundamental. We have shown that an untyped causal
trajectory is insufficient to discriminate the cases already encountered.

\begin{held}[Continuation Under Constraint]
\centering
The question is not whether a later state was caused by an earlier one.\\[3pt]
The question is which transformations preserve the kind of continuation
being claimed.
\end{held}

Card left us with an inaccessible beginning.

Zebrowski has now given us an indefinitely transformable future.

There remains another possibility.

A trajectory can leave its origin and come back.

If it returns, what exactly has returned to what?

The answer cannot simply be

\[
  O\longrightarrow X\longrightarrow O,
\]

because both the traveler and the origin may have histories of their own.

The line is about to become a loop.

And the loop, in turn, will fail to remain a line.

% ============================================================================
%  MOVEMENT III
%  RECURRENCE
% ============================================================================

\section{Recurrence}
\label{mov:recurrence}

\movementquestion{What does it mean for what departed to return?}

\begin{orientation}
Returning somewhere is one of the most ordinary forms of continuity. A person
leaves a house in the morning and returns in the evening. A traveler leaves a
country and comes home years later. A spacecraft departs from a planet,
completes a journey, and returns to its point of departure. The language of
return suggests a loop: something leaves, undergoes a history elsewhere, and
comes back to where it began.

The apparent simplicity depends upon ignoring two changes.

The traveler who returns has acquired a history. The place to which the
traveler returns has acquired one as well.

A childhood home revisited after thirty years is not the state that was left
thirty years earlier. Even if the building remains materially similar, its
inhabitants, surroundings, uses, and relation to the traveler may have
changed. Nor is the returning traveler identical in state to the person who
departed. Return therefore cannot ordinarily mean restoration of the initial
configuration.

Ursula K.\ Le Guin's \emph{The Dispossessed} makes this problem unusually
visible. Shevek's journey from Anarres to Urras and eventual return to Anarres
is geographical, political, intellectual, and relational at once. He crosses
a boundary between worlds that are themselves historically connected. He
carries ideas across that boundary. Those ideas acquire new possible
recipients. When he returns, the apparent loop does not close by restoring its
starting point.

The problem of recurrence is therefore neither Card's problem of reconstructing
an absent beginning nor Zebrowski's problem of continuing indefinitely into an
open future. It asks what sort of relation can connect a later state to an
earlier configuration without pretending that history between them did not
occur.
\end{orientation}

% -----------------------------------------------------------------------------
%  3.1  THE WALL
% -----------------------------------------------------------------------------

\subsection{The Wall}
\label{sec:wall}

\emph{The Dispossessed} begins with a wall \cite{leguin1974}. The object is
physically unimpressive. Its conceptual function is not.

The wall separates the spaceport on Anarres from the surrounding society. Yet
the meaning of the enclosure depends upon which side is taken as inside. From
one perspective the wall encloses the landing field. From another, it encloses
the rest of Anarres by marking the boundary across which contact with Urras
and other worlds is controlled.

A boundary therefore need not merely divide two already constituted regions.
It can participate in determining which region counts as interior, which as
exterior, and which transformations are permitted between them.

This makes the wall more than a metaphor for political separation. It is an
instance of a general computational fact about boundaries: a boundary
classifies crossings.

Let \(A\) denote a region identified with Anarres and \(U\) a region identified
with Urras. A naive representation would write

\[
  A\mid U.
\]

But the bar alone says nothing about what may cross it.

A more useful representation associates the boundary \(B\) with a set of
admissible crossing operations

\[
  \Adm_B.
\]

The boundary is then not merely a location. It participates in determining
which trajectories are realizable.

\begin{definition}[Boundary-relative admissibility]
\label{def:boundary-admissibility}
Let \(B\) separate or relate regions \(X\) and \(Y\). The
\emph{boundary-relative admissibility structure} \(\Adm_B\) specifies the
class of transformations by which states, agents, information, or material
may cross between \(X\) and \(Y\) while remaining recognizable under the
relevant system of constraints.
\end{definition}

This definition is intentionally broad. A wall may physically prevent bodies
while allowing radio signals. A legal boundary may permit persons while
restricting property. An institutional boundary may admit information only
after translation into an accepted vocabulary. A computational interface may
permit data while refusing operations.

The point is not that all boundaries are the same.

It is that a boundary can be characterized partly by what it permits to pass.

\begin{crossing}[Shevek's departure]
\label{cross:shevek-departure}
Shevek's movement from Anarres to Urras is not merely a change of spatial
coordinate. It is a transition across a boundary whose political,
institutional, economic, and epistemic constraints differ on its two sides.
The trajectory therefore changes not only Shevek's location but the set of
transformations available to him.
\end{crossing}

The wall introduces admissibility into space.

Return will introduce it into recurrence.

% -----------------------------------------------------------------------------
%  3.2  RETURN IS NOT RESTORATION
% -----------------------------------------------------------------------------

\subsection{Return Is Not Restoration}
\label{sec:return-not-restoration}

Let \(O\) denote a state at departure.

A simple model of return might be written

\[
  O
  \rightsquigarrow
  X
  \rightsquigarrow
  O.
\]

The notation is misleading because it uses the same symbol for two states
separated by a nontrivial history.

Let the state of origin at departure be \(O_t\), and let the corresponding
state encountered upon return be \(O_{t+n}\). Let the traveler at departure be
\(S_t\), and the traveler upon return \(S_{t+n}\).

Then the relevant configuration is closer to

\[
  (S_t,O_t)
  \longrightarrow
  X
  \longrightarrow
  (S_{t+n},O_{t+n}).
\]

In general,

\[
  S_t\neq S_{t+n}
\]

as states, and

\[
  O_t\neq O_{t+n}.
\]

Yet we may still have reason to describe the event as a return.

This motivates a recurrence relation weaker than state identity.

Write

\[
  O_{t+n}\simconstraint O_t
\]

when the later state stands in the relevant recurrence relation to the earlier
state under a selected family of constraints \(\mathcal C\).

The important fact is

\[
  O_{t+n}\neq O_t
  \qquad\text{while}\qquad
  O_{t+n}\simconstraint O_t.
\]

Return therefore resembles continuation more than restoration.

\begin{definition}[Constraint-relative recurrence]
\label{def:constraint-recurrence}
Given states \(x_t\) and \(x_{t+n}\) belonging to a common historical
structure, \(x_{t+n}\) is a \emph{constraint-relative recurrence} of \(x_t\)
when

\[
  x_{t+n}\sim_{\mathcal C}x_t
\]

for a specified family of recurrence constraints \(\mathcal C\), without
requiring

\[
  x_{t+n}=x_t.
\]
\end{definition}

The definition deliberately refuses to identify recurrence with periodicity.
A periodic system returns to an equivalent dynamical state according to a
specified evolution. Historical return is usually messier. The relevant
equivalence may be institutional, geographical, genealogical, semantic, or
world-relative.

Thus a person can return home even though neither person nor home has remained
unchanged.

\begin{proposition}[Return does not imply restoration]
\label{prop:return-not-restoration}
Constraint-relative recurrence does not in general imply state identity.
That is,

\[
  x_{t+n}\sim_{\mathcal C}x_t
  \nimplies
  x_{t+n}=x_t.
\]
\end{proposition}

\begin{proof}
Choose \(\mathcal C\) to preserve the relations relevant to return while
permitting transformations not relevant to those relations. For example,
\(\mathcal C\) may preserve geographical and institutional continuity while
allowing changes in material composition and internal state. A later state
may then satisfy \(\mathcal C\) without being identical to the earlier state.
Therefore recurrence under \(\mathcal C\) does not imply restoration of the
earlier state.
\end{proof}

The proposition is formally modest. Its importance lies in preventing a loop
from being mistaken for an erasure of history.

% -----------------------------------------------------------------------------
%  3.3  THE RETURNING ORIGIN HAS MOVED
% -----------------------------------------------------------------------------

\subsection{The Returning Origin Has Moved}
\label{sec:returning-origin}

The word \enquote{origin} now becomes unstable.

Anarres is Shevek's origin in the immediate geographical and social sense. He
departs from it and later returns to it. But Anarres is not an origin without
a history. Its society emerged from an earlier revolutionary separation
connected to Urras and to the thought of Odo \cite{leguin1974}.

The origin of one trajectory is therefore itself an endpoint within another.

Instead of

\[
  O
  \longrightarrow
  X
  \longrightarrow
  O',
\]

we obtain something closer to

\[
  O_0
  \longrightarrow
  O_1
  \longrightarrow
  X
  \longrightarrow
  O_1'.
\]

Here \(O_1\) may function as an origin relative to Shevek's journey while
remaining a descendant relative to the longer political history.

There is no contradiction.

Origin is indexed by the trajectory under consideration.

\begin{definition}[Relative origin]
\label{def:relative-origin}
Given a trajectory or provenance structure \(\gamma\), a state \(O\) is an
\emph{origin relative to \(\gamma\)} when it is selected as the relevant
initial state for the derivational question posed over \(\gamma\). This does
not imply that \(O\) has no antecedent history outside \(\gamma\).
\end{definition}

This qualification alters the Originist problem.

Movement~\ref{mov:origin} provisionally treated the origin as the inaccessible
left boundary of a trajectory:

\[
  O
  \longrightarrow
  \cdots
  \longrightarrow
  P.
\]

We can now ask:

\[
  \text{Which }O?
\]

A manuscript may have an origin as a physical artifact, another as a textual
tradition, another as a semantic composition, and another as the endpoint of
earlier oral or documentary transmission.

Likewise, Anarres can be an origin for Shevek while remaining a historical
descendant of Urrasti political conflict.

Origins have provenance.

% -----------------------------------------------------------------------------
%  3.4  THE ORIGINIST TURNS AROUND
% -----------------------------------------------------------------------------

\subsection{The Originist Turns Around}
\label{sec:originist-turns}

The discovery that origins themselves possess histories produces an apparent
regress.

If

\[
  O_1
\]

has provenance, then perhaps we must reconstruct

\[
  O_0\longrightarrow O_1.
\]

But if \(O_0\) also has provenance, we may seek an earlier state \(O_{-1}\),
and so on.

Nothing guarantees a uniquely privileged absolute origin.

This does not make origin-relative reasoning incoherent. It changes what an
origin claim means.

To say that \(O\) is an origin is not necessarily to say

\[
  \nexists x\;(x\longrightarrow O).
\]

It may instead mean that \(O\) is the boundary at which the current
provenance question begins.

The distinction is familiar outside speculative fiction. The origin of a
printed edition need not be the origin of the work. The origin of a species
under one taxonomic criterion need not be a single organism. The origin of an
institution may refer to its legal incorporation even when the practices that
constitute it developed earlier.

Originism therefore requires declared scope.

Without such scope, the command to \enquote{return to the origin} can become
an instruction to pursue an indefinitely receding boundary.

\begin{remark}
The regress is not necessarily vicious. A provenance inquiry may terminate
because an earlier distinction is irrelevant to the task, because evidence
has become unavailable, or because the selected origin is explicitly
relative. What Originism resists is not termination but unacknowledged
termination.
\end{remark}

The distinction is another instance of a principle already encountered.

Forgetting is sometimes necessary.

Forgetting that one has forgotten is more dangerous.

% -----------------------------------------------------------------------------
%  3.5  SHEVEK'S PHYSICS AND THE PROBLEM OF TRANSMISSION
% -----------------------------------------------------------------------------

\subsection{Transmission Without a Single Custodian}
\label{sec:shevek-transmission}

Shevek's scientific work introduces a different complication.

Knowledge can cross a boundary without remaining under the control of the
institution through which it passed. The political significance of Shevek's
physics is tied not merely to what the theory says but to who may possess,
transmit, restrict, or monopolize it. His response to attempted appropriation
is dissemination: the work is communicated beyond the immediate Urrasti
institutions and made available through a wider network \cite{leguin1974}.

For present purposes, the important operation is branching.

Suppose a source artifact or claim \(K\) is transmitted independently to
several recipients:

\[
  K
  \longrightarrow
  K_1,
  \qquad
  K
  \longrightarrow
  K_2,
  \qquad
  K
  \longrightarrow
  K_3.
\]

The provenance structure is no longer adequately represented by a single
chain.

We may draw

\[
\begin{tikzcd}[column sep=3.2em,row sep=1.5em]
& K_1 \\
K \arrow[ur] \arrow[r] \arrow[dr] & K_2 \\
& K_3
\end{tikzcd}
\]

rather than

\[
  K\longrightarrow K_1\longrightarrow K_2\longrightarrow K_3.
\]

The distinction is obvious once drawn, but it changes the formal object
substantially.

A trajectory has become a provenance graph.

% -----------------------------------------------------------------------------
%  3.6  FROM TRAJECTORY TO PROVENANCE GRAPH
% -----------------------------------------------------------------------------

\subsection{The Failure of Linear History}
\label{sec:failure-linear-history}

Let

\[
  \graphG=(V,E)
\]

be a directed graph whose vertices represent relevant states and whose edges
represent recorded or hypothesized derivational transformations.

\begin{definition}[Provenance graph]
\label{def:provenance-graph}
A \emph{provenance graph} is a directed graph

\[
  \graphG=(V,E,\tau),
\]

where \(V\) is a collection of states, \(E\) is a collection of directed
derivational relations among those states, and

\[
  \tau:E\longrightarrow\mathcal T
\]

assigns, when known, a transformation type to each edge.
\end{definition}

A linear trajectory is now a special case: a directed path through
\(\graphG\).

This enlargement permits branching:

\[
  x\longrightarrow y_1,
  \qquad
  x\longrightarrow y_2,
\]

but also convergence:

\[
  x_1\longrightarrow y,
  \qquad
  x_2\longrightarrow y.
\]

It permits recurrence when later vertices stand in a declared recurrence
relation to earlier ones. It permits missing edges, uncertain transformations,
and competing derivational hypotheses.

Most importantly, it prevents the history of an artifact from being confused
with a single privileged chain when the evidence supports a more complicated
structure.

\begin{pathology}[Linear-history insufficiency]
\label{path:linear-history}
If a provenance structure contains branching or convergence relevant to the
question under study, then no single trajectory can represent the complete
structure without discarding derivational alternatives or descendants.
\end{pathology}

\begin{proof}
Suppose \(x\in V\) has distinct outgoing derivational edges

\[
  x\to y_1
  \qquad\text{and}\qquad
  x\to y_2,
\]

where both descendants are relevant. Any single directed path through \(x\)
can follow at most one of these edges immediately after \(x\). Therefore the
path omits at least one relevant branch. The same argument applies in reverse
to relevant convergence. Hence a single trajectory is insufficient to encode
the complete provenance structure.
\end{proof}

Movement~\ref{mov:origin} forced

\[
  x
  \rightsquigarrow
  \gamma.
\]

Le Guin now forces

\[
  \gamma
  \rightsquigarrow
  \graphG.
\]

The trajectory was not wrong.

It was a special case mistaken, temporarily, for a sufficient general form.

% -----------------------------------------------------------------------------
%  3.7  GENERATION BECOMES PATH-RELATIVE
% -----------------------------------------------------------------------------

\subsection{Generation After Branching}
\label{sec:generation-branching}

The move from trajectories to provenance graphs immediately complicates the
generation-depth definition from \cref{def:generation-depth}.

In a single trajectory,

\[
  O
  \to
  x_1
  \to
  \cdots
  \to
  x_n,
\]

generation depth can be counted along the unique represented path.

In a graph, there may be several paths from \(O\) to the same state \(x\).

Let

\[
  \Pi_{\graphG}(O,x)
\]

denote the set of directed paths from \(O\) to \(x\) in \(\graphG\).

Then the expression

\[
  \Depth_O(x)
\]

may be ambiguous whenever there exist

\[
  \pi_1,\pi_2\in\Pi_{\graphG}(O,x)
\]

with

\[
  |\pi_1|\neq|\pi_2|.
\]

Generation depth must therefore be indexed by a path unless an aggregation
rule has been declared:

\[
  \Depth_O(x;\pi).
\]

Possible summaries include

\[
  \Depth_O^{\min}(x)
  =
  \min_{\pi\in\Pi_{\graphG}(O,x)}
  |\pi|,
\]

or

\[
  \Depth_O^{\max}(x)
  =
  \max_{\pi\in\Pi_{\graphG}(O,x)}
  |\pi|.
\]

Neither is intrinsically the correct measure.

The shortest path answers how few recognized transformations are sufficient
to connect the origin to the state.

The longest path answers a different question.

A historically realized path answers another.

A path of highest evidential support answers yet another.

The apparently simple superscript

\[
  \varphi^{[3]}
\]

has therefore acquired an unstated parameter.

We will not repair the notation yet.

The pathology that makes the repair unavoidable belongs to
Movement~\ref{mov:pathologies}.

% -----------------------------------------------------------------------------
%  3.8  CONVERGENCE AFTER BRANCHING
% -----------------------------------------------------------------------------

\subsection{When Descendants Meet Again}
\label{sec:descendants-meet}

Branching is only half of the difficulty.

Suppose a source \(O\) produces two independently transformed descendants:

\[
  O\longrightarrow A
\]

and

\[
  O\longrightarrow B.
\]

Later, material from \(A\) and \(B\) is combined into \(C\):

\[
\begin{tikzcd}[column sep=3em]
O \arrow[r] \arrow[d] & A \arrow[d] \\
B \arrow[r] & C
\end{tikzcd}
\]

What is the generation of \(C\)?

More importantly, what is its origin?

The singular form has become misleading.

\(C\) may possess multiple relevant ancestors. A claim in \(C\) may derive
from \(A\), another from \(B\), and another from a reconstruction produced
only after their convergence.

Thus provenance may need to be tracked below the level of the entire artifact.

For a document \(D\) containing claims

\[
  \varphi_1,\varphi_2,\ldots,\varphi_n,
\]

there need not exist a single provenance value \(p(D)\) adequate for every
claim.

Instead one may require

\[
  \Prov(\varphi_i)
\]

for individual claims or spans.

This observation justifies the otherwise eccentric typographic apparatus of
the present essay. A citation attached to an entire document cannot always
represent the derivational history of each statement within it.

Provenance is potentially local.

% -----------------------------------------------------------------------------
%  3.9  THE RETURN LOOP AS A NON-INVERSE
% -----------------------------------------------------------------------------

\subsection{A Loop That Does Not Undo Itself}
\label{sec:noninverse-loop}

We can now return to Shevek's journey.

Let \(A_t\) denote Anarres at departure, \(S_t\) Shevek at departure, and
\(U_{t+k}\) the relevant Urrasti state encountered during the journey. The
return has the schematic form

\[
  (S_t,A_t)
  \longrightarrow
  (S_{t+k},U_{t+k})
  \longrightarrow
  (S_{t+n},A_{t+n}).
\]

The endpoint is related to the beginning, but it does not invert the path.

There is no operation

\[
  T^{-1}
\]

that removes what Shevek learned, restores Anarres to its earlier state,
recalls every transmission of his work, and returns the entire configuration
to \(t\).

The loop closes relationally while remaining historically open.

This distinction is central to the author's Constraint--Recurrence Principle.
Recurrence need not mean restoration of a prior state. It may instead mean
return to a constraint-compatible region after a nontrivial trajectory.

In the present notation,

\[
  A_{t+n}\simconstraint A_t
\]

can coexist with

\[
  \Hist(A_{t+n})\neq\Hist(A_t).
\]

Indeed, the latter difference is unavoidable if anything happened between the
two times.

\begin{proposition}[Historical non-idempotence of return]
\label{prop:historical-nonidempotence}
A return operation that restores a selected state-level relation need not
restore the historical state of the system.
\end{proposition}

\begin{proof}
Let a system leave \(x_t\), undergo at least one historically recorded
transformation, and later reach \(x_{t+n}\) such that

\[
  x_{t+n}\sim_{\mathcal C}x_t.
\]

The trajectory to \(x_{t+n}\) contains the intervening transformations,
whereas the trajectory terminating at \(x_t\) does not. Hence their histories
differ even when the selected recurrence relation is satisfied.
\end{proof}

A return can close a spatial or structural loop without closing history.

% -----------------------------------------------------------------------------
%  3.10  THREE TEMPORAL PROBES
% -----------------------------------------------------------------------------

\subsection{Three Temporal Probes}
\label{sec:three-probes}

The three literary works can now be placed together without reducing them to
instances of a generic theme of identity.

Card asks a retrospective question:

\[
  \text{Where did this come from?}
\]

Zebrowski asks a prospective question:

\[
  \text{What may this become and still continue?}
\]

Le Guin asks a recurrent question:

\[
  \text{What does it mean for what departed to return?}
\]

These are different operations on historical structure.

Originism begins with traces and constrains candidate ancestry.

Macrolife begins with an extant organization and opens possible continuation.

Recurrence begins with a relation between departure and return and asks what
must remain invariant for the return to be meaningful despite intervening
change.

The three forms can be written schematically as

\[
  \boxed{
  \begin{aligned}
    \textsc{Origin}
      &:\quad
      ?\leadsto P,
      \\[4pt]
    \textsc{Continuation}
      &:\quad
      O\leadsto ?,
      \\[4pt]
    \textsc{Recurrence}
      &:\quad
      O\leadsto X\leadsto O'.
  \end{aligned}
  }
\]

The prime is the entire point of the third line.

% -----------------------------------------------------------------------------
%  3.11  PROVENANCE AS TOPOLOGY RATHER THAN LABEL
% -----------------------------------------------------------------------------

\subsection{Provenance Is Not Merely a Label}
\label{sec:provenance-topology}

It is tempting to treat provenance as metadata attached to a state:

\[
  (x,p_x),
\]

where \(p_x\) might record an author, source, date, or generation number.

The provenance graph shows why this representation can be insufficient.

What matters may not be a label carried by \(x\) but the structure of
relations connecting \(x\) to other states.

If

\[
  x
\]

has two parents,

\[
  a\to x,
  \qquad
  b\to x,
\]

and three descendants,

\[
  x\to c_1,
  \qquad
  x\to c_2,
  \qquad
  x\to c_3,
\]

then its provenance is partly constituted by this position within a network
of descent.

A scalar label cannot, in general, preserve that structure.

This suggests a stronger conception:

\begin{definition}[Relational provenance]
\label{def:relational-provenance}
The \emph{relational provenance} of a state \(x\) relative to a provenance
graph \(\graphG\) is the derivational structure connecting \(x\) to the
relevant ancestors, descendants, transformations, and competing paths
represented in \(\graphG\).
\end{definition}

Provenance is therefore not necessarily something a state \emph{has} in the
same sense that it has a checksum.

It may be something a state occupies.

% -----------------------------------------------------------------------------
%  3.12  THE BEGINNING OF WORLDHOOD
% -----------------------------------------------------------------------------

\subsection{Return to Which World?}
\label{sec:return-world}

The wall now returns in a different form.

If Anarres and Urras were merely two locations in one homogeneous world, then
Shevek's movement could be described primarily as transportation. But the
novel insists upon differences in property, obligation, institutional
organization, language, sexuality, scientific practice, scarcity, abundance,
and political legitimacy \cite{leguin1974}.

The journey crosses between differently organized fields of affordance.

What may be owned, exchanged, refused, published, commanded, or withheld
changes across the boundary.

This suggests that a world cannot be identified merely with a spatial region.

A world, in the sense relevant to the No-World argument, is closer to an
organized field of relations within which distinctions and actions become
available.

The question

\[
  \text{Did Shevek return to Anarres?}
\]

is therefore easier than the question

\[
  \text{Did Shevek return to the same world?}
\]

The first can be answered geographically.

The second requires a theory of which relational structures constitute
world-continuation.

We still do not possess that theory.

But the question has now appeared twice: prospectively in
Movement~\ref{mov:continuation}, and recurrently here.

Its recurrence is evidence that it is not incidental.

% -----------------------------------------------------------------------------
%  3.13  WHAT MOVEMENT III HAS FORCED
% -----------------------------------------------------------------------------

\subsection{What Return Has Broken}
\label{sec:return-forced}

Movement~\ref{mov:origin} began with states and was forced to introduce
trajectories:

\[
  x
  \rightsquigarrow
  \gamma.
\]

Movement~\ref{mov:continuation} showed that trajectories require constraints
on admissible transformation:

\[
  \gamma
  \rightsquigarrow
  (\gamma,\Adm).
\]

The present movement has shown that a single trajectory can itself be
insufficient.

Branching transmission and convergent descent require

\[
  \gamma
  \rightsquigarrow
  \graphG.
\]

Return then shows that recurrence within \(\graphG\) cannot be represented as
simple restoration:

\[
  O
  \leadsto
  X
  \leadsto
  O',
  \qquad
  O'\neq O.
\]

Finally, the discovery that origins themselves possess provenance prevents
the left boundary from being treated as metaphysically privileged merely
because a particular inquiry begins there.

The representational object has therefore changed again.

\begin{held}[Linear-History Insufficiency]
\centering
A trajectory can preserve the history of a path.\\[3pt]
It cannot, by itself, preserve the history of a branching provenance
structure.
\end{held}

The three novels have now supplied the basic operations required by the
argument.

No fourth literary work is needed to add another temporal direction.

What remains is to break the machinery we have constructed.

The next movement therefore changes method. It will not proceed primarily by
novel, chronology, or theory. It will proceed by pathological cases.

Each case will ask whether two objects that one part of the formalism treats
as equivalent can nevertheless be distinguished by another part.

The order matters.

We begin with the easiest deception: a perfect copy with the wrong history.

We end with the harder case: two present states for which every selected
state-level test agrees while their histories still refuse to become one.

Only after those failures have accumulated will it be possible to state the
Originist objection in its stronger form.

% ============================================================================
%  MOVEMENT IV
%  PATHOLOGIES
% ============================================================================

\section{Pathologies}
\label{mov:pathologies}

\movementquestion{Which distinctions disappear when histories are replaced by their endpoints?}

\begin{orientation}
The preceding movements built a vocabulary by following three different
problems. Card required a distinction between a surviving state and the
history that produced it. Zebrowski required a distinction between material
persistence and continuation through transformation. Le Guin required a
distinction between a single trajectory and a branching structure of
departure, transmission, convergence, and return.

It would now be easy to treat these distinctions as a completed theory. That
would be premature.

A formal distinction earns its place only when removing it causes cases that
matter to collapse together. The purpose of the present movement is therefore
destructive. Instead of asking what the formalism can describe, we ask where
simpler descriptions fail.

The cases become progressively less forgiving. A forgery shows that perfect
resemblance does not establish descent. Independent rediscovery shows that
structural agreement does not require copying at all. Compression and
reconstruction show that information can disappear and later reappear without
thereby acquiring an unbroken provenance. Gradual replacement separates
material persistence from historical continuation. Branching and convergence
show that generation count alone cannot characterize provenance. Restoration
shows that successful repair does not reverse the history of damage.
Civilizational migration asks whether continuation of an organization is
enough to preserve its world.

The final case removes the last visual aid. Two states will agree under every
state-level test selected for the problem. Nothing in their present form will
tell them apart.

Their histories will still differ.

That case is not an exotic exception to the argument. It is the case for
which the argument has been preparing.
\end{orientation}

% -----------------------------------------------------------------------------
%  4.1  PERFECT FORGERY
% -----------------------------------------------------------------------------

\subsection{Perfect Forgery}
\label{sec:perfect-forgery}

Consider an original artifact \(O\) and a forgery \(F\).

Ordinary examples make the distinction easy because forgeries contain
imperfections. Pigments are anachronistic. Handwriting differs. Metadata is
wrong. A cryptographic hash fails. Some sufficiently discriminating
state-level test reveals that

\[
  O\neq F.
\]

Those imperfections are irrelevant to the present case.

Assume instead that the forgery is perfect with respect to every state-level
test in the selected family \(\mathcal S\). Write

\[
  O\equiv_{\mathcal S}F.
\]

If the artifacts are digital, we may choose the particularly strong case

\[
  O\eqbyte F.
\]

The forgery is nevertheless a forgery because of how it was produced.

Let

\[
  \gamma_O
\]

denote the history of the original and

\[
  \gamma_F
\]

the history by which the forgery was constructed. Then

\[
  \End(\gamma_O)\equiv_{\mathcal S}\End(\gamma_F)
\]

while

\[
  \gamma_O\neq\gamma_F.
\]

\begin{pathology}[Perfect forgery]
\label{path:perfect-forgery}
For any family of state-level tests \(\mathcal S\) that does not encode
derivational history, agreement under every test in \(\mathcal S\) does not
in general imply common provenance.
\end{pathology}

The qualification concerning derivational history is essential. One can
always define a purportedly state-level property that secretly includes
provenance. That maneuver does not solve the problem. It merely moves history
into the definition of state.

The forgery therefore supplies the first failed implication:

\[
  x\eqbyte y
  \nimplies
  x\eqprov y.
\]

This failure is stronger than the claim that appearances can deceive.

There need be no state-level deception at all.

The state may be exact.

The genealogy may still be false.

% -----------------------------------------------------------------------------
%  4.2  INDEPENDENT REDISCOVERY
% -----------------------------------------------------------------------------

\subsection{Independent Rediscovery}
\label{sec:independent-rediscovery}

Forgery retains an important relation to the original: the forger attempts to
reproduce it.

Independent rediscovery removes even that relation.

Suppose two researchers, \(A\) and \(B\), work from disjoint evidential
collections and arrive independently at the same proposition

\[
  \varphi.
\]

The resulting inscriptions may be semantically equivalent:

\[
  \varphi_A\simI\varphi_B.
\]

They may even be textually identical.

Yet neither result need derive from the other.

Card's account of Leyel Forska and the independently pursued work associated
with Thoren Magolissian provides a literary instance of this broader pattern:
separate research trajectories may converge upon overlapping questions or
results without collapsing into a single derivational chain
\cite{card1989}.

The provenance graph has the form

\[
\begin{tikzcd}[column sep=4em,row sep=1.7em]
E_A \arrow[r] & \varphi_A \\
E_B \arrow[r] & \varphi_B
\end{tikzcd}
\]

rather than

\[
  \varphi_A\longrightarrow\varphi_B.
\]

\begin{pathology}[Independent rediscovery]
\label{path:rediscovery}
Invariant or semantic equivalence between two states does not imply
genealogical dependence between them.
\end{pathology}

Thus

\[
  x\simI y
  \nimplies
  x\leadcausal y.
\]

This failed implication matters because similarity is often treated as
evidence of copying. Sometimes it is excellent evidence. It is never, by
itself, identical with copying.

The Originist therefore distinguishes convergence by descent from convergence
without descent.

The endpoint does not.

% -----------------------------------------------------------------------------
%  4.3  LOSSY COMPRESSION AND RECONSTRUCTION
% -----------------------------------------------------------------------------

\subsection{Lossy Compression Followed by Reconstruction}
\label{sec:lossy-reconstruction}

The next case introduces an operation central to the computational motivation
for the Originist role.

Let \(C\) be a lossy compression operation and \(R\) a reconstruction
operation. Begin with an artifact \(x\):

\[
  x
  \xrightarrow{C}
  c
  \xrightarrow{R}
  \widehat x.
\]

Because \(C\) is lossy, there exist distinct inputs \(x_1\) and \(x_2\) such
that

\[
  C(x_1)=C(x_2).
\]

Consequently, \(C(x)\) does not in general determine \(x\) uniquely.

Reconstruction must therefore use something beyond information preserved
injectively by the compressed representation. It may exploit priors,
regularities, external context, a learned generative model, or constraints on
the class of plausible originals.

Suppose nevertheless that reconstruction succeeds spectacularly:

\[
  \widehat x\eqbyte x.
\]

The present state has returned exactly.

The lost trajectory has not.

\begin{proposition}[Reconstruction does not reverse loss]
\label{prop:reconstruction-not-inverse}
Let \(C:X\rightarrow Y\) be non-injective and let \(R:Y\rightarrow X\) be a
reconstruction operation. Even when

\[
  R(C(x))=x
\]

for a particular \(x\), it does not follow that \(R=C^{-1}\).
\end{proposition}

\begin{proof}
Since \(C\) is non-injective, there exist \(x_1\neq x_2\) such that

\[
  C(x_1)=C(x_2)=y.
\]

If \(R\) is a function, then \(R(y)\) has one value. It cannot simultaneously
satisfy

\[
  R(y)=x_1
  \qquad\text{and}\qquad
  R(y)=x_2.
\]

Therefore \(R\) is not a two-sided inverse of \(C\) on the relevant domain.
The fact that \(R(C(x))=x\) for some particular \(x\) establishes successful
reconstruction in that case, not reversal of the information loss induced by
\(C\).
\end{proof}

This distinction is crucial.

A reconstruction can restore content without restoring provenance.

The sequence

\[
  x
  \xrightarrow{C}
  c
  \xrightarrow{R}
  x
\]

does not reduce historically to

\[
  x.
\]

The endpoint is idempotent with respect to the selected state comparison.

The history is not.

This is the formal situation that motivates generation tracking in repeated
compression--regeneration systems such as the proposed Originist role in the
compression and repair pipeline considered by the present project. If content
passes through

\[
  R_1
  \circ
  C_1
  \circ
  R_2
  \circ
  C_2
\]

and returns to a state indistinguishable from its source, a state-only loss
function may report perfect recovery.

An Originist ledger reports that something happened.

% -----------------------------------------------------------------------------
%  4.4  THE SHIP AFTER THE SHIP
% -----------------------------------------------------------------------------

\subsection{Ship-of-Theseus Replacement}
\label{sec:theseus-pathology}

Movement~\ref{mov:continuation} used gradual replacement to separate material
identity from continuous transformation. We can now sharpen the case.

Let

\[
  S_0
  \to
  S_1
  \to
  \cdots
  \to
  S_n
\]

be a gradual replacement trajectory satisfying an admissibility criterion
\(\Adm\).

Suppose

\[
  M(S_0)\cap M(S_n)=\varnothing.
\]

Now collect the removed components of \(S_0\) and reconstruct a second ship
\(S^\ast\).

The familiar puzzle asks which ship is the original.

The Originist formulation asks a prior question.

What histories are being compared?

We have

\[
  \gamma_1:
  S_0\to S_1\to\cdots\to S_n,
\]

and a second trajectory involving removal, preservation, and reassembly:

\[
  \gamma_2:
  M(S_0)
  \to
  \cdots
  \to
  S^\ast.
\]

The question

\[
  S_n=S^\ast?
\]

is therefore downstream from a richer structure.

The two endpoints inherit different relations to \(S_0\).

\begin{pathology}[Competing continuities]
\label{path:competing-continuities}
A single earlier state may support multiple later claims of continuity under
different preserved relations, without those claims being reducible to
endpoint similarity.
\end{pathology}

This pathology prevents admissibility from being treated as a single universal
predicate.

The gradual ship may preserve operational continuity.

The reconstructed ship may preserve original material.

A legal system might privilege one relation.

A museum might privilege another.

A theory of personal identity might privilege neither.

The formalism must therefore retain the type of continuation being claimed.

% -----------------------------------------------------------------------------
%  4.5  BRANCHING COPY
% -----------------------------------------------------------------------------

\subsection{Branching Copy}
\label{sec:branching-copy}

Suppose a state \(x\) is copied twice:

\[
\begin{tikzcd}[column sep=4em,row sep=1.6em]
& x_A \\
x \arrow[ur,"C_A"] \arrow[dr,"C_B"'] & \\
& x_B
\end{tikzcd}
\]

and suppose both copying operations are exact:

\[
  x_A\eqbyte x
  \qquad\text{and}\qquad
  x_B\eqbyte x.
\]

Then

\[
  x_A\eqbyte x_B.
\]

Yet the copies occupy distinct branches of the provenance graph.

If each branch subsequently changes,

\[
  x_A\to y_A,
  \qquad
  x_B\to y_B,
\]

the common origin does not make the descendants one history.

This matters for the language of continuation. If \(x\) is an agent,
institution, document tradition, or civilization, branching makes the phrase
\enquote{the continuation} ambiguous.

There may be several descendants.

\begin{proposition}[Branching destroys unique succession]
\label{prop:branching-succession}
If a provenance state \(x\) has two distinct admissible descendants \(y_1\)
and \(y_2\) neither of which is derivationally downstream from the other, then
descent from \(x\) does not determine a unique successor.
\end{proposition}

\begin{proof}
Both \(y_1\) and \(y_2\) satisfy the stipulated descent relation from \(x\).
By assumption neither is downstream from the other. Hence the provenance
structure contains no unique descendant selected solely by the fact of descent
from \(x\).
\end{proof}

This is the branching problem already suggested by Shevek's dissemination of
knowledge. Once transmission is plural, genealogy ceases to guarantee
singularity.

% -----------------------------------------------------------------------------
%  4.6  EQUAL GENERATION COUNTS, DIFFERENT HISTORIES
% -----------------------------------------------------------------------------

\subsection{Why Generation Count Is Not Enough}
\label{sec:generation-not-enough}

The provenance notation introduced earlier has so far allowed a statement to
carry a generation number:

\[
  \varphi^{[0]},
  \qquad
  \varphi^{[1]},
  \qquad
  \varphi^{[2]}.
\]

Branching exposes a limitation.

Consider two three-step histories:

\[
  \gamma_A:
  O
  \xrightarrow{R}
  a_1
  \xrightarrow{C}
  a_2
  \xrightarrow{R}
  x_A,
\]

and

\[
  \gamma_B:
  O
  \xrightarrow{C}
  b_1
  \xrightarrow{C}
  b_2
  \xrightarrow{C}
  x_B.
\]

Both endpoints have generation depth three:

\[
  \Depth_O(x_A;\gamma_A)
  =
  \Depth_O(x_B;\gamma_B)
  =
  3.
\]

Yet the transformation histories differ:

\[
  R\circ C\circ R
  \neq
  C\circ C\circ C.
\]

A scalar depth therefore collapses trajectories that an Originist may need to
distinguish.

This is the point at which the fuller provenance notation becomes necessary.

We may write

\[
  \varphi^{[3:R\circ C\circ R]}
\]

rather than merely

\[
  \varphi^{[3]}.
\]

If the transformation history is known only partially, additional fields may
record the invariant class preserved and an evidential support value:

\[
  \varphi^{[3:R\circ C\circ R;\,\Inv^{+};\,p=.71]}.
\]

The notation is intentionally unpleasant.

The unpleasantness records a real increase in the dimensionality of the
problem.

\begin{pathology}[Scalar provenance insufficiency]
\label{path:scalar-provenance}
Equal generation depth does not imply equivalent derivational history.
Consequently, a scalar generation count is insufficient whenever
transformation type matters to the inference being performed.
\end{pathology}

We have now learned that provenance is not adequately represented by a state,
a path length, or even necessarily a single path.

% -----------------------------------------------------------------------------
%  4.7  CONVERGENT DESCENDANTS
% -----------------------------------------------------------------------------

\subsection{Convergent Descendants}
\label{sec:convergent-descendants}

Branching separates descendants.

Convergence brings them together again.

Suppose

\[
  O\to A\to A'
\]

and

\[
  O\to B\to B'
\]

form distinct branches. Later transformations produce

\[
  A'\to Z
\]

and

\[
  B'\to Z.
\]

The state \(Z\) now has more than one relevant derivational path from \(O\).

The provenance graph contains

\[
\begin{tikzcd}[column sep=3.5em,row sep=1.6em]
& A \arrow[r] & A' \arrow[dr] & \\
O \arrow[ur] \arrow[dr] & & & Z \\
& B \arrow[r] & B' \arrow[ur] &
\end{tikzcd}
\]

The question

\[
  \Depth_O(Z)=?
\]

is no longer well formed without specifying a path or aggregation rule.

More importantly, convergence can make later states appear to erase earlier
branching.

If one observes only \(Z\), the two histories have become invisible.

They have not ceased to have occurred.

\begin{pathology}[Convergence hides branching]
\label{path:convergence-hides}
A provenance graph may contain historically relevant branching that is not
recoverable from a convergent descendant considered in isolation.
\end{pathology}

This is the graph-theoretic generalization of the first example in
Movement~\ref{mov:origin}.

There,

\[
  x_A=x_B
\]

concealed different paths.

Here a single state \(Z\) conceals an entire branching structure.

The Originist problem is becoming less like attaching metadata to objects and
more like preserving a topology of descent.

% -----------------------------------------------------------------------------
%  4.8  RESTORATION FROM INCOMPLETE ARCHIVES
% -----------------------------------------------------------------------------

\subsection{Restoration from Incomplete Archives}
\label{sec:restoration}

Suppose an original document \(O\) is damaged, leaving a surviving fragment
\(F\):

\[
  O
  \xrightarrow{D}
  F.
\]

An editor later reconstructs the missing material using independent evidence
\(E\):

\[
  (F,E)
  \xrightarrow{R}
  \widehat O.
\]

Assume the reconstruction is correct:

\[
  \widehat O\eqbyte O.
\]

What has been restored?

The content has been restored under byte identity.

The history has not.

The original reached its state without passing through destruction and
editorial reconstruction. The restored object did.

Thus

\[
  O\eqbyte\widehat O
\]

while

\[
  \Prov(O)\neq\Prov(\widehat O).
\]

This example is important because no error is required.

The reconstruction may be perfect.

Provenance difference is not a synonym for corruption.

\begin{proposition}[Perfect restoration preserves a wound]
\label{prop:restoration-wound}
If a transformation destroys provenance-relevant information and a later
reconstruction restores the endpoint state using additional information, then
perfect state restoration does not erase the historical fact of the
intervening loss.
\end{proposition}

\begin{proof}
Let

\[
  O\xrightarrow{D}F
\]

destroy a distinction required to recover \(O\) uniquely from \(F\). Let
external evidence \(E\) permit a reconstruction

\[
  R(F,E)=O.
\]

The final state agrees with the original. Nevertheless, the realized history
of the reconstructed state contains \(D\), the loss of information, the
introduction of \(E\), and the reconstruction operation \(R\). The original
history does not contain these events. Therefore restoration of state does not
erase the historical distinction.
\end{proof}

This is the first case for which the language of a wound becomes useful.

A wound is not necessarily visible in the healed state.

% -----------------------------------------------------------------------------
%  4.9  WOUND
% -----------------------------------------------------------------------------

\subsection{The Wound}
\label{sec:wound}

Let

\[
  q:\Gamma\longrightarrow\Gamma/{\sim}
\]

be an operation that identifies histories under an equivalence relation
\(\sim\).

Suppose two histories \(\gamma_1\) and \(\gamma_2\) differ with respect to a
provenance distinction \(\delta\), but

\[
  q(\gamma_1)=q(\gamma_2).
\]

After quotienting, \(\delta\) cannot be recovered from the quotient state
alone.

We call this an Originist wound relative to \(\delta\).

\begin{definition}[Originist loss]
\label{def:originist-loss}
For a quotient or information-reducing operation \(q\) and a
provenance-relevant distinction \(\delta\), the Originist loss

\[
  \lossorig(q;\delta)
\]

denotes the loss of discriminative capacity concerning \(\delta\) induced by
applying \(q\).
\end{definition}

No scalar metric is assumed yet. The notation may denote a logical loss
rather than a numerical magnitude.

The minimal condition is

\[
  \lossorig(q;\delta)>0
\]

whenever there exist \(\gamma_1,\gamma_2\) such that

\[
  \delta(\gamma_1)\neq\delta(\gamma_2)
\]

but

\[
  q(\gamma_1)=q(\gamma_2).
\]

\begin{woundbox}
\begin{wounded}
A wound occurs when an operation makes a distinction unavailable that a later
inference may require.

The wound is not identical with damage to the endpoint.

A perfectly restored endpoint may still stand downstream from an irreversible
loss of provenance.
\end{wounded}
\end{woundbox}

This definition explains why compression is not intrinsically objectionable.

Every useful representation forgets something.

The Originist question is whether the forgotten distinction belongs to the
future inferential task.

% -----------------------------------------------------------------------------
%  4.10  CIVILIZATION ON ANOTHER SUBSTRATE
% -----------------------------------------------------------------------------

\subsection{A Civilization Reproduced Elsewhere}
\label{sec:civilization-substrate}

We can now return to the Macrolife limit with more machinery available.

Let a civilization \(C_0\) inhabit an environment \(W_0\). Over time it
constructs new habitats, changes its material substrate, modifies its
institutions, and transmits its organizational structures into descendants:

\[
  (C_0,W_0)
  \longrightarrow
  (C_1,W_1)
  \longrightarrow
  \cdots
  \longrightarrow
  (C_n,W_n).
\]

Suppose the civilizational trajectory satisfies an admissibility criterion

\[
  \gamma_C\in\Adm_C.
\]

We may therefore have reason to assert

\[
  C_0\leadcausal C_n
\]

and perhaps even to regard \(C_n\) as an admissible continuation of \(C_0\).

It does not follow that

\[
  W_0\eqworld W_n.
\]

The distinction becomes especially important when the environment is no longer
a passive background but is itself constructed and transformed by the
civilization.

A civilization may continue while changing worlds.

Alternatively, what appears to be a new world may itself be part of the
continuation of an older one.

The question cannot be settled by civilizational genealogy alone.

\begin{pathology}[Civilizational continuity without established world-continuity]
\label{path:civilization-world}
The existence of an admissible civilizational trajectory

\[
  C_0\rightsquigarrow C_n
\]

does not by itself establish

\[
  W_0\eqworld W_n.
\]
\end{pathology}

We have now reached the distinction that the earlier movements repeatedly
approached without formalizing.

Persistence of an object, persistence of a civilization, and persistence of
a world are separate claims.

% -----------------------------------------------------------------------------
%  4.11  WHEN EVERY STATE-LEVEL TEST AGREES
% -----------------------------------------------------------------------------

\subsection{The Case with Nothing Left to See}
\label{sec:nothing-left-to-see}

We can now construct the limiting pathology.

Let \(\mathcal S\) be the complete family of state-level tests selected for a
particular task. Suppose two states \(x_A\) and \(x_B\) satisfy

\[
  s(x_A)=s(x_B)
  \qquad
  \text{for every }s\in\mathcal S.
\]

Write this as

\[
  x_A\equiv_{\mathcal S}x_B.
\]

Suppose nevertheless that

\[
  \gamma_A\neq\gamma_B.
\]

Further suppose that the difference between \(\gamma_A\) and \(\gamma_B\)
includes a provenance distinction \(\delta\) relevant to a later operation.

No additional test in \(\mathcal S\) can recover \(\delta\), because the
states agree under every test in the selected family.

The only remaining source of the distinction is historical information not
encoded by those state tests.

\begin{theorem}[State-level underdetermination of provenance]
\label{thm:state-underdetermination}
Let \(\mathcal S\) be a family of functions on a state space \(X\), and define

\[
  x\equiv_{\mathcal S}y
  \quad\Longleftrightarrow\quad
  s(x)=s(y)
  \text{ for every }s\in\mathcal S.
\]

Let \(\Gamma(X)\) be a space of derivational histories with endpoint map

\[
  \End:\Gamma(X)\to X.
\]

If there exist distinct histories
\(\gamma_A,\gamma_B\in\Gamma(X)\) such that

\[
  \End(\gamma_A)\equiv_{\mathcal S}\End(\gamma_B),
\]

then no inference using only values of tests in \(\mathcal S\) can, in
general, distinguish \(\gamma_A\) from \(\gamma_B\).
\end{theorem}

\begin{proof}
Let

\[
  x_A=\End(\gamma_A),
  \qquad
  x_B=\End(\gamma_B).
\]

By hypothesis,

\[
  s(x_A)=s(x_B)
\]

for every \(s\in\mathcal S\). Therefore any inference whose input consists
only of the tuple

\[
  \bigl(s(x)\bigr)_{s\in\mathcal S}
\]

receives identical input for \(x_A\) and \(x_B\). Such an inference cannot
deterministically discriminate the two histories on the basis of those inputs
alone. Any successful discrimination must therefore use information not
contained in the selected state-level observations.
\end{proof}

This theorem is elementary enough that its proof could almost be omitted.

It is retained because the triviality is revealing.

Once the problem is stated correctly, the impossibility is immediate.

A state cannot disclose a distinction that the representation of that state
has erased.

% -----------------------------------------------------------------------------
%  4.12  THE ORIGINIST OBJECTION, STRONG FORM
% -----------------------------------------------------------------------------

\subsection{The Originist Objection}
\label{sec:originist-objection-strong}

The examples now permit the preliminary objection from
\cref{obj:preliminary-originist} to be stated more sharply.

\begin{objection}[Strong form]
\label{obj:strong-originist}
Let \(q\) be an abstraction, quotient, compression, canonicalization, or
equivalence operation applied to states or histories. The operation is
epistemically unsafe for an Originist task whenever there exists a
provenance-relevant distinction \(\delta\) such that

\[
  \lossorig(q;\delta)>0
\]

and a subsequent inference requires \(\delta\).

Equivalence is therefore not objectionable because it forgets.

It is objectionable when it forgets a distinction and subsequently behaves as
though no distinction was lost.
\end{objection}

This is the Originist objection to premature quotienting.

It does not demand preservation of all history.

Such preservation would often be impossible and almost always be useless.

It demands an account of the wound.

\begin{held}[The Originist Objection]
\centering
Do not ask only what an equivalence preserves.\\[3pt]
Ask which future questions become unanswerable after the equivalence is
applied.
\end{held}

The shift is subtle but consequential.

The usual mathematical question is often

\[
  \text{Which differences may safely be ignored?}
\]

The Originist adds

\[
  \text{Safe for which later operation?}
\]

Admissibility therefore applies not only to transformations of objects.

It applies to transformations of representations.

% -----------------------------------------------------------------------------
%  4.13  THE CITATION THAT DID NOT CHANGE
% -----------------------------------------------------------------------------

\subsection{The Citation That Did Not Change}
\label{sec:citation-wound}

We can now return to a small piece of the present document.

In Movement~\ref{mov:origin}, a claim concerning Leyel Forska's orientation
toward principles of origin rather than merely the identification of a
planetary coordinate was presented with a provenance-bearing note. At that
point it was marked as directly grounded.

Subsequent inspection of the derivational chain supporting the wording used
there does not establish a sufficiently direct path from the present
formulation to the primary text to justify the generation-zero classification.

The proposition itself need not be withdrawn.

Its truth value has not thereby changed.

Its wording has not changed.

The historical source from which the underlying material ultimately derives
has not changed.

What has changed is the present document's warranted description of its own
relation to that source.

The earlier provenance classification must therefore be retyped:

\[
  \varphi^{[0]}
  \rightsquigarrow
  \varphi\provomega.
\]

This is the unique unknown-depth wound in the present manuscript.

It is deliberately not repaired by assigning an arbitrary finite generation
count.

The available evidence does not warrant one.

\begin{woundbox}
\begin{wounded}
The proposition did not change.

Its history did not change.

What changed was the document's knowledge of its history.

The correction therefore belongs neither to the proposition's semantic
content nor directly to its truth value. It belongs to its epistemic
provenance type.
\end{wounded}
\end{woundbox}

The distinction can now be written

\[
  \Truth(\varphi)
  \neq
  \Prov(\varphi).
\]

More precisely, revision of

\[
  \Prov(\varphi)
\]

does not entail revision of

\[
  \Truth(\varphi).
\]

Nor does the converse hold.

A proposition may be discovered false while its provenance remains perfectly
known.

Truth and provenance are independent coordinates.

This is why provenance cannot be reduced to confidence.

% -----------------------------------------------------------------------------
%  4.14  PROVENANCE AS AN EPISTEMIC TYPE
% -----------------------------------------------------------------------------

\subsection{Provenance as an Epistemic Type}
\label{sec:epistemic-type}

The preceding case suggests a stronger interpretation of the superscript
notation.

Consider

\[
  \varphi^{[0]}
\]

and

\[
  \varphi^{[3:R\circ C\circ R]}.
\]

Suppose the semantic content is identical.

A conventional representation might treat the two propositions as the same
statement accompanied by different metadata.

For Originist reasoning, that description is sometimes too weak.

The provenance class can determine which operations are admissible.

A directly inspected quotation may support a claim about exact wording.

A reconstructed paraphrase may support the semantic claim while failing to
support an inference about lexical choice.

An uncertain derivational chain may support exploratory reasoning while being
inadequate for a claim whose force depends upon primary-source wording.

Thus provenance can alter what one is entitled to infer even when semantic
content is held fixed.

\begin{definition}[Provenance-sensitive epistemic type]
\label{def:epistemic-type}
A proposition has a \emph{provenance-sensitive epistemic type} when its
derivational history constrains the class of inferential operations for which
it is admissible evidence.
\end{definition}

Under this interpretation,

\[
  \varphi^{[0]}
\]

and

\[
  \varphi^{[3:R\circ C\circ R]}
\]

may be semantically equivalent while remaining epistemically non-substitutable.

The superscript is no longer decorative metadata.

It participates in the admissibility structure.

This closes the circuit between provenance and TARTAN-like constraint
reasoning: provenance does not merely tell us where a representation came
from. It constrains what transformations of that representation remain
licensed.

% -----------------------------------------------------------------------------
%  4.15  THE FAILED IMPLICATIONS
% -----------------------------------------------------------------------------

\subsection{An Ecology of Failed Implications}
\label{sec:failed-implications}

The pathologies can now be condensed without replacing their worked examples.

Perfect forgery established

\[
  x\eqbyte y
  \nimplies
  x\eqprov y.
\]

Independent rediscovery established

\[
  x\simI y
  \nimplies
  x\leadcausal y.
\]

The corruption, ruin, and pathological-growth cases established

\[
  x\leadcausal y
  \nimplies
  x\text{ is admissibly continued by }y.
\]

The Macrolife limit established that even admissible civilizational
continuation does not automatically give

\[
  W_t\eqworld W_{t+n}.
\]

Generation analysis established

\[
  \Depth_O(x;\gamma_A)
  =
  \Depth_O(y;\gamma_B)
  \nimplies
  \gamma_A\sim_{\mathrm{gen}}\gamma_B
\]

when \(\sim_{\mathrm{gen}}\) is intended to encode more than equal path length.

Restoration established

\[
  x_{\mathrm{restored}}\eqbyte x_{\mathrm{original}}
  \nimplies
  \Hist(x_{\mathrm{restored}})
  =
  \Hist(x_{\mathrm{original}}).
\]

Return established

\[
  O'\simconstraint O
  \nimplies
  O'=O.
\]

No single one of these failures is especially mysterious.

Their accumulation is the result.

Relations that ordinary language repeatedly compresses into \enquote{same}
have separated under pressure.

% -----------------------------------------------------------------------------
%  4.16  WHAT MOVEMENT IV HAS FORCED
% -----------------------------------------------------------------------------

\subsection{What the Pathologies Leave Standing}
\label{sec:pathologies-forced}

The destructive work is now sufficient.

We began the essay with a state

\[
  x.
\]

Endpoint insufficiency forced the trajectory

\[
  \gamma.
\]

Branching and convergence forced the provenance graph

\[
  \graphG.
\]

Causal insufficiency forced admissibility

\[
  \Adm.
\]

Repeated failures of state-level identity forced explicit invariants

\[
  \Inv.
\]

The civilizational cases repeatedly forced a question not yet answered by any
of them:

\[
  \World.
\]

The progression should not be written as a hierarchy in which each later
object contains the truth of the earlier one.

It is better represented as a sequence of failures:

\[
  \textsc{State}
  \xrightarrow{\text{insufficient}}
  \textsc{Trajectory}
  \xrightarrow{\text{insufficient}}
  \textsc{Provenance}
  \xrightarrow{\text{insufficient}}
  \textsc{Admissibility}
  \xrightarrow{\text{?}}
  \textsc{World}.
\]

The question mark is essential.

Nothing so far establishes that worldhood is simply the next layer in a
well-ordered construction.

It may instead be the point at which the demand for a single underlying
object becomes unstable.

That possibility is the subject of the final movement.

\begin{held}[After the Pathologies]
\centering
State equality does not determine historical equality.\\
Historical descent does not determine admissible continuation.\\
Admissible continuation does not yet determine continuation of a world.
\end{held}

The three literary probes can therefore return one final time.

Card asked how a lost world might be reconstructed from traces.

Zebrowski asked how a civilization might propagate beyond the material world
from which it arose.

Le Guin asked what it means to leave a world, cross its boundaries, transmit
something beyond them, and return to an origin that has continued changing in
one's absence.

These questions no longer require three separate formalisms.

They now occupy different regions of one unresolved problem.

What is it that persists when the state does not?

What is it that returns when the origin has changed?

What is it that has been preserved when a civilization survives but the
relations constituting its world may not?

Movement~V will not introduce another pathology merely to enlarge the
apparatus again.

It will ask whether the apparatus already earned is sufficient to say what
descent means.

% ============================================================================
%  MOVEMENT V
%  DESCENT
% ============================================================================

\section{Descent}
\label{mov:descent}

\movementquestion{If neither matter nor state is fundamental to continuity, what exactly descends?}

\begin{orientation}
The word \enquote{descent} usually sounds biological. Something descends from
something else because reproduction connects them. The parent need not remain
alive for the descendant to exist, and the descendant need not resemble the
parent in every respect. What matters is that there is a history connecting
them.

The same word becomes useful outside biology once the preceding cases are
placed together. A manuscript can descend from another manuscript. A
reconstruction can descend from fragments. A scientific idea can descend
through several translations. A civilization can descend from a civilization
whose material conditions it no longer shares. A political society can return
to an inherited problem without returning to an earlier state.

The difficult question is not whether these things change. They obviously do.
The question is what kind of historical structure makes some changes count as
continuation and others as replacement, corruption, forgery, rediscovery, or
the beginning of something else.

The previous movements did not discover a single substance that survives all
of these cases. They discovered instead a sequence of insufficiencies. States
were insufficient, so trajectories were retained. Trajectories were
insufficient, so branching provenance was retained. Causal provenance was
insufficient, so admissibility became necessary. Admissible continuation at
the level of an object or civilization still failed to answer whether a world
continued.

The temptation now is to name the final object and declare the ontology
finished.

The present movement resists that temptation.

The theories brought into contact here---RSVP, TARTAN, Chain of Memory,
Spherepop, recurrence, irreversibility, and the No-World argument---do not all
identify the same primitive. Their usefulness lies precisely in the fact that
they constrain different parts of the problem.

The synthesis is therefore not a reduction.

It is a geometry of descent.
\end{orientation}

% -----------------------------------------------------------------------------
%  5.1  A SEQUENCE OF INSUFFICIENCIES
% -----------------------------------------------------------------------------

\subsection{The Stack Is Not a Ladder}
\label{sec:stack-not-ladder}

The argument can now be summarized by the sequence

\[
  \State
  \;<\;
  \Trajectory
  \;<\;
  \Provenance
  \;<\;
  \Admissibility
  \;<\;
  \Worldhood.
\]

The symbol \(<\) must not be interpreted as an ontological ordering.

It means only that, for the questions considered in this essay, each term
became necessary after the preceding term failed to discriminate cases that
needed to remain distinct.

The progression is therefore better written

\[
  \State
  \xrightarrow{\mathsf F_1}
  \Trajectory
  \xrightarrow{\mathsf F_2}
  \Provenance
  \xrightarrow{\mathsf F_3}
  \Admissibility
  \xrightarrow{\mathsf F_4}
  \Worldhood,
\]

where each \(\mathsf F_i\) denotes a failure of sufficiency.

The first failure was endpoint convergence:

\[
  x_A=x_B,
  \qquad
  \gamma_A\neq\gamma_B.
\]

The second was linear-history insufficiency: a single path could not preserve
branching and convergence.

The third was causal insufficiency:

\[
  x\leadcausal y
\]

did not guarantee that \(y\) constituted an acceptable continuation of \(x\).

The fourth remains only partially formalized. A civilization may possess a
defensible admissible trajectory while the persistence of the world in which
its practices are meaningful remains undecided.

\begin{proposition}[Non-collapse of the descriptive stack]
\label{prop:noncollapse-stack}
No level of the sequence above may be replaced by the immediately preceding
level without collapsing at least one distinction established in
Movement~\ref{mov:pathologies}.
\end{proposition}

\begin{proof}
State alone collapses perfect forgery and convergent reconstruction when their
selected endpoint properties agree. A single trajectory collapses relevant
branching and convergence in a provenance graph. Causal provenance collapses
admissible continuation with corruption and other causally downstream but
constraint-violating transformations. Finally, admissible continuation of an
object or civilization does not, by itself, settle world-continuation, as
shown by \cref{path:civilization-world}. Hence each enlargement is motivated
by a distinction unavailable at the preceding level.
\end{proof}

The stack is therefore methodological rather than metaphysical.

It records where simpler descriptions broke.

% -----------------------------------------------------------------------------
%  5.2  RSVP: CONTINUITY WITHOUT A FIXED OBJECT
% -----------------------------------------------------------------------------

\subsection{RSVP and the Field Description of Descent}
\label{sec:rsvp-descent}

The Relativistic Scalar--Vector Plenum framework provides one way to loosen the
assumption that persistence must belong to a sharply bounded object.

In its characteristic notation, a configuration is described through a scalar
field \(\Phi\), a vector field \(\mathbf v\), and an entropy-like scalar field
\(S\):

\[
  \mathfrak R
  =
  (\Phi,\mathbf v,S).
\]

For present purposes, the importance of this representation is not any
particular field equation. It is that an entity may be represented as a
structured region within a changing field rather than as an indivisible bearer
of identity.

Let

\[
  \mathfrak R_t
  =
  (\Phi_t,\mathbf v_t,S_t)
\]

and

\[
  \mathfrak R_{t+\Delta t}
  =
  (\Phi_{t+\Delta t},
   \mathbf v_{t+\Delta t},
   S_{t+\Delta t})
\]

denote successive field configurations.

The question of continuation need not then take the form

\[
  x_t=x_{t+\Delta t}.
\]

It can instead ask whether a pattern of constraints, flows, boundaries, and
relations persists through the field evolution.

This is structurally compatible with the Macrolife case. Complete replacement
of material constituents need not terminate a history if the relevant
organization is continuously realized through changing substrate.

It is also compatible with the Originist case. A present configuration may
retain traces of transformations without containing a literal copy of its
earlier state.

The field representation therefore supplies a useful warning against
substantialism:

\begin{held}[Field Persistence]
\centering
Persistence need not require a substance that remains numerically unchanged.\\[3pt]
It may be realized by constrained transformation of a structured field.
\end{held}

This does not establish that every relevant notion of identity is reducible to
RSVP dynamics.

The stronger claim is unnecessary.

RSVP supplies a natural substrate on which trajectory-first reasoning can be
stated without presupposing the very state identity that the pathologies have
made doubtful.

% -----------------------------------------------------------------------------
%  5.3  TARTAN: ADMISSIBILITY IS NOT RESEMBLANCE
% -----------------------------------------------------------------------------

\subsection{TARTAN and the Geometry of Admissibility}
\label{sec:tartan-admissibility}

The transition from causal history to admissible history is where TARTAN
enters most directly.

Suppose the space of possible realizations is \(X\), and let

\[
  \mathcal A\subseteq X
\]

denote an admissible region determined by a collection of constraints.

A transformation

\[
  T:x\mapsto y
\]

is not licensed merely because \(y\) resembles \(x\), nor merely because
\(x\leadcausal y\).

The relevant question is whether the realized trajectory remains compatible
with the constraints defining the task.

For a trajectory

\[
  \gamma:[0,1]\rightarrow X,
\]

a local admissibility criterion might require

\[
  \gamma(t)\in\mathcal A
  \qquad
  \forall t\in[0,1].
\]

But the pathologies already showed why local membership may be insufficient.
Some constraints concern the history as a whole.

We therefore distinguish state admissibility from trajectory admissibility:

\[
  \mathcal A_X
  \subseteq
  X
\]

and

\[
  \mathcal A_\Gamma
  \subseteq
  \Gamma(X).
\]

It is possible that every state on a path satisfies the state-level
constraints while the path itself violates a provenance condition:

\[
  \gamma(t)\in\mathcal A_X
  \quad\forall t,
  \qquad
  \gamma\notin\mathcal A_\Gamma.
\]

A perfect forgery provides the simplest example.

The endpoint may be admissible as an artifact state.

The derivation may be inadmissible for a task requiring authentic descent.

\begin{remark}[Admissibility is constraint-relative, not derived]
\label{rem:admissibility-relativity}
Throughout what follows, admissibility is understood relative to a specified
family of constraints \(\mathcal C\). More explicitly, one may write
\(\mathcal A_{\mathcal C}\) for the admissible states or histories determined
relative to \(\mathcal C\), and \(\mathfrak A_{\mathcal C}\) for the
corresponding admissibility functional. When only one constraint family is
under consideration, the subscript will be suppressed and the shorter
notation \(\mathcal A\) or \(\mathfrak A\) retained. This suppression should
not be read as asserting a universal or independently derived notion of
admissibility: the relevant constraint family remains a parameter of the
construction. Every bare occurrence of \(\mathcal A\) or \(\mathfrak A\)
elsewhere in this document, including in the Path Category movement, is to
be read under this convention.
\end{remark}

\begin{definition}[Trajectory admissibility]
\label{def:trajectory-admissibility}
Let \(\Gamma(X)\) be a space of histories over \(X\). A trajectory-level
admissibility predicate is a map

\[
  \mathfrak A:\Gamma(X)\rightarrow\{0,1\}
\]

or, more generally, a graded functional

\[
  \mathfrak A:\Gamma(X)\rightarrow[0,1]
\]

whose value depends upon properties of the realized history not necessarily
recoverable from its endpoint.
\end{definition}

This is the formal location of the Originist inside a TARTAN-like system.

The Originist does not merely improve the representation of \(x\).

It supplies variables on which

\[
  \mathfrak A(\gamma_x)
\]

may depend.

% -----------------------------------------------------------------------------
%  5.4  CHAIN OF MEMORY
% -----------------------------------------------------------------------------

\subsection{Chain of Memory and the Preservation of Derivational Constraint}
\label{sec:com}

Chain of Memory can now be distinguished sharply from an ordinary store of
past states.

A memory system conceived as an archive may retain

\[
  x_0,x_1,\ldots,x_n.
\]

That is useful, but it does not necessarily preserve why one state followed
another.

A derivational memory requires at least

\[
  x_0
  \xrightarrow{T_1}
  x_1
  \xrightarrow{T_2}
  \cdots
  \xrightarrow{T_n}
  x_n.
\]

A branching memory requires still more:

\[
  \graphG=(V,E,\tau).
\]

Chain of Memory is therefore most naturally interpreted here not as a larger
context window but as persistence of constraints across transformations.

The distinction is important.

Remembering every prior token is not equivalent to remembering the derivation
that licensed the current state.

Conversely, a system may discard enormous quantities of state information
while preserving exactly the derivational relations required for later
reasoning.

This suggests a minimal Originist interpretation of CoM:

\[
  \mathrm{CoM}
  :
  \text{preserve the distinctions required for future admissibility tests}.
\]

The criterion is prospective.

A memory is valuable not because the past deserves exhaustive duplication,
but because future operations may require distinctions that would otherwise
be irrecoverable.

This is precisely the lesson of \(\lossorig\).

% -----------------------------------------------------------------------------
%  5.5  SPHEREPOP: CONTINUATION THROUGH EVENTS
% -----------------------------------------------------------------------------

\subsection{Spherepop and Event-Structured History}
\label{sec:spherepop}

Spherepop contributes a complementary vocabulary.

Where RSVP emphasizes fields and TARTAN emphasizes admissibility, Spherepop
treats histories through events of distinction, binding, refusal, and
collapse.

Let an event alphabet be written schematically as

\[
  \Sigma_{\mathrm{sp}}
  =
  \{
  \Pop,
  \Refuse,
  \Bind,
  \Collapse
  \}.
\]

A trajectory can then carry not only states but event types:

\[
  x_0
  \xrightarrow{\Pop}
  x_1
  \xrightarrow{\Bind}
  x_2
  \xrightarrow{\Refuse}
  x_3.
\]

This makes explicit something the scalar generation count could not.

Two histories of equal length can have different event structure.

Thus

\[
  \Depth_O(x;\gamma_1)
  =
  \Depth_O(y;\gamma_2)
\]

while

\[
  \tau(\gamma_1)\neq\tau(\gamma_2).
\]

The full provenance tag

\[
  [3:R\circ C\circ R]
\]

is a computationally mundane version of the same idea.

History has syntax.

The order and type of transformations matter.

If event composition is noncommutative, then in general

\[
  T_2\circ T_1
  \neq
  T_1\circ T_2.
\]

Consequently, even the multiset of transformations is insufficient.

The path retains order.

This provides another reason why provenance cannot be compressed into a
generation number without loss.

% -----------------------------------------------------------------------------
%  5.6  IRREVERSIBILITY AS HISTORICAL STRUCTURE
% -----------------------------------------------------------------------------

\subsection{Irreversibility Without Mystification}
\label{sec:irreversibility-synthesis}

Irreversibility has appeared repeatedly in the essay, but its role can now be
stated precisely.

Suppose

\[
  T:X\rightarrow Y
\]

identifies states that differ under a provenance-relevant distinction
\(\delta\). Then

\[
  \lossorig(T;\delta)>0.
\]

A later operation may successfully construct a state equivalent to the
earlier one:

\[
  R(T(x))\equiv_{\mathcal S}x.
\]

This does not imply historical reversal.

The sequence

\[
  x
  \xrightarrow{T}
  y
  \xrightarrow{R}
  x'
\]

has occurred.

Even if

\[
  x'\eqbyte x,
\]

the realized trajectory contains an event absent from the history terminating
at \(x\).

The irreversibility relevant here is therefore not simply the impossibility of
returning to a previous state.

A previous state may be restored exactly.

The irreversibility lies in the impossibility of making the restoration
retroactively become the original history.

\begin{proposition}[Historical irreversibility]
\label{prop:historical-irreversibility}
Let a realized history contain a nonidentity transformation \(T\). No later
state-level restoration operation can make the realized history identical to
a history in which \(T\) never occurred.
\end{proposition}

\begin{proof}
Let

\[
  \gamma_1=(x)
\]

denote the history terminating at the initial state prior to \(T\), and let

\[
  \gamma_2=(x\xrightarrow{T}y\xrightarrow{R}x')
\]

be the realized transformed-and-restored history.

Even if \(x'=x\), the ordered event sequence of \(\gamma_2\) contains \(T\)
and \(R\), whereas \(\gamma_1\) does not. Therefore

\[
  \gamma_1\neq\gamma_2.
\]

No property of the endpoint can alter the fact that these events belong to
one realized history and not the other.
\end{proof}

This is the simplest formal statement of the essay's temporal asymmetry.

Return can restore a state.

It cannot un-happen a history.

% -----------------------------------------------------------------------------
%  5.7  RECURRENCE UNDER CONSTRAINT
% -----------------------------------------------------------------------------

\subsection{The Constraint--Recurrence Principle}
\label{sec:constraint-recurrence}

The Le Guinian return can now be expressed without treating recurrence as
literal repetition.

Let

\[
  \gamma:
  x_0
  \leadsto
  x_1
  \leadsto
  \cdots
  \leadsto
  x_n.
\]

Suppose

\[
  x_n\sim_{\mathcal C}x_0
\]

for a selected family of constraints \(\mathcal C\), while

\[
  x_n\neq x_0.
\]

The trajectory has returned to a constraint-compatible region without
restoring the initial state.

This motivates the following formulation.

\begin{proposition}[Constraint--Recurrence Principle]
\label{prop:constraint-recurrence}
Recurrence relevant to historical continuation may be defined by return to an
admissible constraint class rather than by restoration of a prior state.
Accordingly,

\[
  x_n\sim_{\mathcal C}x_0
\]

may constitute recurrence even when

\[
  x_n\neq x_0
\]

and

\[
  \Hist(x_n)\neq\Hist(x_0).
\]
\end{proposition}

The proposition explains why return can be genuine without being reversible.

Shevek returns to Anarres.

The return is not false merely because neither Shevek nor Anarres is exactly
as before.

The difference is what makes the return historical.

% -----------------------------------------------------------------------------
%  5.8  WORLDHOOD AS A PROVENANCE PROBLEM
% -----------------------------------------------------------------------------

\subsection{A World Has a History Too}
\label{sec:world-provenance}

The No-World argument becomes relevant once worldhood is treated not as a
primitive container but as something reconstructed from relations,
distinctions, continuities, and constraints.

If a world were simply a fixed background \(W\), then the persistence question
would be easy to state:

\[
  x_t,x_{t+n}\in W.
\]

Both states occur in the same world because the container has been held fixed
by assumption.

The assumption is precisely what the present framework cannot make without
argument.

Suppose instead that worldhood at time \(t\) is realized through a relational
structure

\[
  W_t
  =
  \mathfrak W
  (
  X_t,
  B_t,
  \Adm_t,
  \Rel_t
  ),
\]

where \(X_t\) denotes realized states, \(B_t\) relevant boundaries,
\(\Adm_t\) admissible transformations, and \(\Rel_t\) the relations through
which states acquire practical or semantic significance.

Then

\[
  W_t\eqworld W_{t+n}
\]

is no longer trivial.

The relation must be earned.

A Macrolife civilization may preserve a genealogy while changing the
boundaries, affordances, material conditions, and admissibility relations
constituting its world.

Likewise, Shevek may return geographically while the relational world of
Anarres has changed.

And Card's Originist may reconstruct not merely the ancestry of a population
but a vanished field of relations in which the surviving traces once made
sense.

The world itself has become an Originist object.

% -----------------------------------------------------------------------------
%  5.9  THE DUALITY OF RECOVERY AND PROPAGATION
% -----------------------------------------------------------------------------

\subsection{Recovery and Propagation}
\label{sec:recovery-propagation}

The relation among the three literary probes can now be stated in its strongest
form.

Originism treats worldhood retrospectively:

\[
  \text{traces}
  \rightsquigarrow
  \text{candidate worlds}.
\]

Macrolife treats worldhood prospectively:

\[
  \text{world}
  \rightsquigarrow
  \text{candidate descendant worlds}.
\]

\emph{The Dispossessed} inserts recurrence:

\[
  W_0
  \rightsquigarrow
  W_1
  \rightsquigarrow
  W_0',
\]

where

\[
  W_0'\neq W_0
\]

need not prevent an appropriate recurrence relation.

Card therefore asks how a world can be recovered from what remains.

Zebrowski asks how a world can propagate through what changes.

Le Guin asks how a world can be left and returned to without pretending that
departure had no consequences.

These are not three definitions of worldhood.

They are three operations that any adequate account of world-continuation
must survive.

% -----------------------------------------------------------------------------
%  5.10  TRAJECTORY CLASSES
% -----------------------------------------------------------------------------

\subsection{Perhaps the Object Is a Class of Histories}
\label{sec:trajectory-classes}

The essay began by replacing state comparison with trajectory comparison.

We can now ask whether individual trajectories are themselves too specific.

Let

\[
  \gamma_1,\gamma_2\in\Gamma(X).
\]

Suppose they differ in details irrelevant to a selected continuation criterion
while preserving the same admissibility structure.

Define an equivalence relation

\[
  \gamma_1\sim_{\Adm}\gamma_2
\]

when the two trajectories are equivalent under the selected admissibility
criterion.

The relevant object becomes

\[
  [\gamma]_{\Adm}.
\]

This is the trajectory-class proposal:

\[
  [O\leadsto X]_{\Adm}.
\]

It has an immediate attraction.

If state identity is too strict and unrestricted causal ancestry too weak,
perhaps continuity belongs neither to the endpoint nor to the exact path but
to an admissibility-quotiented class of histories.

Yet the Originist objection applies immediately.

What does the quotient forget?

If

\[
  q_{\Adm}:
  \Gamma(X)
  \rightarrow
  \Gamma(X)/{\sim_{\Adm}},
\]

then we must ask

\[
  \lossorig(q_{\Adm};\delta)
\]

for every provenance distinction \(\delta\) relevant to subsequent inference.

The trajectory-class proposal therefore cannot simply be declared the final
ontology.

Its own quotient carries a wound.

% -----------------------------------------------------------------------------
%  5.11  OR PERHAPS THE GRAPH
% -----------------------------------------------------------------------------

\subsection{Perhaps the Object Is the Graph}
\label{sec:perhaps-graph}

The provenance graph avoids one weakness of the trajectory class by retaining
branching and convergence.

Instead of

\[
  [\gamma]_{\Adm},
\]

we might privilege

\[
  \graphG=(V,E,\tau,\Adm).
\]

This object can represent multiple origins, competing reconstructions,
independent rediscoveries, copied descendants, convergence, and recurrence.

It is therefore more expressive than a single trajectory.

But expressiveness is not fundamentality.

The graph itself depends upon decisions about what counts as a vertex, which
transformations deserve an edge, how finely transformations are typed, which
uncertain relations are represented, and where the graph is truncated.

A graph of provenance has provenance.

The representation is not outside the problem it represents.

% -----------------------------------------------------------------------------
%  5.12  THE PRIMITIVE-OBJECT QUESTION
% -----------------------------------------------------------------------------

\subsection{No Final Primitive Is Required}
\label{sec:no-final-primitive}

Three candidates now remain visible:

\[
  \gamma,
  \qquad
  [\gamma]_{\Adm},
  \qquad
  \graphG.
\]

The first preserves a realized path.

The second suppresses distinctions declared irrelevant under an admissibility
relation.

The third preserves branching structure at the cost of requiring explicit
choices about granularity and representation.

Nothing in the preceding argument requires one of these to be declared
metaphysically primitive.

Indeed, doing so would reproduce the error against which the essay has argued:
a representation would be selected, its losses would be forgotten, and the
result would then be mistaken for the object itself.

The more defensible conclusion is conditional.

\begin{held}[Primitive-Object Refusal]
\centering
Choose the weakest historical object that preserves every distinction required
by the inference being performed.\\[3pt]
Record what the choice forgets.
\end{held}

This is less satisfying than a final ontology.

It is also more consistent with the Originist objection.

% -----------------------------------------------------------------------------
%  5.13  NOTHING ORIGINAL NEED SURVIVE
% -----------------------------------------------------------------------------

\subsection{Nothing Original Need Survive}
\label{sec:nothing-original}

We can now return to the title without metaphor.

Let an origin \(O\) undergo a long admissible history

\[
  O
  \to
  x_1
  \to
  x_2
  \to
  \cdots
  \to
  x_n.
\]

Suppose no material component of \(O\) remains in \(x_n\):

\[
  M(O)\cap M(x_n)=\varnothing.
\]

Suppose further that no byte-exact representation of \(O\) survives, and that
some details of its history are irrecoverable.

It may nevertheless be the case that

\[
  [O\leadsto x_n]_{\Adm}
\]

constitutes a defensible continuation under the relevant constraints.

If so, nothing original in the narrow material or state-identical sense need
survive.

What survives is not necessarily a thing.

It may be a constrained descent relation.

This is why the title should not be read nihilistically.

The claim is not

\[
  \text{nothing need survive}.
\]

The claim is

\[
  \text{nothing \emph{original} need survive}.
\]

Something must constrain the descent if the word \enquote{continuation} is to
do more work than \enquote{caused by}.

But the constraint need not be an unchanged fragment carried intact from the
beginning.

% -----------------------------------------------------------------------------
%  5.14  CARD, ZEBROWSKI, LE GUIN
% -----------------------------------------------------------------------------

\subsection{Three Ways of Losing the Original}
\label{sec:three-losses}

Card removes the origin by temporal distance.

The present contains traces, records, reconstructions, and hypotheses, but the
origin itself is not available for inspection. Originism therefore asks what
can be recovered without pretending that reconstruction is direct possession
of the past \cite{card1989}.

Zebrowski removes the origin by transformation.

Civilization persists by becoming capable of inhabiting forms and environments
radically unlike those from which it began. Continuation therefore cannot be
identified with indefinite conservation of the initial substrate
\cite{zebrowski1979}.

Le Guin removes the origin by return.

Shevek comes back, but the existence of the journey prevents return from being
restoration. The origin encountered after recurrence is historically later
than the origin that was left \cite{leguin1974}.

Thus the original disappears in three different ways:

\[
  \text{distance},
  \qquad
  \text{transformation},
  \qquad
  \text{recurrence}.
\]

The same formal demand survives all three.

Do not infer identity of history from adequacy of state.

% -----------------------------------------------------------------------------
%  5.15  THE ESSAY AS AN INSTANCE
% -----------------------------------------------------------------------------

\subsection{This Document Has a Trajectory}
\label{sec:document-trajectory}

The provenance apparatus of this essay has deliberately made the document
slightly harder to read.

That cost requires justification.

The apparatus is not intended to simulate scientific precision by attaching
numbers to uncertain claims. Nor is it intended as a typographic joke about
the subject matter.

Its function is to make one distinction locally visible.

The document presented to the reader is a state.

Its sources, drafts, paraphrases, reconstructions, corrections, and
compilations constitute a history.

Those objects are not identical.

The same final PDF could, in principle, have been produced by many different
histories. One could be typed directly from inspected sources. Another could
be reconstructed from damaged notes. Another could be repeatedly summarized
and regenerated before returning to exactly the same final wording.

If the PDFs are byte-identical, no inspection of their bytes distinguishes
those histories.

This manuscript therefore refuses to treat its final state as a complete
description of itself.

The provenance tags are incomplete.

That incompleteness is not a defect hidden by the apparatus.

It is one of the things the apparatus records.

% -----------------------------------------------------------------------------
%  5.16  CLOSURE
% -----------------------------------------------------------------------------

\subsection{Closure Without Exhaustion}
\label{sec:closure}

The argument can stop here without claiming that the topic has been exhausted.

Closure is supplied by a different criterion.

The three literary probes initially forced a growing sequence of distinctions.
Each enlargement answered a failure produced by the preceding representation.

State required trajectory.

Trajectory required provenance structure.

Provenance required admissibility.

Admissibility raised the unresolved problem of world-continuation.

The pathology catalog then tested those distinctions against forgery,
rediscovery, compression, reconstruction, replacement, branching, convergence,
restoration, civilizational migration, and exact endpoint agreement.

The final movement has not produced another forced extension of the stack.

Instead, the existing apparatus has been sufficient to locate the remaining
uncertainty in the choice of primitive itself.

That is a legitimate stopping condition.

The question remains open whether the most useful fundamental object for a
given theory is

\[
  \gamma,
\]

\[
  [\gamma]_{\Adm},
\]

or

\[
  \graphG.
\]

The essay need not resolve this choice because its central claim survives all
three representations.

In each case, history contains distinctions not guaranteed to be recoverable
from the endpoint.

In each case, quotienting can erase distinctions required by later inference.

In each case, continuation must be evaluated over transformations rather than
assumed from resemblance.

And in each case, restoration of state does not undo descent.

\begin{held}[Final Proposition]
\centering
Continuity is not, in general, a property of states considered in isolation.\\[4pt]
It is a property attributed under constraints to histories of transformation.
\end{held}

Card's Originist looks backward along those histories.

Zebrowski's Macrolife extends them forward.

Le Guin bends them toward recurrence.

None supplies an unchanged object capable of carrying identity by itself.

None needs to.

The history is not an inconvenience surrounding the state.

It is part of what the state cannot say.

\bigskip

\begin{center}
  \Large\bfseries
  \(\State \neq \History\)
\end{center}

% ============================================================================
%  Epilogue
%  To be appended after Movement V (Descent). Uses only apparatus already
%  defined in originist-preamble.tex — no new macros introduced here, on the
%  grounds that an epilogue which needs new machinery has not actually closed
%  anything.
% ============================================================================

\section*{Epilogue: What This Document Cannot Tell You About Itself}
\movementquestion{Does the state of a text determine what kind of text it is?}

\epigraph{\uncertain{The proposition did not change. Its history did.}}{---provisional, pending the reader's own reconstruction}

This document has, at every point, been an instance of its own subject rather than an illustration standing outside it. \recon{The claims made about Leyel's archive, about the Anarres wall, about the substrate migrations of a macrolife civilization} have themselves passed through paraphrase, compression, and re-expression before arriving here; none of them carry a \provn{0} tag, and none should. What the apparatus above has tried to hold open is not the possibility of avoiding reconstruction but the requirement of \emph{knowing where it happened}.

\begin{held}[Closing claim, held rather than asserted outright]
\[
  \traj{x} \;\neq\; x, \qquad\text{and the difference is not decorative.}
\]
Two objects presently indistinguishable at the level of \(\eqbyte\) may remain permanently distinguishable at the level of \(\eqprov\). The essay has not tried to collapse this gap. It has tried to make the gap load-bearing.
\end{held}

The three probes converge here without merging. Card's question, \(?\leadsto P\), remains open in the ordinary sense that most origins are not recoverable and the essay has not claimed otherwise. Zebrowski's question, \(P\leadsto\,?\), remains open in the stronger sense that admissible continuation is a standing constraint problem, not a fact settled once for a given substrate. Le Guin's question, \(O\leadsto X\leadsto O'\), is the one this document has quietly been performing rather than merely describing: whatever returns at the end of an epilogue is \simconstraint\ what began the essay, not \eqbyte\ it.

\begin{trajectorydiagram}
  O \arrow[r] & X \arrow[r] & O' \arrow[loop right, distance=2em, "{O' \simconstraint O}"]
\end{trajectorydiagram}

\begin{wounded}
\begin{woundbox}
Consider, once more, the proposition attached earlier to \provomegadisplay. Nothing asserted by it has been retracted in the course of this document. What changed is not its content but the set of operations a careful reader is now entitled to perform with it — citation, extension, quotation, reliance. \(\lossorig\) was defined so that this sentence could be written truthfully rather than merely dramatically: some quotient taken early in the essay's own drafting made a distinction unrecoverable, and this paragraph is the residue of that loss, not a demonstration staged for effect.
\end{woundbox}
\end{wounded}

\bigskip
\noindent
\verbzero{Continuity is a property of trajectories, not of states.} That sentence has now appeared \recon{in at least one earlier form,} and this final instance of it is not guaranteed to be \eqbyte\ that earlier one. Whether it is \simI\ it — whether the invariant the essay set out to preserve has in fact survived the compression of writing it down — is, appropriately, not a question this document is in a position to answer about itself. \uncertain{That determination, like every determination of this kind, belongs to whoever reads it next.}

\bigskip
\noindent\hfill\footnotesize\itshape reconstructed spans in this epilogue: \thereconspanledger\ (running total, whole document)\hfill\null

\newpage 
\section*{Appendices}

% ============================================================================
%  APPENDICES
% ============================================================================

\clearpage
\appendix

\section{Path Spaces, Endpoint Maps, and Provenance Graphs}
\label{app:path-spaces}

\movementquestion{What mathematical structure is minimally required to distinguish a state from the histories that terminate in it?}

\begin{orientation}
The body of the essay deliberately used only as much mathematics as each
argument required. The present appendix collects the more technical
consequences of treating histories rather than states as first-class objects.

The central construction is elementary. A state space tells us what may be
present. A path space tells us how one state may become another. An endpoint
map forgets that history and retains only where it finished. Once written this
way, the Originist problem becomes a problem about the information lost by
that forgetting operation.

This appendix also makes precise the transition from a single trajectory to a
provenance graph. That transition matters whenever histories branch, converge,
or admit more than one derivational route between the same states. In those
cases there is no unique generation depth unless a path or aggregation rule
has first been selected.

Nothing here requires the state space to be continuous. The same definitions
apply to manuscripts, model outputs, software artifacts, archival records, or
civilizational states. When a continuous structure is useful, it will be
introduced explicitly rather than presumed.
\end{orientation}

% -----------------------------------------------------------------------------
%  A.1  STATE AND HISTORY SPACES
% -----------------------------------------------------------------------------

\subsection{States and Histories}
\label{app:states-histories}

Let \(X\) be a nonempty set of states. Let \(\mathcal T\) be a collection of
admissible transformation labels. A finite history over \(X\) is an expression

\[
  \gamma
  =
  \left(
  x_0
  \xrightarrow{T_1}
  x_1
  \xrightarrow{T_2}
  \cdots
  \xrightarrow{T_n}
  x_n
  \right),
\]

where

\[
  x_i\in X
  \qquad\text{and}\qquad
  T_i\in\mathcal T.
\]

Denote the collection of such histories by

\[
  \Gamma_{\mathcal T}(X).
\]

When the transformation labels are irrelevant or understood, write simply

\[
  \Gamma(X).
\]

Define the origin and endpoint maps

\[
  \Orig:\Gamma(X)\to X,
  \qquad
  \End:\Gamma(X)\to X
\]

by

\[
  \Orig(\gamma)=x_0,
  \qquad
  \End(\gamma)=x_n.
\]

The map \(\End\) is the simplest formal model of state-oriented observation.

It forgets everything except the terminal state.

\begin{definition}[Endpoint fibre]
\label{def:endpoint-fibre}
For \(x\in X\), the \emph{endpoint fibre} over \(x\) is

\[
  \End^{-1}(x)
  =
  \{
    \gamma\in\Gamma(X)
    :
    \End(\gamma)=x
  \}.
\]
\end{definition}

An Originist question is therefore frequently a question about the internal
structure of an endpoint fibre.

If

\[
  |\End^{-1}(x)|>1,
\]

then the endpoint does not uniquely determine its history.

\begin{proposition}[Endpoint non-identifiability]
\label{prop:endpoint-nonidentifiability}
If there exist distinct histories

\[
  \gamma_1\neq\gamma_2
\]

such that

\[
  \End(\gamma_1)=\End(\gamma_2),
\]

then no function

\[
  f:X\to\Gamma(X)
\]

can recover the realized history from the endpoint alone for every history in
\(\Gamma(X)\).
\end{proposition}

\begin{proof}
Assume such a function \(f\) exists. Let

\[
  \End(\gamma_1)
  =
  \End(\gamma_2)
  =
  x.
\]

Correct recovery would require simultaneously

\[
  f(x)=\gamma_1
\]

and

\[
  f(x)=\gamma_2.
\]

Since \(f\) is a function and \(\gamma_1\neq\gamma_2\), this is impossible.
\end{proof}

The proposition formalizes the most elementary Originist obstruction.

The problem is not insufficient computational power.

The information is absent from the endpoint representation.

% -----------------------------------------------------------------------------
%  A.2  OBSERVATIONAL QUOTIENTS
% -----------------------------------------------------------------------------

\subsection{Observation as Quotienting}
\label{app:observational-quotients}

In practice, an observer rarely has access even to the complete state \(x\).
Instead, a measurement or representation map

\[
  m:X\to Y
\]

produces an observable state \(m(x)\).

This induces an equivalence relation

\[
  x\sim_m y
  \quad\Longleftrightarrow\quad
  m(x)=m(y).
\]

The observer therefore works not directly with \(X\) but with equivalence
classes

\[
  [x]_m
  =
  \{
    y\in X:m(y)=m(x)
  \}.
\]

Composing observation with the endpoint map gives

\[
  \Gamma(X)
  \xrightarrow{\End}
  X
  \xrightarrow{m}
  Y.
\]

Thus two distinct losses may occur.

First,

\[
  \End
\]

forgets history.

Second,

\[
  m
\]

may forget distinctions among states.

The composite

\[
  m\circ\End
\]

therefore identifies histories whenever their endpoints are observationally
equivalent.

\begin{proposition}[Nested underdetermination]
\label{prop:nested-underdetermination}
For histories \(\gamma_1,\gamma_2\in\Gamma(X)\),

\[
  \End(\gamma_1)=\End(\gamma_2)
\]

implies

\[
  m(\End(\gamma_1))
  =
  m(\End(\gamma_2)),
\]

but the converse need not hold.
\end{proposition}

\begin{proof}
The forward implication follows immediately from applying \(m\) to equal
arguments. For the converse, choose distinct states \(x_1\neq x_2\) satisfying

\[
  m(x_1)=m(x_2),
\]

and histories terminating at \(x_1\) and \(x_2\), respectively. Their
observations agree although their endpoints need not.
\end{proof}

The distinction matters whenever apparent identity is defined relative to a
measurement regime.

Byte equality, semantic similarity, invariant equivalence, and visual
resemblance correspond to different choices of \(m\).

None should be mistaken for history itself.

% -----------------------------------------------------------------------------
%  A.3  TRAJECTORY EQUIVALENCE
% -----------------------------------------------------------------------------

\subsection{Equivalence Relations on Histories}
\label{app:trajectory-equivalence}

Exact equality of histories is often too strong.

Suppose two histories differ only in transformations irrelevant to a specified
task. We may introduce an equivalence relation

\[
  \sim_{\mathcal A}
\]

on \(\Gamma(X)\) and form the quotient

\[
  \Gamma(X)/{\sim_{\mathcal A}}.
\]

For \(\gamma\in\Gamma(X)\), write

\[
  [\gamma]_{\mathcal A}
\]

for its equivalence class.

The central warning of the essay can now be stated as a condition on such
quotients.

Let

\[
  \delta:
  \Gamma(X)\to D
\]

be a provenance-relevant observable.

If there exist

\[
  \gamma_1\sim_{\mathcal A}\gamma_2
\]

while

\[
  \delta(\gamma_1)\neq\delta(\gamma_2),
\]

then \(\delta\) does not descend to a well-defined function on the quotient.

\begin{theorem}[Descent criterion for provenance observables]
\label{thm:descent-criterion}
Let \(q:\Gamma(X)\to\Gamma(X)/{\sim}\) be the canonical quotient map. A
function

\[
  \delta:\Gamma(X)\to D
\]

factors through the quotient, so that there exists

\[
  \bar\delta:
  \Gamma(X)/{\sim}\to D
\]

satisfying

\[
  \delta=\bar\delta\circ q,
\]

if and only if \(\delta\) is constant on every equivalence class of
\(\sim\).
\end{theorem}

\begin{proof}
Suppose first that

\[
  \delta=\bar\delta\circ q.
\]

If

\[
  \gamma_1\sim\gamma_2,
\]

then

\[
  q(\gamma_1)=q(\gamma_2).
\]

Hence

\[
  \delta(\gamma_1)
  =
  \bar\delta(q(\gamma_1))
  =
  \bar\delta(q(\gamma_2))
  =
  \delta(\gamma_2).
\]

Thus \(\delta\) is constant on equivalence classes.

Conversely, suppose \(\delta\) is constant on every equivalence class. Define

\[
  \bar\delta([\gamma])
  =
  \delta(\gamma).
\]

Constancy on equivalence classes guarantees that this definition does not
depend upon the representative \(\gamma\). Therefore \(\bar\delta\) is
well-defined and

\[
  \delta=\bar\delta\circ q.
\]
\end{proof}

This familiar quotient criterion gives the Originist objection a precise
mathematical form.

A provenance quantity survives quotienting exactly when it is constant on the
histories that the quotient identifies.

If it is not constant, the quotient destroys the ability to recover it.

% -----------------------------------------------------------------------------
%  A.4  PROVENANCE GRAPHS
% -----------------------------------------------------------------------------

\subsection{From Paths to Directed Graphs}
\label{app:directed-graphs}

A single history assumes a linear ordering of derivation.

Branching requires a directed graph

\[
  \graphG=(V,E,\tau),
\]

where

\[
  \tau:E\to\mathcal T
\]

labels edges by transformation type.

For vertices \(u,v\in V\), let

\[
  \Pi_{\graphG}(u,v)
\]

denote the collection of directed paths from \(u\) to \(v\).

A provenance graph may therefore encode

\[
  |\Pi_{\graphG}(u,v)|=0,
\]

\[
  |\Pi_{\graphG}(u,v)|=1,
\]

or

\[
  |\Pi_{\graphG}(u,v)|>1.
\]

These correspond respectively to no represented derivation, a unique
represented path, or multiple represented derivational paths.

The final case is the one for which scalar generation depth becomes
ambiguous.

% -----------------------------------------------------------------------------
%  A.5  GENERATION DEPTH
% -----------------------------------------------------------------------------

\subsection{Generation Depth Is a Path Functional}
\label{app:generation-depth}

For a directed path

\[
  \pi
  =
  (
  v_0\to v_1\to\cdots\to v_n
  ),
\]

define its unweighted generation depth by

\[
  d(\pi)=n.
\]

If \(\Pi_{\graphG}(O,x)\) contains more than one path, then there is no unique
value

\[
  d(O,x)
\]

without an additional convention.

One may define

\[
  d_{\min}(O,x)
  =
  \min_{\pi\in\Pi_{\graphG}(O,x)}d(\pi),
\]

or, when finite,

\[
  d_{\max}(O,x)
  =
  \max_{\pi\in\Pi_{\graphG}(O,x)}d(\pi).
\]

One may instead privilege a historically realized path

\[
  \pi_{\mathrm{real}}
\]

and define

\[
  d_{\mathrm{real}}(O,x)
  =
  d(\pi_{\mathrm{real}}).
\]

These quantities answer different questions.

Shortest-path depth measures minimal represented derivational separation.

Realized-path depth measures the length of a selected history.

Neither encodes transformation type.

For that purpose, let

\[
  \tau(\pi)
  =
  \bigl(
    \tau(e_1),
    \tau(e_2),
    \ldots,
    \tau(e_n)
  \bigr)
\]

denote the ordered transformation word associated with the path.

Thus two paths may satisfy

\[
  d(\pi_1)=d(\pi_2)
\]

while

\[
  \tau(\pi_1)\neq\tau(\pi_2).
\]

This is the formal reason the notation

\[
  [3]
\]

may need to be refined to

\[
  [3:R\circ C\circ R].
\]

% -----------------------------------------------------------------------------
%  A.6  WEIGHTED PROVENANCE
% -----------------------------------------------------------------------------

\subsection{Weighted Derivational Distance}
\label{app:weighted-provenance}

Equal treatment of every transformation is rarely justified.

Let

\[
  w:\mathcal T\to\mathbb R_{\geq 0}
\]

assign a nonnegative provenance cost to each transformation type.

For a path

\[
  \pi=(e_1,\ldots,e_n),
\]

define

\[
  D_w(\pi)
  =
  \sum_{i=1}^{n}w(\tau(e_i)).
\]

A verbatim transfer might receive a small cost, a lossy compression a larger
cost, and an unconstrained reconstruction a still larger cost.

The mathematical construction does not determine those values.

They are task-relative.

More generally, provenance cost need not be additive. Let

\[
  D:\Pi(\graphG)\to\mathbb R_{\geq 0}
\]

be any path functional capable of depending upon ordering and interaction
among transformations.

Then

\[
  D(R\circ C)
\]

need not equal

\[
  D(C\circ R).
\]

This is important whenever reconstruction before compression and compression
before reconstruction preserve different information.

\begin{remark}
The essay's generation notation should therefore not be interpreted as a
universal metric. It is a deliberately crude coordinate that becomes useful
only after the relevant transformation vocabulary and inferential task have
been specified.
\end{remark}

% -----------------------------------------------------------------------------
%  A.7  CONVERGENCE AND NON-RECOVERABILITY
% -----------------------------------------------------------------------------

\subsection{Convergence}
\label{app:convergence}

Let distinct paths

\[
  \pi_1,\pi_2\in\Pi_{\graphG}(O,x)
\]

terminate at the same vertex \(x\).

If only \(x\) is retained, the map

\[
  \pi\mapsto\End(\pi)
\]

identifies \(\pi_1\) and \(\pi_2\).

The endpoint is therefore a quotient of the path information.

This observation generalizes immediately to any representation

\[
  r:\Pi(\graphG)\to Z.
\]

Whenever

\[
  r(\pi_1)=r(\pi_2)
\]

for provenance-distinct paths, the representation loses the distinction
between them.

\begin{proposition}[Convergent provenance loss]
\label{prop:convergent-provenance-loss}
Let \(\delta\) be a provenance observable with

\[
  \delta(\pi_1)\neq\delta(\pi_2).
\]

If a representation \(r\) satisfies

\[
  r(\pi_1)=r(\pi_2),
\]

then no function of \(r(\pi)\) alone can recover \(\delta\) correctly on both
paths.
\end{proposition}

\begin{proof}
If such a function \(h\) existed, then

\[
  h(r(\pi_1))
  =
  \delta(\pi_1)
\]

and

\[
  h(r(\pi_2))
  =
  \delta(\pi_2).
\]

But

\[
  r(\pi_1)=r(\pi_2),
\]

so the two left-hand sides are equal while the two right-hand sides are not,
a contradiction.
\end{proof}

This is the information-theoretic core of the Originist position in its
simplest deterministic form.

% -----------------------------------------------------------------------------
%  A.8  BRANCHING AND IDENTITY
% -----------------------------------------------------------------------------

\subsection{Branching Prevents Identity from Selecting a Unique Future}
\label{app:branching-identity}

Suppose an origin \(O\) has two admissible descendants

\[
  O\leadsto A
\]

and

\[
  O\leadsto B.
\]

If descent from \(O\) were sufficient to establish numerical identity with
\(O\), then both \(A\) and \(B\) would inherit that identity.

By transitivity of identity,

\[
  A=O
  \qquad\text{and}\qquad
  B=O
\]

would imply

\[
  A=B.
\]

But branching permits

\[
  A\neq B.
\]

Thus genealogical descent cannot by itself be numerical identity.

\begin{proposition}[Branching obstruction]
\label{prop:branching-obstruction}
No transitive numerical identity relation can identify an origin with every
member of a genuinely branching set of distinct descendants.
\end{proposition}

\begin{proof}
Let \(A\neq B\) be descendants of \(O\). Assume

\[
  A=O
  \qquad\text{and}\qquad
  B=O.
\]

By symmetry,

\[
  O=B,
\]

and by transitivity,

\[
  A=B,
\]

contradicting the assumption that the descendants are distinct.
\end{proof}

The result is elementary but useful for consciousness-transfer and copying
problems.

A continuation relation may branch.

Numerical identity cannot branch in the same way without contradiction.

% -----------------------------------------------------------------------------
%  A.9  RECURRENCE
% -----------------------------------------------------------------------------

\subsection{Recurrence as Return to a Class}
\label{app:recurrence-class}

Let

\[
  q_{\mathcal C}:X\to X/{\sim_{\mathcal C}}
\]

be the quotient induced by a constraint-relative recurrence relation.

A history

\[
  \gamma
  =
  (
  x_0\to\cdots\to x_n
  )
\]

is recurrent relative to \(\mathcal C\) when

\[
  q_{\mathcal C}(x_n)
  =
  q_{\mathcal C}(x_0).
\]

Equivalently,

\[
  x_n\sim_{\mathcal C}x_0.
\]

This does not require

\[
  x_n=x_0.
\]

Nor does it imply equality of histories.

Indeed, if \(n>0\),

\[
  \gamma
\]

contains an ordered sequence of transformations absent from the zero-length
history at \(x_0\).

Recurrence therefore closes a loop only after quotienting by
\(\sim_{\mathcal C}\).

At the history level, the path remains extended.

This supplies a concise formal interpretation of the Le Guinian return:

\[
  q_{\mathcal C}(O')
  =
  q_{\mathcal C}(O),
\]

while

\[
  O'\neq O
\]

and

\[
  \Hist(O')\neq\Hist(O).
\]

% -----------------------------------------------------------------------------
%  A.10  CONTINUOUS PATH SPACES
% -----------------------------------------------------------------------------

\subsection{Continuous Generalization}
\label{app:continuous-paths}

When \(X\) carries a topology, a continuous trajectory may be represented as

\[
  \gamma:[0,1]\to X.
\]

For fixed endpoints \(x,y\in X\), define

\[
  \Gamma(x,y)
  =
  \{
    \gamma:[0,1]\to X
    :
    \gamma(0)=x,\;
    \gamma(1)=y
  \}.
\]

Two trajectories may share endpoints while differing as maps:

\[
  \gamma_1(0)=\gamma_2(0)=x,
\]

\[
  \gamma_1(1)=\gamma_2(1)=y,
\]

yet

\[
  \gamma_1\neq\gamma_2.
\]

If a homotopy relation is imposed, one may further quotient

\[
  \Gamma(x,y)/{\simeq}.
\]

The Originist question immediately reappears:

\[
  \text{Which distinctions are erased by the homotopy quotient?}
\]

For some tasks, homotopic paths should indeed be identified.

For others, the difference between them may encode the very history under
investigation.

Thus topology does not solve the provenance problem.

It provides another family of principled quotients to which the same question
must be applied.

% -----------------------------------------------------------------------------
%  A.11  THE FORMAL CORE
% -----------------------------------------------------------------------------

\subsection{The Formal Core of the Trajectory Thesis}
\label{app:formal-core}

The mathematical content of the trajectory-first claim can now be stated
without literary language.

Let

\[
  \End:\Gamma(X)\to X
\]

be non-injective.

Then there exist histories

\[
  \gamma_A\neq\gamma_B
\]

with

\[
  \End(\gamma_A)=\End(\gamma_B).
\]

Any property

\[
  P:\Gamma(X)\to Z
\]

that is not constant on endpoint fibres cannot be represented exactly as a
function of endpoints alone.

\begin{theorem}[Trajectory necessity criterion]
\label{thm:trajectory-necessity}
A history-dependent property

\[
  P:\Gamma(X)\to Z
\]

admits an exact state-level representation

\[
  \widetilde P:X\to Z
\]

satisfying

\[
  P=\widetilde P\circ\End
\]

if and only if \(P\) is constant on every endpoint fibre.
\end{theorem}

\begin{proof}
This is \cref{thm:descent-criterion} applied to the equivalence relation

\[
  \gamma_1\sim_{\End}\gamma_2
  \quad\Longleftrightarrow\quad
  \End(\gamma_1)=\End(\gamma_2).
\]

The quotient by \(\sim_{\End}\) is represented by endpoint states. Therefore
\(P\) factors through \(\End\) exactly when it is constant on the fibres of
\(\End\).
\end{proof}

This theorem states the central issue in its most economical form.

If the property of interest varies among histories that end in the same state,
then the property is not a property of the endpoint alone.

Provenance is one such property.

Under the argument developed in this essay, some forms of continuation are
another.

The mathematical statement does not decide which historical properties matter.

That is the work of admissibility.

% ============================================================================
%  APPENDIX B
%  QUOTIENTS AND THE ORIGINIST LOSS OPERATOR
% ============================================================================

\section{Quotients and the Originist Loss Operator}
\label{app:originist-loss}

\movementquestion{How can one say precisely what a representation forgets without pretending that every kind of forgetting has the same magnitude?}

\begin{orientation}
The phrase \enquote{information loss} is dangerously broad. A representation
can forget punctuation, ordering, authorship, transformation history,
confidence, material composition, or an entire branch of descent. These are
not interchangeable losses, and there is no reason to assume that they admit
a common numerical scale.

The Originist loss operator introduced in the body of the essay is therefore
not initially a number. Its most primitive form records whether a distinction
required by a later inference survives a representation at all. Only after
that logical question has been answered does it make sense to introduce
cardinalities, probabilities, entropies, or task-relative costs.

This appendix develops that hierarchy. It begins with a binary notion of
recoverability, extends it to families of provenance observables, and then
shows how probabilistic loss can be defined when a distribution over histories
is available. The final sections return to compression and reconstruction.

The important result is not that lossy compression is bad. Almost every useful
abstraction is lossy. The result is that reconstruction of an endpoint cannot,
by itself, make a provenance-destroying quotient provenance-neutral.
\end{orientation}

% -----------------------------------------------------------------------------
%  B.1  QUOTIENTS AS FORGETTING OPERATIONS
% -----------------------------------------------------------------------------

\subsection{Representations and Their Fibres}
\label{app:representation-fibres}

Let \(\Gamma\) be a space of histories and let

\[
  q:\Gamma\to Q
\]

be any representation map.

The codomain \(Q\) may contain compressed records, canonical forms, endpoint
states, equivalence classes, summaries, embeddings, or any other objects used
in place of the full histories.

For \(z\in Q\), define the fibre

\[
  q^{-1}(z)
  =
  \{
    \gamma\in\Gamma:q(\gamma)=z
  \}.
\]

The fibre contains every history rendered indistinguishable by the
representation \(q\).

A representation is injective precisely when every fibre contains at most one
history.

In most cases of interest here, \(q\) is deliberately non-injective.

The question is therefore not whether distinctions have been erased.

The question is which ones.

% -----------------------------------------------------------------------------
%  B.2  PROVENANCE OBSERVABLES
% -----------------------------------------------------------------------------

\subsection{Provenance Observables}
\label{app:provenance-observables}

Let

\[
  \delta:\Gamma\to D
\]

be a function encoding some provenance-relevant distinction.

Examples include the original source,

\[
  \delta_{\mathrm{src}}(\gamma),
\]

generation depth,

\[
  \delta_{\mathrm{depth}}(\gamma),
\]

the ordered transformation word,

\[
  \delta_{\mathrm{word}}(\gamma),
\]

the presence of a lossy operation,

\[
  \delta_{\mathrm{loss}}(\gamma)\in\{0,1\},
\]

or membership in an admissible trajectory class,

\[
  \delta_{\Adm}(\gamma).
\]

No assumption is made that all provenance questions can be reduced to one
observable.

Indeed, the argument of the essay suggests the opposite.

Let

\[
  \Delta
  =
  \{
    \delta_\alpha
  \}_{\alpha\in A}
\]

be a family of provenance observables selected for a task.

A representation may preserve some members of \(\Delta\) and destroy others.

% -----------------------------------------------------------------------------
%  B.3  LOGICAL ORIGINIST LOSS
% -----------------------------------------------------------------------------

\subsection{Logical Originist Loss}
\label{app:logical-originist-loss}

The weakest useful notion of loss asks only whether a provenance observable
survives the representation.

\begin{definition}[Logical Originist loss]
\label{def:logical-originist-loss}
For a representation

\[
  q:\Gamma\to Q
\]

and provenance observable

\[
  \delta:\Gamma\to D,
\]

define

\[
  \lossorig^{(0)}(q;\delta)
  =
  \begin{cases}
    0,
      & \text{if }\delta\text{ factors through }q,\\[4pt]
    1,
      & \text{otherwise.}
  \end{cases}
\]
\end{definition}

Thus

\[
  \lossorig^{(0)}(q;\delta)=0
\]

exactly when there exists a function

\[
  \bar\delta:Q\to D
\]

such that

\[
  \delta=\bar\delta\circ q.
\]

By \cref{thm:descent-criterion}, this occurs exactly when \(\delta\) is
constant on every fibre of \(q\).

Hence

\[
  \lossorig^{(0)}(q;\delta)=1
\]

if and only if there exist histories \(\gamma_1,\gamma_2\) satisfying

\[
  q(\gamma_1)=q(\gamma_2)
\]

while

\[
  \delta(\gamma_1)\neq\delta(\gamma_2).
\]

This is the formal version of an Originist wound.

% -----------------------------------------------------------------------------
%  B.4  TASK-RELATIVE SAFETY
% -----------------------------------------------------------------------------

\subsection{Loss Is Relative to a Future Question}
\label{app:task-relative-safety}

A representation can be perfectly safe for one task and unusable for another.

Let

\[
  F:\Gamma\to Y
\]

represent a future inferential task.

If \(F\) factors through \(q\), then the representation retains everything
required to compute \(F\) exactly:

\[
  F=\bar F\circ q.
\]

If it does not, then some histories requiring different outputs have become
indistinguishable after representation.

\begin{definition}[Task sufficiency]
\label{def:task-sufficiency}
A representation \(q:\Gamma\to Q\) is \emph{sufficient for the task \(F\)}
when there exists

\[
  \bar F:Q\to Y
\]

such that

\[
  F=\bar F\circ q.
\]
\end{definition}

The Originist objection can now be written in one line:

\[
  q\text{ is unsafe for }F
  \quad\Longleftrightarrow\quad
  F\text{ does not factor through }q.
\]

This is stronger than the vague instruction to preserve provenance.

Only provenance distinctions on which \(F\) depends need be retained.

The framework therefore does not imply archival maximalism.

% -----------------------------------------------------------------------------
%  B.5  A FAMILY OF DISTINCTIONS
% -----------------------------------------------------------------------------

\subsection{Loss Profiles}
\label{app:loss-profiles}

Given a family

\[
  \Delta
  =
  \{
    \delta_\alpha
  \}_{\alpha\in A},
\]

define the logical loss profile

\[
  \mathbf L_q
  =
  \left(
    \lossorig^{(0)}(q;\delta_\alpha)
  \right)_{\alpha\in A}.
\]

This vector is more informative than a single scalar whenever the provenance
dimensions are heterogeneous.

For example, a representation might preserve

\[
  \delta_{\mathrm{src}}
\]

and

\[
  \delta_{\mathrm{depth}},
\]

while destroying

\[
  \delta_{\mathrm{word}}.
\]

Its loss profile could then take the schematic form

\[
  \mathbf L_q=(0,0,1).
\]

No meaningful conclusion follows from replacing this vector with the number
\(1\) unless weights and a task have first been specified.

The refusal to scalarize prematurely is deliberate.

Different wounds matter to different future operations.

% -----------------------------------------------------------------------------
%  B.6  WEIGHTED ORIGINIST LOSS
% -----------------------------------------------------------------------------

\subsection{Task-Weighted Loss}
\label{app:weighted-loss}

Suppose \(A\) is finite and let

\[
  w_\alpha\geq 0
\]

represent the importance of preserving \(\delta_\alpha\) for a specified
task.

Define

\[
  \lossorig^{(w)}(q;\Delta)
  =
  \sum_{\alpha\in A}
  w_\alpha
  \lossorig^{(0)}(q;\delta_\alpha).
\]

This scalar has no task-independent interpretation.

Changing the weights changes what the representation is being evaluated for.

If exact wording is irrelevant, one may set

\[
  w_{\mathrm{word}}=0.
\]

If authorship is legally decisive, one may assign

\[
  w_{\mathrm{src}}
\]

a dominant value.

The weights therefore encode a normative or operational choice.

They are not discovered merely by inspecting \(q\).

% -----------------------------------------------------------------------------
%  B.7  MONOTONICITY UNDER FURTHER QUOTIENTING
% -----------------------------------------------------------------------------

\subsection{Further Forgetting Cannot Restore Logical Recoverability}
\label{app:monotonicity}

Suppose

\[
  \Gamma
  \xrightarrow{q_1}
  Q_1
  \xrightarrow{q_2}
  Q_2
\]

are successive representations, and define

\[
  q=q_2\circ q_1.
\]

If a provenance distinction is already unrecoverable after \(q_1\), applying
\(q_2\) cannot make it exactly recoverable from the representation alone.

\begin{theorem}[Monotonicity of logical loss]
\label{thm:logical-loss-monotonicity}
For every provenance observable \(\delta\),

\[
  \lossorig^{(0)}(q_1;\delta)
  \leq
  \lossorig^{(0)}(q_2\circ q_1;\delta).
\]
\end{theorem}

\begin{proof}
If

\[
  \lossorig^{(0)}(q_1;\delta)=0,
\]

the inequality is automatic because the left side is zero.

Suppose instead

\[
  \lossorig^{(0)}(q_1;\delta)=1.
\]

Then there exist \(\gamma_1,\gamma_2\in\Gamma\) such that

\[
  q_1(\gamma_1)=q_1(\gamma_2)
\]

while

\[
  \delta(\gamma_1)\neq\delta(\gamma_2).
\]

Applying \(q_2\) gives

\[
  q_2(q_1(\gamma_1))
  =
  q_2(q_1(\gamma_2)).
\]

Thus the composite also identifies two histories on which \(\delta\) differs.
Therefore

\[
  \lossorig^{(0)}(q_2\circ q_1;\delta)=1.
\]
\end{proof}

The theorem concerns information available from the representation alone.

External evidence can, of course, restore knowledge.

That restoration is the next problem.

% -----------------------------------------------------------------------------
%  B.8  RECONSTRUCTION WITH EXTERNAL INFORMATION
% -----------------------------------------------------------------------------

\subsection{Reconstruction Is Not Inversion}
\label{app:reconstruction-not-inversion}

Let

\[
  q:\Gamma\to Q
\]

be non-injective.

Suppose additional evidence belongs to a space \(E\), and let

\[
  R:Q\times E\to\widehat\Gamma
\]

be a reconstruction procedure.

For some realized history \(\gamma\), it may happen that

\[
  R(q(\gamma),e)=\gamma.
\]

This does not imply that \(q\) was lossless.

The evidence \(e\) is doing inferential work.

To make this explicit, define the augmented representation

\[
  q_e:\Gamma\to Q\times E
\]

only when the evidence itself is generated as part of the represented
information.

If \(e\) arrives independently, then reconstruction uses information outside
\(q(\gamma)\).

\begin{proposition}[External repair does not erase internal loss]
\label{prop:external-repair}
Suppose

\[
  \lossorig^{(0)}(q;\delta)=1.
\]

The existence of an external evidence variable \(e\) and reconstruction
procedure \(R\) that correctly recovers \(\delta(\gamma)\) for a particular
history does not imply

\[
  \lossorig^{(0)}(q;\delta)=0.
\]
\end{proposition}

\begin{proof}
Logical Originist loss is defined by whether \(\delta\) factors through \(q\).
The existence of a procedure depending upon both \(q(\gamma)\) and an
additional variable \(e\) establishes, at most, factorization through an
augmented representation. It does not establish factorization through \(q\)
alone.
\end{proof}

This is the exact distinction between preservation and successful
reconstruction.

% -----------------------------------------------------------------------------
%  B.9  BAYESIAN RECONSTRUCTION
% -----------------------------------------------------------------------------

\subsection{Probabilistic Reconstruction}
\label{app:bayesian-reconstruction}

When exact recovery is impossible, an Originist may maintain a distribution
over candidate histories.

Let \(\Gamma\) be finite or countable for simplicity, with prior

\[
  P(\gamma).
\]

After observing

\[
  z=q(\gamma),
\]

Bayes' rule gives

\[
  P(\gamma\mid z)
  =
  \frac{
    P(z\mid\gamma)P(\gamma)
  }{
    P(z)
  }.
\]

For deterministic \(q\),

\[
  P(z\mid\gamma)
  =
  \begin{cases}
    1,&q(\gamma)=z,\\
    0,&\text{otherwise.}
  \end{cases}
\]

Hence the posterior is supported on the fibre

\[
  q^{-1}(z).
\]

If that fibre contains multiple histories with nonzero posterior mass, the
representation leaves provenance uncertainty.

A reconstruction algorithm may select

\[
  \widehat\gamma_{\mathrm{MAP}}
  =
  \arg\max_{\gamma\in q^{-1}(z)}
  P(\gamma\mid z).
\]

But

\[
  \widehat\gamma_{\mathrm{MAP}}
\]

is an estimate.

It is not a recovered historical fact merely because it is probable.

This distinction is particularly important for generative reconstruction,
where plausible completions can be much easier to produce than derivationally
warranted ones.

% -----------------------------------------------------------------------------
%  B.10  CONDITIONAL ENTROPY
% -----------------------------------------------------------------------------

\subsection{An Information-Theoretic Originist Loss}
\label{app:conditional-entropy}

When a probability distribution is available, provenance uncertainty can be
quantified.

Let the random variable \(G\) range over histories and let

\[
  Z=q(G)
\]

be the represented state.

The conditional entropy

\[
  H(G\mid Z)
\]

measures residual uncertainty about history after observing the
representation.

If

\[
  H(G\mid Z)=0,
\]

then \(G\) is determined almost surely by \(Z\) under the chosen distribution.

If

\[
  H(G\mid Z)>0,
\]

some historical uncertainty remains.

For a particular provenance observable

\[
  D=\delta(G),
\]

the more task-specific quantity is

\[
  H(D\mid Z).
\]

This suggests the probabilistic Originist loss

\[
  \lossorig^{(H)}(q;\delta)
  =
  H(\delta(G)\mid q(G)).
\]

This quantity depends upon the distribution over histories.

It therefore differs fundamentally from the logical loss

\[
  \lossorig^{(0)}.
\]

Logical loss asks whether exact recovery is possible in principle over the
specified domain.

Entropy-based loss asks how much uncertainty remains on average under a
specified probability model.

% -----------------------------------------------------------------------------
%  B.11  LOGICAL LOSS WITHOUT LARGE ENTROPY
% -----------------------------------------------------------------------------

\subsection{Rare Ambiguity}
\label{app:rare-ambiguity}

The distinction between logical and probabilistic loss is not cosmetic.

Suppose two histories \(\gamma_1\) and \(\gamma_2\) satisfy

\[
  q(\gamma_1)=q(\gamma_2),
\]

and

\[
  \delta(\gamma_1)\neq\delta(\gamma_2).
\]

Then

\[
  \lossorig^{(0)}(q;\delta)=1.
\]

But if

\[
  P(\gamma_2)
\]

is extremely small, then

\[
  H(\delta(G)\mid q(G))
\]

may also be extremely small.

The representation is almost always sufficient probabilistically while still
failing exact recoverability.

This is precisely the sort of case in which a system optimized for average
performance may appear reliable until a low-probability provenance dispute
becomes the only case that matters.

Archival authentication, legal evidence, scientific priority, and safety
auditing can all have this structure.

% -----------------------------------------------------------------------------
%  B.12  DATA PROCESSING
% -----------------------------------------------------------------------------

\subsection{A Provenance Form of Data Processing}
\label{app:data-processing}

Suppose

\[
  G\to Z_1\to Z_2
\]

forms a Markov chain in which

\[
  Z_1=q_1(G)
\]

and

\[
  Z_2=q_2(Z_1).
\]

For any provenance variable

\[
  D=\delta(G),
\]

the data-processing inequality gives

\[
  I(D;Z_2)
  \leq
  I(D;Z_1),
\]

where \(I\) denotes mutual information.

Equivalently,

\[
  H(D\mid Z_2)
  \geq
  H(D\mid Z_1)
\]

when the relevant entropies are finite.

Further processing cannot increase information about provenance unless new
information enters the system.

This is the probabilistic counterpart of
\cref{thm:logical-loss-monotonicity}.

A reconstruction model can appear to increase information because it contains
priors learned elsewhere.

That information is genuinely new relative to the compressed artifact.

It should not be misdescribed as information that the artifact itself
preserved.

% -----------------------------------------------------------------------------
%  B.13  COMPRESSION AND GENERATIVE REPAIR
% -----------------------------------------------------------------------------

\subsection{Repeated Compression--Repair Cycles}
\label{app:compression-repair-cycles}

Consider a sequence

\[
  x_0
  \xrightarrow{C_1}
  z_1
  \xrightarrow{R_1}
  x_1
  \xrightarrow{C_2}
  z_2
  \xrightarrow{R_2}
  x_2
  \to\cdots.
\]

Let

\[
  \gamma_n
\]

denote the history terminating at \(x_n\).

Suppose the invariant loss function used by the repair system satisfies

\[
  \mathcal L_{\Inv}(x_0,x_n)\approx 0.
\]

This may be an excellent result.

It means the selected invariants have been successfully maintained.

It does not imply

\[
  \Prov(x_0)=\Prov(x_n).
\]

Nor does it imply

\[
  \lossorig(C_i;\delta)=0
\]

for every provenance observable \(\delta\).

The repair system and the Originist answer different questions.

The repair system asks whether a selected content property has survived or
been reconstructed.

The Originist asks what derivational claims remain warranted about the result.

% -----------------------------------------------------------------------------
%  B.14  GENERATION COUNT AS A LOWER-DIMENSIONAL SUMMARY
% -----------------------------------------------------------------------------

\subsection{Generation Count and Transformation Words}
\label{app:generation-summary}

Let

\[
  \tau(\gamma)
  =
  T_n\circ\cdots\circ T_1
\]

denote the ordered transformation word associated with a history.

Define generation count

\[
  g(\gamma)=n.
\]

Then \(g\) is a projection from transformation words to their lengths.

Consequently,

\[
  g(\gamma_1)=g(\gamma_2)
\]

does not imply

\[
  \tau(\gamma_1)=\tau(\gamma_2).
\]

Generation count is therefore itself a quotient.

Let

\[
  q_g:\Gamma\to\mathbb N
\]

be given by

\[
  q_g(\gamma)=g(\gamma).
\]

If a future task depends on transformation order, then

\[
  \lossorig^{(0)}(q_g;\delta_{\mathrm{word}})=1.
\]

This is why the provenance notation had to expand from

\[
  [3]
\]

to

\[
  [3:R\circ C\circ R].
\]

The richer notation is not necessarily final.

It merely preserves another distinction.

% -----------------------------------------------------------------------------
%  B.15  PROVENANCE CONFIDENCE
% -----------------------------------------------------------------------------

\subsection{Confidence Is Not Generation}
\label{app:confidence-generation}

The full provenance notation introduced in the essay permits expressions such
as

\[
  \varphi^{[3:R\circ C\circ R;\,\Inv^+;\,p=.71]}.
\]

The final coordinate must be interpreted carefully.

A confidence value

\[
  p=.71
\]

is not a generation count.

It is not a truth value.

It is not an invariant score.

It is not a measure of semantic similarity.

It represents an evidential assessment under some specified model.

Two claims may therefore satisfy

\[
  \Truth(\varphi_1)=\Truth(\varphi_2)
\]

and

\[
  g(\varphi_1)=g(\varphi_2)
\]

while carrying different confidence values because the evidence for their
respective derivational histories differs.

Conversely, two provenance hypotheses may have equal posterior probability
while describing entirely different transformation chains.

The coordinates must remain separate unless a model explicitly relates them.

% -----------------------------------------------------------------------------
%  B.16  THE UNKNOWN-DEPTH CASE
% -----------------------------------------------------------------------------

\subsection{Why \texorpdfstring{\([\omega?]\)}{[omega?]} Is Not Infinity}
\label{app:omega-question}

The notation

\[
  [\omega?]
\]

used for the wound in Movement~IV is intentionally suggestive but should not
be read as a literal transfinite generation count.

Its meaning is epistemic:

\[
  \text{derivational depth is not presently recoverable}.
\]

Thus

\[
  [\omega?]
  \neq
  [\infty].
\]

Nor does it mean that the derivation is known to be arbitrarily long.

It marks refusal to invent a finite depth when the available evidence does not
warrant one.

Formally, if

\[
  g:\Gamma\to\mathbb N
\]

is generation depth and the evidence \(E\) leaves multiple values possible,

\[
  |\{g(\gamma):\gamma\in\Gamma_E\}|>1,
\]

then assigning a unique \(n\in\mathbb N\) would overstate what is known.

The wound tag therefore records underdetermination rather than magnitude.

% -----------------------------------------------------------------------------
%  B.17  REPAIR CANNOT MAKE A QUOTIENT INJECTIVE
% -----------------------------------------------------------------------------

\subsection{The Non-Neutrality of Perfect Repair}
\label{app:repair-nonneutrality}

We can now state the strongest result needed for the compression argument.

Let

\[
  q:\Gamma\to Q
\]

be a provenance-losing representation and let

\[
  R:Q\to X
\]

be any deterministic repair operation producing endpoint states.

Suppose there exist histories

\[
  \gamma_1,\gamma_2
\]

such that

\[
  q(\gamma_1)=q(\gamma_2)
\]

but

\[
  \delta(\gamma_1)\neq\delta(\gamma_2).
\]

Then

\[
  R(q(\gamma_1))
  =
  R(q(\gamma_2)).
\]

No quality of endpoint repair changes the fact that the two histories were
identified before reconstruction.

\begin{theorem}[Perfect repair cannot neutralize provenance loss]
\label{thm:perfect-repair}
If

\[
  \lossorig^{(0)}(q;\delta)=1,
\]

then for every deterministic repair map

\[
  R:Q\to X,
\]

the composite

\[
  R\circ q
\]

also has positive logical Originist loss with respect to \(\delta\):

\[
  \lossorig^{(0)}(R\circ q;\delta)=1.
\]
\end{theorem}

\begin{proof}
Since

\[
  \lossorig^{(0)}(q;\delta)=1,
\]

there exist \(\gamma_1,\gamma_2\) such that

\[
  q(\gamma_1)=q(\gamma_2)
\]

and

\[
  \delta(\gamma_1)\neq\delta(\gamma_2).
\]

Applying \(R\) to the equal representations yields

\[
  R(q(\gamma_1))
  =
  R(q(\gamma_2)).
\]

Thus \(R\circ q\) also identifies histories distinguished by \(\delta\).
Therefore \(\delta\) does not factor through \(R\circ q\), and

\[
  \lossorig^{(0)}(R\circ q;\delta)=1.
\]
\end{proof}

The theorem does not say that repair is useless.

It says something narrower.

Repair can restore an endpoint property.

It cannot make an earlier non-injective provenance representation become
injective after the fact.

% -----------------------------------------------------------------------------
%  B.18  SIDE INFORMATION
% -----------------------------------------------------------------------------

\subsection{How Provenance Can Be Recovered}
\label{app:side-information}

The preceding theorem might appear stronger than ordinary archival practice
allows. Lost provenance is sometimes recovered.

The apparent contradiction disappears once side information is represented
explicitly.

Suppose

\[
  q(\gamma_1)=q(\gamma_2),
\]

but an independent variable \(E\) satisfies

\[
  E(\gamma_1)\neq E(\gamma_2).
\]

The augmented representation

\[
  q^+(\gamma)
  =
  \bigl(q(\gamma),E(\gamma)\bigr)
\]

may now distinguish the histories.

It is possible that

\[
  \lossorig^{(0)}(q^+;\delta)=0
\]

even though

\[
  \lossorig^{(0)}(q;\delta)=1.
\]

This is not reversal of loss inside \(q\).

It is recovery through another channel.

An archive, witness, checksum, external manuscript, independent sensor, or
second model may supply exactly such a channel.

The Archivist and Originist roles therefore become complementary.

The Archivist preserves side information before it is needed.

The Originist determines when later claims depend upon it.

% -----------------------------------------------------------------------------
%  B.19  MINIMAL SUFFICIENT PROVENANCE
% -----------------------------------------------------------------------------

\subsection{Minimal Sufficient Provenance}
\label{app:minimal-sufficient-provenance}

The theory should not terminate in the instruction to store everything.

Let

\[
  F:\Gamma\to Y
\]

be the future task of interest.

We seek a representation

\[
  q_F:\Gamma\to Q_F
\]

such that

\[
  F=\bar F\circ q_F
\]

while \(q_F\) discards distinctions irrelevant to \(F\).

This resembles the search for a sufficient statistic.

Two histories may be identified whenever they induce the same task-relevant
output:

\[
  \gamma_1\sim_F\gamma_2
  \quad\Longleftrightarrow\quad
  F(\gamma_1)=F(\gamma_2).
\]

The corresponding quotient

\[
  \Gamma/{\sim_F}
\]

is maximally compressed relative to the single task \(F\).

But future tasks are rarely known completely in advance.

This is where archival judgment enters.

The more uncertain the future task family, the more dangerous aggressive
quotienting becomes.

% -----------------------------------------------------------------------------
%  B.20  AN ORIGINIST RATE--DISTORTION ANALOGUE
% -----------------------------------------------------------------------------

\subsection{Compression Under Provenance Constraints}
\label{app:originist-rate-distortion}

The preceding observation suggests a rate--distortion-like formulation.

Let

\[
  R(q)
\]

denote the storage or description cost of representation \(q\), and let

\[
  D_{\mathrm{state}}(q)
\]

measure distortion of task-relevant endpoint properties.

Add an Originist term

\[
  D_{\mathrm{prov}}(q)
  =
  \lossorig^{(w)}(q;\Delta).
\]

One may then consider an objective of the form

\[
  \mathcal J(q)
  =
  R(q)
  +
  \lambda
  D_{\mathrm{state}}(q)
  +
  \mu
  D_{\mathrm{prov}}(q),
\]

where

\[
  \lambda,\mu\geq 0.
\]

The point is not the particular linear form.

The point is that provenance loss becomes an explicit optimization variable
rather than an accidental side effect.

Setting

\[
  \mu=0
\]

recovers a system indifferent to derivational distinctions except insofar as
they affect endpoint distortion.

A positive \(\mu\) makes some histories expensive to collapse even when their
current states are equivalent.

This is the computational expression of the essay's central objection.

% -----------------------------------------------------------------------------
%  B.21  WHY THE OPERATOR IS TASK-RELATIVE
% -----------------------------------------------------------------------------

\subsection{There Is No Universal Provenance Metric}
\label{app:no-universal-provenance-metric}

Nothing in the construction licenses a universal scalar

\[
  \lossorig(q)
\]

independent of a provenance vocabulary, task, probability model, or
admissibility criterion.

The same quotient may be harmless in one context and catastrophic in another.

A normalization that removes whitespace may be irrelevant to semantic
retrieval but unacceptable in forensic authorship analysis.

A paraphrase may preserve a theorem while destroying evidence about its
original formulation.

A reconstructed civilization may preserve institutions while losing
environmental relations required for world-continuation.

The correct expression is therefore normally

\[
  \lossorig(q;\delta),
\]

\[
  \lossorig(q;\Delta),
\]

or

\[
  \lossorig(q;F).
\]

The second argument is not optional decoration.

It states what one is refusing to forget.

% -----------------------------------------------------------------------------
%  B.22  THE FORMAL ORIGINIST OBJECTION
% -----------------------------------------------------------------------------

\subsection{The Objection as a Commutative Failure}
\label{app:commutative-failure}

The entire appendix can be summarized through a diagram.

Suppose we would like a provenance observable

\[
  \delta:\Gamma\to D
\]

to remain computable after representation by \(q\).

We seek a map

\[
  \bar\delta:Q\to D
\]

making

\[
\begin{tikzcd}[column sep=4em,row sep=2.5em]
\Gamma \arrow[r,"q"] \arrow[dr,"\delta"'] &
Q \arrow[d,dashed,"\bar\delta"] \\
& D
\end{tikzcd}
\]

commute.

If such a map exists, the provenance distinction survives the quotient.

If no such map exists, the dashed arrow cannot be filled.

That missing arrow is the wound.

\begin{held}[Originist Loss]
\centering
A representation loses a provenance distinction exactly when the desired
historical observable no longer factors through the representation.
\end{held}

This formulation is intentionally austere.

It covers perfect forgery, independent rediscovery, lossy compression,
generative reconstruction, branching descent, convergent histories, archival
restoration, and the document's own uncertain citation without requiring those
cases to be mathematically identical.

They share only the relevant structural failure:

\[
  \text{different histories}
  \longmapsto
  \text{one representation}.
\]

Once that identification has occurred, no operation internal to the
representation can discover which distinction was erased.

Something else must have survived elsewhere.

% ============================================================================
%  APPENDIX C
%  ADMISSIBILITY, GLUING, AND OBSTRUCTION
% ============================================================================

\section{Admissibility, Gluing, and Obstruction}
\label{app:admissibility}

\movementquestion{Why can every local step look acceptable while the resulting history still fails to constitute one coherent continuation?}

\begin{orientation}
Causal descent is too permissive. A corrupted file descends causally from the
file it corrupts. A malignant growth descends causally from healthy tissue. A
ruin descends causally from a city. None of these observations is sufficient
to establish the stronger relation that the essay has called continuation.

Admissibility supplies the missing restriction, but it creates a second
problem. A long history is rarely evaluated all at once. Different portions
of it may be inspected under different local constraints, by different
observers, or with access to different records. Every local segment can appear
acceptable while the collection fails to assemble into a globally coherent
history.

This appendix develops that possibility using a deliberately restrained
gluing formalism. The purpose is not to claim that every provenance problem is
literally a sheaf-cohomology problem. The purpose is to isolate a recurring
structural pattern: local realizations exist, neighboring realizations agree
partially, yet no single global realization satisfies all of the local data at
once.

That failure is especially important for TARTAN-like reasoning, Chain of
Memory, archival reconstruction, and world-continuation. It also explains why
checking transformations one at a time cannot in general certify an entire
trajectory.
\end{orientation}

% -----------------------------------------------------------------------------
%  C.1  STATE AND TRAJECTORY ADMISSIBILITY
% -----------------------------------------------------------------------------

\subsection{Two Kinds of Admissibility}
\label{app:two-admissibilities}

Let \(X\) be a state space and let

\[
  \Gamma(X)
\]

denote a collection of histories over \(X\).

A state-level admissibility predicate is a map

\[
  A_X:X\to\{0,1\}.
\]

The admissible state set is

\[
  \mathcal A_X
  =
  A_X^{-1}(1).
\]

A trajectory-level admissibility predicate is instead

\[
  A_\Gamma:\Gamma(X)\to\{0,1\},
\]

with

\[
  \mathcal A_\Gamma
  =
  A_\Gamma^{-1}(1).
\]

There is no general reason for the latter to be determined pointwise by the
former.

Indeed, one may have a history

\[
  \gamma
  =
  (
  x_0\to x_1\to\cdots\to x_n
  )
\]

such that

\[
  A_X(x_i)=1
  \qquad
  \text{for every }i,
\]

while

\[
  A_\Gamma(\gamma)=0.
\]

The simplest provenance example is a perfect forgery.

Every state-level property relevant to the artifact may be acceptable.

The history is nevertheless inadmissible for a task requiring authentic
descent.

\begin{proposition}[Pointwise admissibility is insufficient]
\label{prop:pointwise-insufficient}
Without an additional locality assumption on \(A_\Gamma\),

\[
  \bigl(
    A_X(x_i)=1
    \text{ for all }i
  \bigr)
  \centernot\Rightarrow
  A_\Gamma(\gamma)=1.
\]
\end{proposition}

The proposition is almost definitional, but its consequence is substantial.

A trajectory constraint cannot in general be replaced by repeated endpoint
checks.

% -----------------------------------------------------------------------------
%  C.2  TRANSITION ADMISSIBILITY
% -----------------------------------------------------------------------------

\subsection{Admissible Transitions}
\label{app:admissible-transitions}

Between state and trajectory admissibility lies a third possibility.

Let

\[
  A_E:X\times\mathcal T\times X\to\{0,1\}
\]

judge individual transformations.

For

\[
  x\xrightarrow{T}y,
\]

write

\[
  x\xrightarrow[\mathcal A]{T}y
\]

when

\[
  A_E(x,T,y)=1.
\]

A history is edgewise admissible when

\[
  A_E(x_{i-1},T_i,x_i)=1
\]

for every \(i\).

Again, edgewise admissibility need not imply trajectory admissibility.

A global constraint may depend upon accumulated drift, ordering, recurrence,
resource consumption, provenance depth, or a relation between distant points
of the path.

For example, suppose every transformation is allowed to change a scalar
quantity \(f\) by at most \(\varepsilon\):

\[
  |f(x_i)-f(x_{i-1})|
  \leq
  \varepsilon.
\]

Every edge may satisfy this bound while after \(n\) steps

\[
  |f(x_n)-f(x_0)|
  \gg
  \varepsilon.
\]

Local acceptability therefore does not imply global preservation.

% -----------------------------------------------------------------------------
%  C.3  ACCUMULATED DRIFT
% -----------------------------------------------------------------------------

\subsection{The Simplest Global Obstruction: Drift}
\label{app:drift}

Let

\[
  d:X\times X\to\mathbb R_{\geq 0}
\]

be a task-relative discrepancy function.

Suppose every edge of a trajectory satisfies

\[
  d(x_{i-1},x_i)\leq\varepsilon.
\]

If \(d\) obeys the triangle inequality, then only

\[
  d(x_0,x_n)
  \leq
  \sum_{i=1}^{n}d(x_{i-1},x_i)
  \leq
  n\varepsilon
\]

follows.

Thus an arbitrarily long trajectory can move arbitrarily far from its origin
while every local step remains small.

This is the elementary form of a problem that later appears as gluing failure.

Local compatibility constrains neighboring pieces.

It does not automatically establish global identity.

\begin{pathology}[Locally small, globally distant]
A trajectory may satisfy every local distortion bound while violating a
global continuation constraint relative to its origin.
\end{pathology}

The relevance to repeated regeneration is immediate.

A system can make only tiny changes per generation while eventually reaching a
state that no longer satisfies the invariant class that motivated the process.

% -----------------------------------------------------------------------------
%  C.4  LOCAL WINDOWS
% -----------------------------------------------------------------------------

\subsection{Histories Seen Through Windows}
\label{app:local-windows}

Let a history be parameterized over an interval

\[
  I=[0,1].
\]

Choose an open cover

\[
  \mathcal U
  =
  \{
    U_i
  \}_{i\in I_0}
\]

of the parameter domain.

Each \(U_i\) represents a local window into the history.

Suppose an observer reconstructs a local realization

\[
  s_i\in\mathscr F(U_i),
\]

where \(\mathscr F(U_i)\) denotes the admissible descriptions available over
that window.

Whenever

\[
  U_i\cap U_j\neq\varnothing,
\]

the two local realizations can be restricted to the overlap:

\[
  s_i|_{U_i\cap U_j},
  \qquad
  s_j|_{U_i\cap U_j}.
\]

The strongest local compatibility condition is

\[
  s_i|_{U_i\cap U_j}
  =
  s_j|_{U_i\cap U_j}.
\]

If all compatible local sections assemble uniquely into a global section

\[
  s\in\mathscr F(I),
\]

then the local data glue.

This is the ideal case.

% -----------------------------------------------------------------------------
%  C.5  WHY EXACT GLUING IS TOO STRONG
% -----------------------------------------------------------------------------

\subsection{Compatibility Up to Transformation}
\label{app:compatibility-up-to-transformation}

Provenance problems often do not provide literal equality on overlaps.

Two reconstructions may differ by a known translation, coordinate change,
normalization, renaming, or other admissible transformation.

Let \(G\) be a group, or more generally a groupoid of admissible
re-identifications.

On an overlap

\[
  U_i\cap U_j,
\]

suppose

\[
  s_i
  =
  g_{ij}\cdot s_j
\]

for some transition element

\[
  g_{ij}\in G.
\]

Consistency requires

\[
  g_{ii}=e
\]

and

\[
  g_{ji}=g_{ij}^{-1}.
\]

On triple overlaps, a further condition appears:

\[
  g_{ij}g_{jk}g_{ki}=e.
\]

When this condition fails, pairwise compatibility does not extend to coherent
three-way compatibility.

This is the first genuinely nonlocal obstruction.

% -----------------------------------------------------------------------------
%  C.6  THE COCYCLE CONDITION
% -----------------------------------------------------------------------------

\subsection{Cocycles}
\label{app:cocycles}

Assume for simplicity that the transition structure is represented by an
abelian group \(G\).

A collection

\[
  \{
    g_{ij}
  \}
\]

on pairwise overlaps satisfying

\[
  g_{ij}+g_{jk}+g_{ki}=0
\]

on every triple overlap is a Čech \(1\)-cocycle.

Using multiplicative notation, the same condition is

\[
  g_{ij}g_{jk}g_{ki}=e.
\]

A change of local representatives by elements

\[
  h_i\in G
\]

modifies transition data according to

\[
  g_{ij}
  \mapsto
  h_i^{-1}g_{ij}h_j.
\]

In the abelian case,

\[
  g_{ij}
  \mapsto
  g_{ij}+h_j-h_i.
\]

If the transition data can be eliminated entirely by such local
redefinitions, the cocycle is a coboundary.

Its cohomology class vanishes:

\[
  [g]=0
  \in
  \check H^1(\mathcal U,G).
\]

If

\[
  [g]\neq 0,
\]

the local pieces cannot be globally trivialized by the permitted local
re-identifications.

% -----------------------------------------------------------------------------
%  C.7  INTERPRETING THE OBSTRUCTION
% -----------------------------------------------------------------------------

\subsection{What a Nonzero Class Means Here}
\label{app:interpret-h1}

The formal analogy to provenance should be used carefully.

The statement

\[
  [g]\neq 0
\]

does not mean simply that the local data are false.

Each local realization may be individually admissible.

Each pairwise relation may also be legitimate.

The obstruction says that no single choice of representatives makes all of
those local relations simultaneously disappear.

In the present framework, that means there is no globally coherent
identification of all local realizations under the selected transition
structure.

This is stronger than local disagreement.

It is an obstruction to global continuation.

\begin{definition}[Gluing obstruction]
\label{def:gluing-obstruction}
Given locally admissible realizations \(\{s_i\}\) and transition data
\(\{g_{ij}\}\), a \emph{gluing obstruction} is a nontrivial equivalence class
of transition data preventing the local realizations from being represented as
restrictions of one globally coherent realization under the permitted
re-identifications.
\end{definition}

The definition is deliberately broader than the abelian Čech model.

The latter is one tractable realization of the idea.

% -----------------------------------------------------------------------------
%  C.8  PROVENANCE GLUING
% -----------------------------------------------------------------------------

\subsection{A Provenance Example}
\label{app:provenance-gluing}

Suppose three archival fragments \(A\), \(B\), and \(C\) overlap pairwise.

Fragment \(A\) can be reconciled with \(B\) by a transformation

\[
  g_{AB}.
\]

Fragment \(B\) can be reconciled with \(C\) by

\[
  g_{BC}.
\]

Fragment \(C\) can be reconciled with \(A\) by

\[
  g_{CA}.
\]

Pairwise reconciliation therefore succeeds.

But if

\[
  g_{AB}g_{BC}g_{CA}\neq e,
\]

the three reconciliations cannot all describe one globally consistent
derivation under the chosen transformation model.

One of several things may be true.

The fragments may derive from different branches.

One transition model may be wrong.

A missing transformation may have occurred.

The assumed common origin may be false.

Or the representation may be too coarse to express the distinction required
for gluing.

The obstruction does not decide among these explanations.

It establishes that the current local account is globally insufficient.

% -----------------------------------------------------------------------------
%  C.9  CHAIN OF MEMORY AS A GLUING PROBLEM
% -----------------------------------------------------------------------------

\subsection{Chain of Memory and Local Consistency}
\label{app:com-gluing}

A Chain of Memory system can be modeled as retaining a family of local
derivational records

\[
  s_i\in\mathscr F(U_i).
\]

If each new context window preserves only compatibility with the immediately
preceding window, then the system verifies conditions of the form

\[
  s_i|_{U_i\cap U_{i+1}}
  \sim
  s_{i+1}|_{U_i\cap U_{i+1}}.
\]

This is not enough to guarantee global consistency.

Long-range constraints may fail.

The system can therefore exhibit what may be called a \emph{CLIO stall}: local
repair continues to produce acceptable neighboring states, but the global
obstruction does not decrease.

Let

\[
  \Omega_n
\]

denote a measure or class of unresolved gluing obstruction after generation
\(n\).

A stalled regime has schematically

\[
  \Omega_{n+1}\not\prec\Omega_n
\]

despite continuing local repair.

In a cohomological realization this may appear as persistent nonvanishing

\[
  [g_n]\neq 0
\]

across successive repair cycles.

The system is locally active but globally stuck.

% -----------------------------------------------------------------------------
%  C.10  REPAIR CAN MOVE THE OBSTRUCTION
% -----------------------------------------------------------------------------

\subsection{Repair Without Resolution}
\label{app:repair-moves-obstruction}

Suppose a repair operation changes local representatives

\[
  s_i\mapsto s_i'.
\]

This induces new transition data

\[
  g_{ij}\mapsto g_{ij}'.
\]

A repair can make some overlaps easier to reconcile while making others
harder.

Thus a decrease in a local mismatch functional

\[
  \mathcal L_{\mathrm{local}}
\]

does not imply disappearance of the global obstruction.

One may have

\[
  \mathcal L_{\mathrm{local}}^{(n+1)}
  <
  \mathcal L_{\mathrm{local}}^{(n)}
\]

while

\[
  [g_{n+1}]
  =
  [g_n]
  \neq 0.
\]

The representation looks progressively smoother.

The obstruction survives.

This is an important pathology for generative repair.

Optimization of local coherence can conceal persistence of a global
incompatibility.

% -----------------------------------------------------------------------------
%  C.11  LOOP CLOSURE
% -----------------------------------------------------------------------------

\subsection{Loop Closure as a Global Test}
\label{app:loop-closure}

Suppose a trajectory returns to a region equivalent to its origin under an
invariant relation:

\[
  x_n\simI x_0.
\]

This establishes a form of endpoint recurrence.

It does not guarantee that the intervening history closes consistently.

Let the path traverse local charts

\[
  U_1,U_2,\ldots,U_n
\]

with transition maps

\[
  g_{12},
  g_{23},
  \ldots,
  g_{n1}.
\]

The holonomy around the loop is

\[
  H(\gamma)
  =
  g_{12}
  g_{23}
  \cdots
  g_{n1}.
\]

If

\[
  H(\gamma)=e,
\]

the local identifications close trivially under the chosen structure.

If

\[
  H(\gamma)\neq e,
\]

the system returns to an apparently equivalent state while carrying a residual
transformation accumulated around the loop.

This gives a more exact version of

\[
  O\rightsquigarrow X\rightsquigarrow O'.
\]

The relation

\[
  O'\simI O
\]

does not imply trivial historical holonomy.

% -----------------------------------------------------------------------------
%  C.12  LE GUIN AND NONTRIVIAL RETURN
% -----------------------------------------------------------------------------

\subsection{Return with Holonomy}
\label{app:leguin-holonomy}

The metaphorical usefulness of holonomy for \emph{The Dispossessed} is now
clear.

Shevek's return to Anarres is not interesting because the journey forms a
geographical loop.

It is interesting because traversal changes the relation between traveler,
origin, knowledge, and political possibility \cite{leguin1974}.

Schematically,

\[
  O
  \rightsquigarrow
  X
  \rightsquigarrow
  O'
\]

with

\[
  O'\sim_{\mathcal C}O
\]

but

\[
  O'\neq O.
\]

If the return carries a residual transformation

\[
  H(\gamma)\neq e,
\]

then recurrence is not restoration.

The loop closes spatially or constraint-relatively while remaining
historically nontrivial.

This is precisely the kind of case for which trajectory information survives
after endpoint recurrence.

% -----------------------------------------------------------------------------
%  C.13  LOCAL ADMISSIBILITY AND GLOBAL WORLDHOOD
% -----------------------------------------------------------------------------

\subsection{World-Continuation as a Gluing Problem}
\label{app:world-gluing}

The same structure can be lifted from artifacts to worlds.

Suppose a world is locally realized through overlapping domains

\[
  U_i
\]

with local structures

\[
  W_i.
\]

Each \(W_i\) may be internally coherent.

Neighboring domains may admit compatible translations.

Yet global worldhood requires more than a collection of locally acceptable
pieces.

They must glue into a structure supporting coherent cross-domain relations.

A Macrolife civilization could preserve local organizational patterns across
many habitats while the total system no longer admits the global relations
that characterized its ancestral world.

Thus

\[
  \forall i,\quad
  W_i^{(0)}
  \sim
  W_i^{(n)}
\]

does not automatically imply

\[
  W_0\eqworld W_n.
\]

World-continuation can fail globally while succeeding locally.

This is one reason the No-World problem cannot be reduced to persistence of
objects within a presumed container.

% -----------------------------------------------------------------------------
%  C.14  ADMISSIBILITY AS A PATH-SPACE SUBSET
% -----------------------------------------------------------------------------

\subsection{The Admissible Path Space}
\label{app:admissible-path-space}

Return now to the path space

\[
  \Gamma(X).
\]

Let

\[
  \mathcal A_\Gamma
  \subseteq
  \Gamma(X)
\]

denote the admissible histories.

For endpoints \(x,y\in X\), define

\[
  \Gamma_{\mathcal A}(x,y)
  =
  \{
    \gamma\in\mathcal A_\Gamma:
    \Orig(\gamma)=x,\;
    \End(\gamma)=y
  \}.
\]

Then continuation from \(x\) to \(y\) is possible under the selected
constraints exactly when

\[
  \Gamma_{\mathcal A}(x,y)\neq\varnothing.
\]

This is stronger than causal reachability.

If

\[
  \Gamma(x,y)\neq\varnothing
\]

but

\[
  \Gamma_{\mathcal A}(x,y)=\varnothing,
\]

then \(y\) is reachable from \(x\) but not by any admissible history.

This distinction formalizes the difference between descent and acceptable
continuation.

% -----------------------------------------------------------------------------
%  C.15  MULTIPLE ADMISSIBLE CLASSES
% -----------------------------------------------------------------------------

\subsection{Continuation Need Not Be Unique}
\label{app:multiple-admissible-classes}

Even when

\[
  \Gamma_{\mathcal A}(x,y)\neq\varnothing,
\]

there may be several inequivalent admissible histories.

Suppose

\[
  \gamma_1,\gamma_2
  \in
  \Gamma_{\mathcal A}(x,y)
\]

while

\[
  \gamma_1\not\sim_{\mathcal A}\gamma_2.
\]

Then admissibility establishes that both histories are permitted.

It does not establish that they are the same continuation.

This matters especially under branching.

A civilization may possess several admissible descendants.

A theory may have several independently reconstructed formulations.

An archive may support several histories consistent with all surviving
evidence.

Admissibility therefore constrains possibility without necessarily selecting
uniqueness.

% -----------------------------------------------------------------------------
%  C.16  AN OBSTRUCTION FUNCTIONAL
% -----------------------------------------------------------------------------

\subsection{A Minimal Obstruction Functional}
\label{app:obstruction-functional}

For applications not requiring full cohomological machinery, define a
nonnegative functional

\[
  \mathfrak O:
  \Gamma(X)\to\mathbb R_{\geq 0}
\]

such that

\[
  \mathfrak O(\gamma)=0
\]

when all required local compatibility conditions glue globally, and

\[
  \mathfrak O(\gamma)>0
\]

when some unresolved incompatibility remains.

The precise functional is domain-specific.

It may combine invariant mismatch, provenance inconsistency, boundary
violation, or failed recurrence.

A repair trajectory

\[
  \gamma_0,
  \gamma_1,
  \gamma_2,
  \ldots
\]

can then be evaluated by

\[
  \mathfrak O(\gamma_n).
\]

Successful repair ideally produces

\[
  \mathfrak O(\gamma_{n+1})
  <
  \mathfrak O(\gamma_n).
\]

A stall occurs when local losses decrease while

\[
  \mathfrak O(\gamma_n)
\]

fails to approach zero.

This is the operational analogue of persistent nontrivial gluing obstruction.

% -----------------------------------------------------------------------------
%  C.17  TARTAN INTERPRETATION
% -----------------------------------------------------------------------------

\subsection{TARTAN as Constraint-Preserving Realization}
\label{app:tartan-interpretation}

Within the present essay, TARTAN can therefore be read as a framework for
distinguishing realization from mere resemblance.

A state \(x\) does not count as an acceptable realization merely because it
matches selected features.

A realization must inhabit an admissibility structure.

Likewise, a trajectory does not count as continuation merely because every
neighboring pair is causally related.

The entire path must remain compatible with the constraints whose preservation
defines the relevant form of descent.

This suggests the schematic condition

\[
  \Continuation_{\mathcal A}(x,y)
  \quad\Longleftrightarrow\quad
  \Gamma_{\mathcal A}(x,y)\neq\varnothing.
\]

When global compatibility is additionally required, one may strengthen this
to

\[
  \Continuation_{\mathcal A,\mathfrak O}(x,y)
  \quad\Longleftrightarrow\quad
  \exists\gamma\in\Gamma_{\mathcal A}(x,y)
  \text{ such that }
  \mathfrak O(\gamma)=0.
\]

This is not proposed as a universal definition of identity.

It is a formal model of the narrower claim that continuation is constrained
reachability rather than endpoint resemblance.

% -----------------------------------------------------------------------------
%  C.18  ORIGINIST LOSS AND GLUING OBSTRUCTION
% -----------------------------------------------------------------------------

\subsection{Two Different Failures}
\label{app:loss-vs-obstruction}

Originist loss and gluing obstruction should not be conflated.

Originist loss concerns distinctions erased by a representation:

\[
  \lossorig(q;\delta)>0.
\]

Gluing obstruction concerns local data that fail to assemble globally:

\[
  \mathfrak O(\gamma)>0
\]

or, in the cohomological model,

\[
  [g]\neq 0.
\]

The two can interact.

A representation may erase exactly the distinction required to diagnose a
gluing failure.

Conversely, a persistent gluing obstruction may motivate preserving more
provenance information.

But neither implies the other in general.

One is a failure of recoverability.

The other is a failure of global compatibility.

Their separation is essential to the larger architecture.

% -----------------------------------------------------------------------------
%  C.19  A COMMUTATIVE FORM OF CONTINUATION
% -----------------------------------------------------------------------------

\subsection{Continuation as Successful Gluing}
\label{app:continuation-gluing}

The desired situation can be summarized schematically.

Local histories

\[
  \gamma_i\in\mathscr F(U_i)
\]

should admit a global history

\[
  \gamma\in\mathscr F(U)
\]

such that

\[
  \gamma|_{U_i}
  \sim
  \gamma_i
\]

for every \(i\).

Diagrammatically,

\[
\begin{tikzcd}[column sep=4em,row sep=2.5em]
& \gamma \arrow[dl] \arrow[d] \arrow[dr] & \\
\gamma_1 & \gamma_2 & \gamma_3
\end{tikzcd}
\]

represents a common global realization whose restrictions recover the local
ones up to the permitted equivalence.

When no such \(\gamma\) exists, the local descriptions may remain individually
useful.

What fails is their claim to constitute one history.

This distinction will matter directly in Appendix~\ref{app:worldhood}, where
the object to be glued is no longer merely a record or trajectory but a
candidate world.

% -----------------------------------------------------------------------------
%  C.20  THE ADMISSIBILITY PRINCIPLE
% -----------------------------------------------------------------------------

\subsection{The Result Needed by the Main Argument}
\label{app:admissibility-principle}

The technical machinery of this appendix can be condensed into one principle.

\begin{theorem}[Admissibility does not localize automatically]
\label{thm:admissibility-does-not-localize}
In the absence of a theorem establishing that a global admissibility predicate
is determined by local restrictions, verification of every local state,
transition, or overlap is insufficient to establish global trajectory
admissibility.
\end{theorem}

\begin{proof}
It suffices to exhibit a class of counterexamples. Let local realizations
\(\{s_i\}\) be individually admissible and pairwise related by transition data
\(\{g_{ij}\}\). If the resulting cocycle represents a nontrivial obstruction,

\[
  [g]\neq 0,
\]

then no permitted choice of local representatives produces a single globally
compatible realization. Thus every stipulated local test may succeed while
the global gluing condition fails.
\end{proof}

The theorem identifies the precise burden of proof hidden by claims of local
continuity.

One must either demonstrate a valid local-to-global principle or retain the
possibility of obstruction.

\begin{held}[Local-to-Global Warning]
\centering
Every step may be admissible.\\
Every overlap may be repairable.\\
The whole may still fail to constitute one continuation.
\end{held}

That warning is the mathematical counterpart of the essay's literary
triangulation.

Card supplies incomplete local traces of an absent past.

Zebrowski supplies indefinitely extended local transformations toward a
future.

Le Guin supplies the loop in which departure and return test whether locally
continuous history closes globally.

The next appendix asks what happens when the thing whose coherence must be
tested is not an artifact, a person, or a civilization considered in
isolation, but a world.

% ============================================================================
%  APPENDIX D
%  WORLD-CONTINUATION AND THE NO-WORLD FORMULATION
% ============================================================================

\section{World-Continuation and the No-World Formulation}
\label{app:worldhood}

\movementquestion{What could it mean for a world to continue if worldhood is not assumed in advance as a container within which continuing things reside?}

\begin{orientation}
Ordinary persistence questions begin too late. They assume a world and then
ask whether some object inside it has remained the same. That order is useful
for everyday reasoning, but it conceals the harder question raised by the
cases considered in this essay. If objects, organisms, archives, institutions,
habitats, and civilizations can all be replaced, reconstructed, distributed,
or moved across substrates, then the persistence of their surrounding world
cannot simply be taken for granted.

The No-World formulation reverses the explanatory order. A world is not first
posited as a primitive container. Worldhood is reconstructed from a sufficiently
coherent organization of relations, constraints, recurrences, boundaries, and
possible continuations. What persists, on this view, need not be a material
inventory. But neither is persistence secured merely by reproducing an abstract
pattern somewhere else.

This appendix formalizes only the portion of that thesis required by the
present essay. It distinguishes object-continuation, civilization-continuation,
and world-continuation; defines a world signature as a structured relational
object rather than a list of contents; and shows why preservation of every
selected component need not imply preservation of the world constituted by
their relations.

The result is especially important for \emph{Macrolife}. A civilization may
continue through radical material transformation while the question of whether
its world continues remains independently open \cite{zebrowski1979}.
\end{orientation}

% -----------------------------------------------------------------------------
%  D.1  FROM OBJECTS TO RELATIONAL REALIZATIONS
% -----------------------------------------------------------------------------

\subsection{Relational Realizations}
\label{app:relational-realizations}

Let \(X\) denote a collection of realized entities.

A purely extensional description of a state would specify only a set

\[
  X_t
  =
  \{
    x_1,\ldots,x_n
  \}.
\]

For worldhood this is insufficient.

Introduce a family of relations

\[
  \mathcal R_t
  =
  \{
    R_t^{(1)},
    R_t^{(2)},
    \ldots,
    R_t^{(k)}
  \},
\]

a family of constraints

\[
  \mathcal C_t
  =
  \{
    C_t^{(1)},
    C_t^{(2)},
    \ldots,
    C_t^{(m)}
  \},
\]

and a family of admissible transformations

\[
  \mathcal A_t.
\]

Define a relational realization by

\[
  W_t
  =
  \left(
    X_t,
    \mathcal R_t,
    \mathcal C_t,
    \mathcal A_t
  \right).
\]

The notation \(W_t\) should not yet be read as asserting that a world exists.

It denotes candidate world-constituting structure.

Worldhood is the further claim that this structure closes sufficiently to
support the distinctions, recurrences, and continuations by which its
constituents are jointly realized.

% -----------------------------------------------------------------------------
%  D.2  WORLD SIGNATURES
% -----------------------------------------------------------------------------

\subsection{World Signatures}
\label{app:world-signatures}

Let

\[
  \Sigma(W)
\]

denote a world signature extracted from a relational realization \(W\).

A useful schematic form is

\[
  \Sigma(W)
  =
  \left(
    \mathcal B,
    \mathcal R,
    \mathcal C,
    \mathcal P,
    \mathcal Q
  \right),
\]

where \(\mathcal B\) records boundary structure, \(\mathcal R\) relational
organization, \(\mathcal C\) operative constraints, \(\mathcal P\) admissible
path structure, and \(\mathcal Q\) recurrence or closure structure.

The components need not be represented identically in every application.

The important point is that

\[
  \Sigma(W)
\]

is relational.

It is not merely an inventory of objects.

Two candidate worlds may therefore contain equivalent components while
possessing inequivalent world signatures.

% -----------------------------------------------------------------------------
%  D.3  COMPONENT EQUIVALENCE DOES NOT IMPLY WORLD EQUIVALENCE
% -----------------------------------------------------------------------------

\subsection{The Failure of Componentwise Persistence}
\label{app:componentwise-failure}

Suppose

\[
  W
  =
  (X,\mathcal R,\mathcal C,\mathcal A)
\]

and

\[
  W'
  =
  (X',\mathcal R',\mathcal C',\mathcal A').
\]

Assume there exists a bijection

\[
  f:X\to X'
\]

such that every object in \(X\) has a corresponding object in \(X'\).

This alone does not imply

\[
  W\eqworld W'.
\]

The relations may differ.

The constraints may differ.

The available trajectories may differ.

The boundary conditions under which the entities jointly realize one another
may differ.

\begin{proposition}[Component preservation is insufficient]
\label{prop:component-preservation-insufficient}
A bijection between the constituent sets of two relational realizations is not
sufficient to establish world-equivalence.
\end{proposition}

\begin{proof}
Choose

\[
  X=\{a,b\}
\]

and

\[
  X'=\{a',b'\},
\]

with bijection

\[
  f(a)=a',
  \qquad
  f(b)=b'.
\]

Let \(R\) hold between \(a\) and \(b\), while the corresponding relation
\(R'\) does not hold between \(a'\) and \(b'\). The constituent sets are
bijectively matched, but their relational structures are not preserved.
Therefore constituent correspondence alone cannot establish equivalence of
the complete relational realizations.
\end{proof}

The result is elementary.

Its philosophical consequence is not.

Replacing every component of a world with an individually adequate
counterpart does not by itself replace the world.

% -----------------------------------------------------------------------------
%  D.4  OBJECT, CIVILIZATION, AND WORLD CONTINUATION
% -----------------------------------------------------------------------------

\subsection{Three Persistence Questions}
\label{app:three-persistence-questions}

Let

\[
  o_t
\]

denote an object,

\[
  C_t
\]

a civilization, and

\[
  W_t
\]

a candidate world at time \(t\).

Define three distinct continuation predicates:

\[
  \Cont_O(o_t,o_{t'})
\]

for object-continuation,

\[
  \Cont_C(C_t,C_{t'})
\]

for civilization-continuation, and

\[
  \Cont_W(W_t,W_{t'})
\]

for world-continuation.

No implication among these predicates should be assumed without argument.

In particular,

\[
  \Cont_O(o_t,o_{t'})
  \centernot\Rightarrow
  \Cont_C(C_t,C_{t'}),
\]

\[
  \Cont_C(C_t,C_{t'})
  \centernot\Rightarrow
  \Cont_W(W_t,W_{t'}),
\]

and

\[
  \Cont_W(W_t,W_{t'})
  \centernot\Rightarrow
  \Cont_O(o_t,o_{t'})
\]

in general.

A world may continue while particular objects disappear.

A civilization may continue while the ecological and relational world from
which it arose does not.

An object may persist through the collapse of the civilization that once gave
it significance.

These are different questions because their admissibility constraints are
different.

% -----------------------------------------------------------------------------
%  D.5  MACROLIFE AS A SEPARATING CASE
% -----------------------------------------------------------------------------

\subsection{\emph{Macrolife} and the Separation of Scales}
\label{app:macrolife-separation}

Zebrowski's \emph{Macrolife} supplies an unusually strong thought experiment
because the relevant unit of persistence expands beyond the organism and even
beyond the planetary society \cite{zebrowski1979}.

Suppose a civilization undergoes a trajectory

\[
  C_0
  \leadcausal
  C_1
  \leadcausal
  \cdots
  \leadcausal
  C_n
\]

through migration, habitat construction, technological transformation, and
changes of material substrate.

It may be defensible that

\[
  \Cont_C(C_0,C_n)=1.
\]

But let the ancestral world be

\[
  W_0
  =
  (
    X_0,
    \mathcal R_0,
    \mathcal C_0,
    \mathcal A_0
  \right)
\]

and the later world be

\[
  W_n
  =
  (
    X_n,
    \mathcal R_n,
    \mathcal C_n,
    \mathcal A_n
  \right).
\]

Civilizational continuation does not settle whether

\[
  \Cont_W(W_0,W_n)=1.
\]

The civilization may have survived precisely by ceasing to participate in the
relations that constituted its former world.

The success of one continuation can therefore coincide with the failure of
another.

% -----------------------------------------------------------------------------
%  D.6  WORLD-EQUIVALENCE
% -----------------------------------------------------------------------------

\subsection{Constraint-Relative World-Equivalence}
\label{app:constraint-world-equivalence}

Let

\[
  \mathfrak W
\]

be a space of candidate world realizations.

Define a task-relative relation

\[
  \eqworld^{\mathcal K}
\]

by

\[
  W\eqworld^{\mathcal K}W'
\]

when \(W\) and \(W'\) preserve the world-constituting structures selected by a
constraint family \(\mathcal K\).

The superscript is important.

There is no reason to expect world-equivalence to be absolute.

For one inquiry, preserving causal accessibility may suffice.

For another, ecological reciprocity may be essential.

For another, the relevant structure may include institutional memory,
boundary relations, or recurrence classes.

Thus

\[
  W\eqworld^{\mathcal K_1}W'
\]

need not imply

\[
  W\eqworld^{\mathcal K_2}W'.
\]

Worldhood, like provenance, is typed by the distinctions one requires the
representation to preserve.

% -----------------------------------------------------------------------------
%  D.7  WORLD SIGNATURE PRESERVATION
% -----------------------------------------------------------------------------

\subsection{Signature-Preserving Maps}
\label{app:signature-preserving-maps}

Let

\[
  F:W\to W'
\]

be a transformation between candidate world realizations.

Suppose there exists an induced map

\[
  F_\Sigma:\Sigma(W)\to\Sigma(W').
\]

\begin{remark}[Inducedness is assumed, not derived]
\label{rem:F-sigma-assumed}
\(\Sigma\) is given schematically in \cref{app:world-signatures} as a
5-tuple and is not equipped here with a functorial construction. The
existence of \(F_\Sigma\) for a given \(F\) is therefore a standing
assumption of this subsection, component by component, rather than a
consequence derived from the definition of \(\Sigma\) itself. A fully
general construction of \(F_\Sigma\) from \(F\) is left open.
\end{remark}

A strong form of world preservation would require \(F_\Sigma\) to preserve all
relations designated by \(\mathcal K\).

Schematically,

\[
  \Sigma_{\mathcal K}(W)
  \simeq
  \Sigma_{\mathcal K}(W').
\]

This condition is stronger than preservation of the objects in \(W\).

It is also weaker than literal identity of every constituent.

\begin{remark}[\(\mathcal P\) is local; \(\mathcal A_W\) is global]
\label{rem:P-vs-AW}
One component of \(\Sigma(W)\) is \(\mathcal P\), the admissible path
structure available at \(W\) — this is a per-state, local property, in the
sense of \(\mathcal A_X\) from \cref{sec:tartan-admissibility}: which
elementary transformations are admissible starting from this state. It is
not \(\mathcal A_W\), the global, trajectory-level property of which
complete histories count as admissible world-continuations. Preserving
\(\mathcal P\) under \(F_\Sigma\) says something about the local menu of
transitions available at corresponding states; it does not by itself place
any history through those states inside \(\mathcal A_W\), for the same
reason \(\gamma(t)\in\mathcal A_X\) for all \(t\) does not imply
\(\gamma\in\mathcal A_\Gamma\).
\end{remark}

This supplies an intermediate notion of world-structural preservation
appropriate to a theory in which nothing original need survive. It does not,
by itself, establish world-continuation. Two terminal realizations may
preserve the same \(\mathcal K\)-selected world signature while differing in
provenance, or while failing to be connected by any admissible world
history: continuation requires the further trajectory-level condition
introduced in \cref{app:world-path-space} below, and

\[
  \Sigma_{\mathcal K}(W_0)
  \simeq
  \Sigma_{\mathcal K}(W_1)
  \centernot\Rightarrow
  \Cont_W(W_0,W_1).
\]

% -----------------------------------------------------------------------------
%  D.8  THE REPLACEMENT THEOREM
% -----------------------------------------------------------------------------

\subsection{Total Material Replacement}
\label{app:total-replacement}

Suppose a world trajectory

\[
  \Gamma_W
  =
  (
    W_0\to W_1\to\cdots\to W_n
  )
\]

replaces every original material constituent.

Let

\[
  M(W_i)
\]

denote the material inventory of \(W_i\).

Assume

\[
  M(W_0)\cap M(W_n)=\varnothing.
\]

This fact alone proves neither world-discontinuity nor world-continuation.

\begin{proposition}[Material overlap is neither necessary nor sufficient]
\label{prop:material-overlap}
Within a relational account of worldhood, nonzero material overlap is neither
a necessary nor a sufficient condition for world-continuation.
\end{proposition}

\begin{proof}
For non-necessity, consider a trajectory in which every material component is
replaced while all world-defining relations designated by \(\mathcal K\) are
continuously preserved and the trajectory remains within the admissible
world-history family \(\mathcal A_W\) introduced in
\cref{app:world-path-space} below. Under such a criterion, continuation may
hold despite

\[
  M(W_0)\cap M(W_n)=\varnothing.
\]

For non-sufficiency, preserve a single material component while destroying the
relations designated by \(\mathcal K\) or leaving the admissible family
\(\mathcal A_W\). Material overlap remains nonempty, but the selected
world-continuation criterion fails.
\end{proof}

This is the formal core of the title's claim.

The word \enquote{need} matters.

The proposition does not say that material continuity is irrelevant.

It says that material overlap cannot serve as a universal criterion.

% -----------------------------------------------------------------------------
%  D.9  CONTINUATION AS AN ADMISSIBLE WORLD TRAJECTORY
% -----------------------------------------------------------------------------

\subsection{The World Path Space}
\label{app:world-path-space}

Let

\[
  \Gamma(\mathfrak W)
\]

denote histories of candidate world realizations.

Let

\[
  \mathcal A_W
  \subseteq
  \Gamma(\mathfrak W)
\]

be the subset satisfying the selected world-continuation constraints.

\begin{remark}[\(\mathcal K\) and \(\mathcal A_W\) are distinct constraint families]
\label{rem:K-vs-AW}
\(\mathcal K\) (\cref{app:constraint-world-equivalence}) selects which
\emph{states} count as world-equivalent; \(\mathcal A_W\) selects which
\emph{histories} count as admissible world-continuations. Nothing in their
definitions identifies one with the other, and a trajectory can satisfy
either without satisfying the other: a history can preserve every
\(\mathcal K\)-designated relation at each endpoint while failing to lie in
\(\mathcal A_W\), and conversely. Any stated relationship between them —
for instance, that \(\mathcal A_W\) consists exactly of the histories along
which \(\mathcal K\)-equivalence holds at every intermediate state — is a
bridge principle in the sense of \cref{thm:persistence-separation}, and is
not assumed here.
\end{remark}

For \(W_0,W_1\in\mathfrak W\), define

\[
  \Gamma_{\mathcal A_W}(W_0,W_1)
  =
  \{
    \gamma\in\mathcal A_W:
    \Orig(\gamma)=W_0,\;
    \End(\gamma)=W_1
  \}.
\]

Then a trajectory-based world-continuation predicate may be defined by

\[
  \Cont_W(W_0,W_1)
  =
  1
\]

exactly when

\[
  \Gamma_{\mathcal A_W}(W_0,W_1)
  \neq
  \varnothing.
\]

The definition makes world-continuation a property of admissible histories
between world realizations.

It is not a binary comparison of isolated endpoints.

% -----------------------------------------------------------------------------
%  D.10  EQUAL WORLDS WITH DIFFERENT HISTORIES
% -----------------------------------------------------------------------------

\subsection{World-Level Convergent Reconstruction}
\label{app:world-convergent-reconstruction}

The Originist problem now reappears at world scale.

Suppose

\[
  W_A\eqworld^{\mathcal K}W_B.
\]

It does not follow that their histories are equivalent.

One world may be the uninterrupted continuation of an earlier realization.

Another may have been independently reconstructed from records.

A third may be a deliberate simulation.

A fourth may have converged upon the same organization without any common
genealogical path.

Thus

\[
  W_A\eqworld^{\mathcal K}W_B
\]

can coexist with

\[
  \traj{W_A}\neq\traj{W_B}.
\]

This is the world-level analogue of the essay's original case

\[
  x_A=x_B,
  \qquad
  \gamma_A\neq\gamma_B.
\]

The same distinction has survived every increase in scale.

% -----------------------------------------------------------------------------
%  D.11  WORLD PROVENANCE
% -----------------------------------------------------------------------------

\subsection{Worlds Have Provenance}
\label{app:world-provenance}

If worldhood is reconstructed from relational organization, then candidate
worlds themselves inherit provenance.

Define

\[
  \Prov_W:
  \Gamma(\mathfrak W)
  \to
  D_W
\]

as a world-provenance observable.

Its values may encode ancestral worlds, transformation classes, branching
events, reconstruction episodes, or uncertainty over these.

Two worlds can therefore satisfy

\[
  \Sigma_{\mathcal K}(W_A)
  =
  \Sigma_{\mathcal K}(W_B)
\]

while

\[
  \Prov_W(\traj{W_A})
  \neq
  \Prov_W(\traj{W_B}).
\]

Consequently,

\[
  \text{world-equivalence}
  \neq
  \text{world-provenance equivalence}.
\]

The distinction is not an exotic extension of the Originist thesis.

It is that thesis applied consistently.

% -----------------------------------------------------------------------------
%  D.12  THE NO-WORLD MOVE
% -----------------------------------------------------------------------------

\subsection{Worldhood Is Not Presupposed}
\label{app:no-world-move}

The expression \enquote{No-World} can now be stated negatively and positively.

Negatively, the framework refuses the inference

\[
  \text{entities exist}
  \quad\Rightarrow\quad
  \text{one determinate world containing them is primitive}.
\]

Positively, it treats worldhood as a successful organization of relations,
constraints, boundaries, and admissible continuations.

Let

\[
  \World:
  \mathfrak W\to\{0,1\}
\]

be a worldhood predicate.

Then

\[
  \World(W)=1
\]

only when \(W\) satisfies the selected closure conditions.

The exact closure conditions belong to the underlying theory and need not be
fixed universally here.

What matters is the order of explanation:

\[
  \text{relational realization}
  \longrightarrow
  \text{closure}
  \longrightarrow
  \text{worldhood},
\]

rather than

\[
  \text{world}
  \longrightarrow
  \text{relations inside it}.
\]

% -----------------------------------------------------------------------------
%  D.13  WORLDHOOD AS AN INFERRED ORIGIN
% -----------------------------------------------------------------------------

\subsection{The Originist Reading of Worldhood}
\label{app:worldhood-originist}

This reversal reveals the deeper relation to Card.

An Originist begins with surviving traces and attempts to infer a structure
that is no longer directly available \cite{card1989}.

A No-World theorist begins with local realizations, recurrences, constraints,
and relations and asks whether they warrant reconstruction of a coherent
world.

The formal shapes are similar.

Let \(E\) denote available evidence and let

\[
  \mathcal W(E)
\]

be the class of candidate worlds compatible with it.

Then world reconstruction is an inverse problem:

\[
  E
  \rightsquigarrow
  \mathcal W(E).
\]

If

\[
  |\mathcal W(E)|>1,
\]

the evidence underdetermines worldhood.

A single inferred world may still be selected by additional constraints or
priors.

But selection should not be confused with direct possession of the world as a
primitive datum.

% -----------------------------------------------------------------------------
%  D.14  BRANCHING WORLDS
% -----------------------------------------------------------------------------

\subsection{Branching World-Continuation}
\label{app:branching-worlds}

Suppose a world realization \(W_0\) admits two admissible continuations:

\[
  W_0
  \leadcausal
  W_A
\]

and

\[
  W_0
  \leadcausal
  W_B.
\]

Assume

\[
  W_A\neq W_B.
\]

If both trajectories satisfy the selected world-admissibility conditions, then
world-continuation branches.

The resulting structure is not a single line

\[
  W_0\to W_1\to W_2,
\]

but a provenance graph

\[
  \graphG_W.
\]

Worldhood therefore inherits the branching problem established for personal
and civilizational continuation.

A unique future world cannot be inferred merely from continuity with the
present.

% -----------------------------------------------------------------------------
%  D.15  DISSEMINATION AND MULTIPLE WORLDS
% -----------------------------------------------------------------------------

\subsection{Le Guin's Dissemination Problem}
\label{app:leguin-worlds}

The dissemination of Shevek's work in \emph{The Dispossessed} supplies a
different kind of branching \cite{leguin1974}.

Knowledge leaves one political world and becomes available beyond the
institutions that previously constrained its circulation.

The relevant trajectory is therefore not merely

\[
  \gamma:
  O\to X\to O'.
\]

It also branches:

\[
  \gamma
  \longrightarrow
  \{
    \gamma_1,
    \gamma_2,
    \ldots
  \}.
\]

Once the knowledge becomes multiply instantiated, no single receiving world
exhausts its continuation.

This matters because dissemination can preserve content while breaking the
exclusive provenance relation through which one institution previously
controlled it.

The return and the branching are therefore coupled.

The origin remains part of the history.

It ceases to be the sole container of the future.

% -----------------------------------------------------------------------------
%  D.16  WORLD CONTINUATION AND RECURRENCE
% -----------------------------------------------------------------------------

\subsection{Recurrence Without Restoration}
\label{app:world-recurrence}

Suppose a world trajectory satisfies

\[
  W_0
  \rightsquigarrow
  W_1
  \rightsquigarrow
  \cdots
  \rightsquigarrow
  W_n
\]

and

\[
  W_n\eqworld^{\mathcal K}W_0.
\]

Even then,

\[
  \traj{W_n}
  \neq
  \traj{W_0}
\]

in the relevant historical sense.

If the loop carries nontrivial holonomy as in
\cref{app:loop-closure}, the return can be world-equivalent while historically
nontrivial.

Thus a restored world is not an untouched world.

A reconstructed world is not an unbroken world.

A returned world is not a world that never departed.

The endpoint relation cannot erase the loop that produced it.

% -----------------------------------------------------------------------------
%  D.17  WORLDHOOD AND RSVP
% -----------------------------------------------------------------------------

\subsection{An RSVP Reading}
\label{app:rsvp-worldhood}

The RSVP framework offers one possible continuous realization of these ideas.

Let a candidate world state be represented schematically by

\[
  \mathfrak R_t
  =
  \bigl(
    \Phi_t,
    \mathbf v_t,
    S_t
  \bigr),
\]

where \(\Phi_t\) is a scalar field, \(\mathbf v_t\) a vector field, and \(S_t\)
an entropy-like field.

The present essay does not require a specific RSVP dynamics.

It requires only that world-relevant structure be encoded relationally across
the fields and their evolution.

A trajectory is then

\[
  \Gamma_{\mathrm{RSVP}}
  =
  \{
    \mathfrak R_t
  \}_{t\in I}.
\]

World-continuation cannot be reduced to equality

\[
  \mathfrak R_{t_0}
  =
  \mathfrak R_{t_1}.
\]

Nor need it require such equality.

Instead, one asks whether the field evolution remains within an admissible
class

\[
  \Gamma_{\mathrm{RSVP}}
  \in
  \mathcal A_W
\]

and whether the relations constitutive of the relevant world signature remain
realizable along that history.

This places RSVP on the trajectory side of the state/history distinction.

% -----------------------------------------------------------------------------
%  D.18  IRREVERSIBILITY
% -----------------------------------------------------------------------------

\subsection{World-Continuation Is Not Temporal Inversion}
\label{app:world-irreversibility}

Suppose an admissible transformation carries

\[
  W_0\to W_1.
\]

Even if another admissible transformation later produces

\[
  W_1\to W_0'
\]

with

\[
  W_0'\eqworld^{\mathcal K}W_0,
\]

it does not follow that the second transformation inverted the first.

The composite history

\[
  W_0
  \to
  W_1
  \to
  W_0'
\]

contains events absent from the zero-length history at \(W_0\).

If provenance is retained, then

\[
  [W_0]
  \neq
  [W_0']_{\mathrm{prov}}.
\]

Irreversibility here need not mean that no state can recur.

It means that recurrence of state does not annihilate occurrence of history.

This is the minimal irreversibility thesis required by the essay.

% -----------------------------------------------------------------------------
%  D.19  THREE INDEPENDENT QUESTIONS
% -----------------------------------------------------------------------------

\subsection{The Separation Theorem}
\label{app:separation-theorem}

The distinctions developed above can now be stated formally.

The three persistence predicates used below are instances of the same
trajectory-relative construction at different scales. Let \(S\in\{O,C,W\}\)
denote respectively the object, civilizational, and world scales. For each
\(S\), let \(X_S\) be the corresponding state space, let \(\Gamma(X_S)\) be
its history space, and let

\[
  \mathcal A_S=\mathcal A_{\mathcal C_S}\subseteq\Gamma(X_S)
\]

be the histories admissible relative to a specified scale-appropriate
constraint family \(\mathcal C_S\), in the sense of
\cref{rem:admissibility-relativity}. Define

\[
  \Gamma_{\mathcal A_S}(a,b)
  =
  \left\{
    \gamma\in\mathcal A_S:
    \Orig(\gamma)=a,\;
    \End(\gamma)=b
  \right\},
\]

and

\[
  \Cont_S(a,b)=1
  \quad\Longleftrightarrow\quad
  \Gamma_{\mathcal A_S}(a,b)\neq\varnothing.
\]

Thus \(\Cont_O\), \(\Cont_C\), and \(\Cont_W\) share a common formal shape
without thereby sharing a common admissibility criterion: \(\Cont_W\) is
exactly the \(S=W\) instance already defined in
\cref{app:world-path-space}, and \(\Cont_O\), \(\Cont_C\) are its \(S=O\),
\(S=C\) counterparts. A \emph{bridge principle} between two persistence
scales is additional structure relating their state spaces, histories, or
admissibility constraints. No such bridge is contained merely in the fact
that continuation holds at one scale.

\begin{theorem}[Separation of persistence scales]
\label{thm:persistence-separation}
For \(S,S'\in\{O,C,W\}\) with \(S\neq S'\), no bridge principle between
scale \(S\) and scale \(S'\) is derivable from the schema above alone.
Consequently, without additional assumptions relating \(\mathcal A_O\),
\(\mathcal A_C\), and \(\mathcal A_W\), none of \(\Cont_O\), \(\Cont_C\),
\(\Cont_W\) determines either of the other two.
\end{theorem}

\begin{proof}
It is enough to provide separating cases.

An object can persist while its civilization and world collapse, so
\(\Cont_O\) does not determine \(\Cont_C\) or \(\Cont_W\).

A civilization can preserve its organizational continuity while migrating
away from the ecological and relational structure constitutive of its prior
world, so \(\Cont_C\) does not determine \(\Cont_W\).

A world can preserve its relevant relational organization while particular
objects and civilizations disappear and are replaced, so \(\Cont_W\) does not
determine \(\Cont_O\) or \(\Cont_C\).

In each case the schema above supplies the reason such a case is even
statable: \(X_O\), \(X_C\), and \(X_W\) are distinct state spaces with
distinct admissible families \(\mathcal A_O\), \(\mathcal A_C\),
\(\mathcal A_W\), so a nonempty \(\Gamma_{\mathcal A_S}(a,b)\) at one scale
carries no formal content about \(\Gamma_{\mathcal A_{S'}}(a',b')\) at
another unless a bridge principle explicitly relates \(\mathcal A_S\) to
\(\mathcal A_{S'}\). Therefore no implication follows without additional
assumptions connecting the respective admissibility criteria.
\end{proof}

The theorem is deliberately modest.

It does not claim that the scales are unrelated.

It says that their relation must be supplied rather than smuggled into the
word \enquote{continuity}.

% -----------------------------------------------------------------------------
%  D.20  THE NO-WORLD CONSEQUENCE
% -----------------------------------------------------------------------------

\subsection{Persistence Without a Primitive Container}
\label{app:no-world-consequence}

If the preceding analysis is accepted, then the familiar question

\[
  \text{Does }x\text{ remain in the same world?}
\]

decomposes into at least two questions.

First,

\[
  \Cont_O(x_0,x_1)?
\]

Second,

\[
  \Cont_W(W_0,W_1)?
\]

The second cannot be answered merely by pointing to a container in which both
states are located, because the persistence of that container is precisely
what is at issue.

Worldhood must itself possess a trajectory.

Thus

\[
  W_0
  \leadsto
  W_1
\]

is not background notation.

It is part of the object of inquiry.

% -----------------------------------------------------------------------------
%  D.21  THE WORLD-ORIGINIST PRINCIPLE
% -----------------------------------------------------------------------------

\subsection{The Result Needed by the Essay}
\label{app:world-originist-principle}

The technical argument of this appendix can be compressed into the following
statement.

\begin{theorem}[World-Originist principle]
\label{thm:world-originist}
Suppose worldhood is constituted by relational and trajectory-dependent
structure. Then equality or equivalence of terminal world states is
insufficient, in general, to determine world provenance. If, in addition,
world-continuation is evaluated by an admissibility criterion that
distinguishes histories with equivalent terminal states, then terminal
world equivalence is likewise insufficient to determine world-continuation.
\end{theorem}

\begin{proof}
Let

\[
  \End_W:
  \Gamma(\mathfrak W)
  \to
  \mathfrak W
\]

be the world-level endpoint map.

Suppose two distinct world histories satisfy

\[
  \gamma_A\neq\gamma_B
\]

while

\[
  \End_W(\gamma_A)
  \eqworld^{\mathcal K}
  \End_W(\gamma_B).
\]

Any world-provenance observable that differs on \(\gamma_A\) and \(\gamma_B\)
is not determined by the equivalent endpoints. By
\cref{thm:trajectory-necessity}, such an observable cannot factor through the
endpoint representation alone.

If the admissibility of continuation also depends upon the histories, then
world-continuation likewise cannot be inferred from terminal equivalence
alone.
\end{proof}

\begin{held}[No-World Consequence]
\centering
A world is not what remains after history is discarded.\\[3pt]
If worldhood persists, its persistence has a history too.
\end{held}

This completes the mathematical ascent begun with the smallest Originist
example.

At the artifact level,

\[
  x_A=x_B
\]

did not imply equality of histories.

At the civilizational level, material replacement did not settle
continuation.

At the world level, even world-equivalent endpoints fail to determine whether
one world has continued into another.

The scale has changed.

The structural distinction has not.

% ============================================================================
%  APPENDIX E
%  IRREVERSIBILITY, BRANCHING, AND THE PATH CATEGORY
% ============================================================================

\section{Irreversibility, Branching, and the Path Category}
\label{app:path-category}

\movementquestion{What changes when histories are treated as arrows rather than when states alone are treated as the objects of metaphysical comparison?}

\begin{orientation}
The preceding appendices have repeatedly used the same structure without yet
naming it categorically. There are states. There are transformations between
states. Transformations can be composed when the endpoint of one is the origin
of another. Different histories may share the same endpoints. Some histories
can be reversed; others cannot. Branches can depart from a common origin and
later converge.

This is naturally represented by a category of paths.

The categorical language is useful here for a specific reason. It prevents
several distinctions from being collapsed into the notation for endpoint
identity. A path from \(x\) to \(y\) is not the equation \(x=y\). Two paths
with the same source and target need not be the same path. The existence of a
path from \(x\) to \(y\) does not imply the existence of a path from \(y\) to
\(x\). Even when both exist, they need not be inverses.

Those elementary facts provide a compact formal home for the essay's
irreversibility thesis. They also expose the danger of passing too quickly to
a groupoid, where every arrow is invertible. Groupoid completion can be
mathematically useful, but an Originist must ask what historical distinction
is being changed when an irreversible path is supplied with a formal inverse.

The appendix ends by returning to the three literary probes. Card asks about
arrows whose sources are partially hidden. Zebrowski asks which outgoing
arrows count as continuation. Le Guin asks what it means for a path to return
to an object or its equivalence class without thereby becoming the identity
arrow.
\end{orientation}

% -----------------------------------------------------------------------------
%  E.1  THE PATH CATEGORY
% -----------------------------------------------------------------------------

\subsection{States as Objects and Histories as Morphisms}
\label{app:path-category-definition}

As in \cref{rem:admissibility-relativity}, admissibility here is relative to
a specified constraint family \(\mathcal C\). Thus the category constructed
below is, more explicitly, \(\mathbf{Path}_{\mathcal A_{\mathcal C}}(X)\).
When a single constraint family is fixed by context, we suppress this
dependence and write \(\mathbf{Path}_{\mathcal A}(X)\). Accordingly, every
result stated below in terms of \(\mathbf{Path}_{\mathcal A}(X)\) —
including faithfulness and cancellation — is conditional on the
admissibility structure induced by that fixed \(\mathcal C\). (The
Separation Theorem of \cref{app:separation-theorem} is not among these: its
continuation predicates \(\Cont_O\), \(\Cont_C\), \(\Cont_W\) are introduced
independently of \(\mathcal A\) and are not currently defined in terms of
it.) In particular,

\[
  \mathcal C_1\neq\mathcal C_2
  \quad\Longrightarrow\quad
  \mathbf{Path}_{\mathcal A_{\mathcal C_1}}(X)
  \text{ need not equal }
  \mathbf{Path}_{\mathcal A_{\mathcal C_2}}(X).
\]

Let \(X\) be a state space and let \(\mathcal A\) specify the admissible
elementary transformations between states.

Construct a directed graph whose vertices are states

\[
  x\in X
\]

and whose directed edges are admissible elementary transformations

\[
  x\xrightarrow{T}y.
\]

The free path category generated by this graph will be denoted

\[
  \mathbf{Path}_{\mathcal A}(X).
\]

Its objects are the states of \(X\).

Its morphisms

\[
  \Hom_{\mathbf{Path}_{\mathcal A}(X)}(x,y)
\]

are finite admissible paths from \(x\) to \(y\).

A morphism may therefore be written

\[
  \gamma:
  x
  \longrightarrow
  y,
\]

where

\[
  \gamma
  =
  \left(
  x=x_0
  \xrightarrow{T_1}
  x_1
  \xrightarrow{T_2}
  \cdots
  \xrightarrow{T_n}
  x_n=y
  \right).
\]

Composition is concatenation.

If

\[
  \gamma:x\to y
\]

and

\[
  \eta:y\to z,
\]

then

\[
  \eta\circ\gamma:x\to z
\]

is the history obtained by traversing \(\gamma\) and then \(\eta\).

For each state \(x\), the identity morphism

\[
  1_x:x\to x
\]

is the zero-length path at \(x\).

% -----------------------------------------------------------------------------
%  E.2  ENDPOINT EQUALITY IS NOT MORPHISM EQUALITY
% -----------------------------------------------------------------------------

\subsection{Parallel Histories}
\label{app:parallel-histories}

The central Originist distinction appears categorically as the possibility of
parallel morphisms.

Suppose

\[
  \gamma_A,\gamma_B:x\to y.
\]

Both histories have the same origin and endpoint.

Nevertheless,

\[
  \gamma_A\neq\gamma_B
\]

may hold.

Thus

\[
  \Hom(x,y)
\]

need not contain only one element.

A state-only representation effectively forgets which member of

\[
  \Hom(x,y)
\]

was realized.

\begin{proposition}[Endpoint coincidence does not determine a morphism]
\label{prop:endpoint-not-morphism}
If

\[
  |\Hom(x,y)|>1,
\]

then knowledge of the ordered pair

\[
  (x,y)
\]

is insufficient to determine the realized history.
\end{proposition}

\begin{proof}
Let

\[
  \gamma_1,\gamma_2\in\Hom(x,y)
\]

with

\[
  \gamma_1\neq\gamma_2.
\]

Both correspond to the same source-target pair. Therefore the pair \((x,y)\)
cannot uniquely select which morphism was realized.
\end{proof}

This is the categorical form of

\[
  x_A=x_B
  \qquad\text{while}\qquad
  \gamma_A\neq\gamma_B.
\]

The endpoint state may converge.

The morphisms remain distinct.

% -----------------------------------------------------------------------------
%  E.3  THE IDENTITY MORPHISM
% -----------------------------------------------------------------------------

\subsection{Remaining Is Not Returning}
\label{app:identity-vs-return}

The distinction between the identity morphism and a nontrivial loop is
especially important.

For a state \(x\),

\[
  1_x:x\to x
\]

contains no intervening transformation.

Now suppose

\[
  \gamma:x\to x
\]

is a nontrivial loop.

Although

\[
  \Orig(\gamma)=\End(\gamma)=x,
\]

it does not follow that

\[
  \gamma=1_x.
\]

Indeed, in the free path category,

\[
  \gamma=1_x
\]

only if \(\gamma\) is the empty path.

\begin{held}[Categorical Form of the Return Problem]
\centering
To return to the same object is not to have taken the identity morphism.
\end{held}

This is the categorical core of recurrence without restoration.

% -----------------------------------------------------------------------------
%  E.4  IRREVERSIBILITY
% -----------------------------------------------------------------------------

\subsection{Invertibility Is an Additional Property}
\label{app:path-invertibility}

Let

\[
  \gamma:x\to y.
\]

The existence of \(\gamma\) does not imply the existence of any morphism

\[
  \eta:y\to x.
\]

Even when such an \(\eta\) exists, it does not follow that

\[
  \eta\circ\gamma=1_x
\]

and

\[
  \gamma\circ\eta=1_y.
\]

Only when both equations hold is \(\gamma\) an isomorphism with inverse
\(\eta\).

Thus causal reachability

\[
  x\leadcausal y
\]

is weaker than categorical invertibility.

\begin{definition}[Historical irreversibility]
\label{def:historical-irreversibility}
An admissible history

\[
  \gamma:x\to y
\]

is \emph{historically irreversible} in a path category when it is not an
isomorphism.
\end{definition}

This definition is intentionally structural.

It does not require a thermodynamic claim.

A process may be historically irreversible simply because the category of
admissible transformations contains no genuine inverse for it.

% -----------------------------------------------------------------------------
%  E.5  REVERSE REACHABILITY IS NOT INVERSION
% -----------------------------------------------------------------------------

\subsection{Going Back Is Not Undoing}
\label{app:going-back}

Suppose

\[
  \gamma:x\to y
\]

and

\[
  \eta:y\to x.
\]

The existence of both directions establishes mutual reachability.

It does not establish

\[
  \eta=\gamma^{-1}.
\]

Their composite may be a nontrivial loop:

\[
  \eta\circ\gamma:x\to x,
\]

with

\[
  \eta\circ\gamma\neq1_x.
\]

Similarly,

\[
  \gamma\circ\eta\neq1_y
\]

may hold.

This provides a categorical version of the distinction already developed
through recurrence and holonomy.

One may return to an equivalent state without undoing the history by which one
left it.

% -----------------------------------------------------------------------------
%  E.6  A PROVENANCE FUNCTOR
% -----------------------------------------------------------------------------

\subsection{Recording Histories Functorially}
\label{app:provenance-functor}

Let \(\mathbf P\) be a category whose objects encode provenance states and
whose morphisms encode provenance transformations.

A provenance assignment may be represented by a functor

\[
  \mathsf{Prov}:
  \mathbf{Path}_{\mathcal A}(X)
  \to
  \mathbf P.
\]

Functoriality requires

\[
  \mathsf{Prov}(1_x)
  =
  1_{\mathsf{Prov}(x)}
\]

and

\[
  \mathsf{Prov}(\eta\circ\gamma)
  =
  \mathsf{Prov}(\eta)
  \circ
  \mathsf{Prov}(\gamma).
\]

The construction makes composition history-sensitive.

For example, if provenance morphisms record transformation words, then

\[
  \mathsf{Prov}(\gamma)
  =
  R\circ C\circ R
\]

may differ from

\[
  \mathsf{Prov}(\gamma')
  =
  C\circ R\circ R
\]

even when

\[
  \Orig(\gamma)=\Orig(\gamma')
\]

and

\[
  \End(\gamma)=\End(\gamma').
\]

Generation count can then be obtained by a further forgetting map that retains
only word length.

% -----------------------------------------------------------------------------
%  E.7  FORGETFUL FUNCTORS
% -----------------------------------------------------------------------------

\subsection{Forgetting the Arrow}
\label{app:forgetful-arrow}

A state-oriented system can be modeled by a representation that forgets some
or all morphism structure.

At the crudest level, consider the endpoint operation

\[
  \End:
  \Mor(\mathbf{Path}_{\mathcal A}(X))
  \to
  \Ob(\mathbf{Path}_{\mathcal A}(X)).
\]

This is not itself ordinarily a functor between the same categories, because
it maps morphisms to objects, but it captures the relevant forgetting
operation.

More structured quotients may instead be represented by functors

\[
  F:
  \mathbf{Path}_{\mathcal A}(X)
  \to
  \mathbf C
\]

that identify multiple morphisms.

If

\[
  F(\gamma_1)=F(\gamma_2)
\]

while

\[
  \mathsf{Prov}(\gamma_1)
  \neq
  \mathsf{Prov}(\gamma_2),
\]

then \(F\) is not faithful with respect to the provenance distinction.

This suggests a categorical reformulation of Originist loss.

% -----------------------------------------------------------------------------
%  E.8  FAITHFULNESS AND ORIGINIST LOSS
% -----------------------------------------------------------------------------

\subsection{Faithfulness}
\label{app:faithfulness}

Recall that a functor

\[
  F:\mathbf C\to\mathbf D
\]

is faithful when, for every pair of objects \(x,y\), the induced map

\[
  F_{x,y}:
  \Hom_{\mathbf C}(x,y)
  \to
  \Hom_{\mathbf D}(F(x),F(y))
\]

is injective.

If \(F\) is not faithful, then there exist distinct parallel morphisms

\[
  \gamma_1\neq\gamma_2
\]

such that

\[
  F(\gamma_1)=F(\gamma_2).
\]

This is almost exactly the Originist failure mode.

\begin{definition}[Originist-faithful representation]
\label{def:originist-faithful}
A functorial representation \(F\) is \emph{Originist-faithful relative to a
provenance distinction} when it does not identify morphisms that differ under
that provenance distinction.
\end{definition}

Full categorical faithfulness may preserve more information than a task
requires.

Originist faithfulness is therefore relative to the distinctions whose loss
would affect the intended inference.

% -----------------------------------------------------------------------------
%  E.9  BRANCHING
% -----------------------------------------------------------------------------

\subsection{Branching as Multiple Outgoing Morphisms}
\label{app:categorical-branching}

Suppose

\[
  \alpha:x\to y
\]

and

\[
  \beta:x\to z
\]

are admissible histories with

\[
  y\neq z.
\]

The common source \(x\) has branched into distinct descendants.

Nothing in category theory requires one branch to be more authentic than the
other.

Both morphisms exist.

If both satisfy the relevant admissibility constraints, then both are
continuations in the weaker relational sense.

The question of numerical identity cannot be settled by the mere existence of
these arrows, as established in \cref{prop:branching-obstruction}.

Categorically, the structure is simply

\[
\begin{tikzcd}[column sep=4em,row sep=2em]
& x \arrow[dl,"\alpha"'] \arrow[dr,"\beta"] & \\
y && z.
\end{tikzcd}
\]

The fork is not a defect in the model.

It is information the model must preserve.

% -----------------------------------------------------------------------------
%  E.10  CONVERGENCE
% -----------------------------------------------------------------------------

\subsection{Convergence as Multiple Incoming Histories}
\label{app:categorical-convergence}

Dually, suppose

\[
  \alpha:x\to z
\]

and

\[
  \beta:y\to z
\]

with

\[
  x\neq y.
\]

Then distinct origins converge upon the same endpoint.

\[
\begin{tikzcd}[column sep=4em,row sep=2em]
x \arrow[dr,"\alpha"'] && y \arrow[dl,"\beta"] \\
& z &
\end{tikzcd}
\]

This is the categorical shape of independent rediscovery and convergent
reconstruction.

The terminal object \(z\), considered only as a state, does not tell us
whether it was reached through \(\alpha\), through \(\beta\), or through some
other morphism.

The convergence problem is therefore not exceptional.

It is built into any category in which distinct arrows may share a codomain.

% -----------------------------------------------------------------------------
%  E.11  COMMUTING SQUARES AND DISTINCT HISTORIES
% -----------------------------------------------------------------------------

\subsection{Commutativity Can Erase a Distinction}
\label{app:commuting-histories}

Suppose

\[
\begin{tikzcd}[column sep=4em,row sep=3em]
O \arrow[r,"a"] \arrow[d,"b"'] &
A \arrow[d,"c"] \\
B \arrow[r,"d"'] &
X
\end{tikzcd}
\]

commutes.

Then

\[
  c\circ a
  =
  d\circ b
\]

as morphisms in the category.

This equality should not automatically be interpreted as historical identity.

Whether the two composites are genuinely to be identified depends upon what
relations were imposed when constructing the category.

In a free path category, the two paths remain distinct unless an explicit
relation equates them.

In a quotient category, they may become equal.

Thus the statement that a diagram commutes can itself encode a decision to
forget path distinctions.

The Originist question is therefore prior:

\[
  \text{Why are these composites being identified?}
\]

% -----------------------------------------------------------------------------
%  E.12  PRESENTATIONS OF HISTORY
% -----------------------------------------------------------------------------

\subsection{Relations Are Historical Commitments}
\label{app:relations-as-commitments}

A category may be presented by generators and relations.

Let \(G\) be a directed graph of elementary transformations and let

\[
  \mathbf{Free}(G)
\]

be its free category.

Now impose a collection of relations

\[
  \mathcal R
\]

among paths.

The resulting category is schematically

\[
  \mathbf C
  =
  \mathbf{Free}(G)/{\mathcal R}.
\]

Every relation of the form

\[
  \gamma_1\sim_{\mathcal R}\gamma_2
\]

declares two histories equivalent for the purposes of \(\mathbf C\).

This is mathematically ordinary.

From the Originist perspective, however, the relation set \(\mathcal R\)
becomes an epistemic object.

Each relation states which historical distinctions the formalism has chosen
not to preserve.

The quotient-loss operator of Appendix~\ref{app:originist-loss} can therefore
be applied to categorical presentations themselves.

% -----------------------------------------------------------------------------
%  E.13  THE PATH CATEGORY AND THE PATH GROUPOID
% -----------------------------------------------------------------------------

\subsection{Formal Reversibility}
\label{app:path-groupoid}

Given a category \(\mathbf C\), one may formally invert a selected collection
of morphisms.

In the extreme case, invert every morphism to obtain a groupoid completion

\[
  \mathbf C
  \longrightarrow
  \mathbf C[\Mor(\mathbf C)^{-1}].
\]

For the path category, denote this schematically by

\[
  \mathbf{Path}_{\mathcal A}(X)
  \longrightarrow
  \mathbf{Gpd}_{\mathcal A}(X).
\]

Every path

\[
  \gamma:x\to y
\]

then acquires a formal inverse

\[
  \gamma^{-1}:y\to x
\]

satisfying

\[
  \gamma^{-1}\circ\gamma=1_x
\]

and

\[
  \gamma\circ\gamma^{-1}=1_y.
\]

This construction is mathematically legitimate.

Its ontological interpretation requires caution.

% -----------------------------------------------------------------------------
%  E.14  FORMAL INVERSION IS NOT PHYSICAL REVERSAL
% -----------------------------------------------------------------------------

\subsection{The Groupoid Warning}
\label{app:groupoid-warning}

Suppose the original path category contains an irreversible history

\[
  \gamma:x\to y
\]

with no inverse.

The groupoid completion introduces

\[
  \gamma^{-1}
\]

formally.

It does not thereby establish that an admissible physical, historical, or
computational process exists that literally undoes \(\gamma\).

The inverse belongs to the completed formal structure.

It need not belong to the original process ontology.

\begin{held}[Groupoid Warning]
\centering
Formal invertibility is a property of a chosen completion.\\[3pt]
It is not evidence that history itself ran backward.
\end{held}

This distinction matters whenever groupoid language is used in the No-World
framework.

A groupoid can represent equivalences among world presentations without
implying that every realized transformation is historically reversible.

% -----------------------------------------------------------------------------
%  E.15  WHAT GROUPOID COMPLETION FORGETS
% -----------------------------------------------------------------------------

\subsection{Irreversibility as Lost Structure}
\label{app:groupoid-loss}

Consider a category \(\mathbf C\) containing

\[
  \gamma:x\to y
\]

without an inverse.

After groupoid completion, \(\gamma\) becomes invertible.

The completed category therefore no longer distinguishes between

\[
  \text{``an inverse existed in the original process category''}
\]

and

\[
  \text{``an inverse was supplied by localization.''}
\]

Unless that provenance is recorded separately, the completion erases a
historical distinction about the status of the inverse.

This is itself an Originist phenomenon.

Two arrows can be extensionally identical inside the completed groupoid while
possessing different derivational status.

One was realized.

The other was formally adjoined.

% -----------------------------------------------------------------------------
%  E.16  A TAGGED COMPLETION
% -----------------------------------------------------------------------------

\subsection{Preserving the Provenance of Formal Inverses}
\label{app:tagged-completion}

An Originist-aware formalism can preserve this distinction by enriching the
completed arrows with provenance tags.

Write

\[
  \gamma^{-1}_{[\mathrm{formal}]}
\]

for an inverse introduced by localization and

\[
  \eta_{[0]}
\]

for a reverse transformation independently present in the original category.

Even if the completed category identifies

\[
  \eta
  =
  \gamma^{-1},
\]

their provenance annotations may remain distinct:

\[
  \Prov(\eta)
  \neq
  \Prov(\gamma^{-1}_{[\mathrm{formal}]}).
\]

This is another instance of the general principle

\[
  \text{semantic or structural equality}
  \centernot\Rightarrow
  \text{provenance equality}.
\]

The provenance layer does not invalidate the quotient.

It records what the quotient would otherwise conceal.

% -----------------------------------------------------------------------------
%  E.17  CANCELLATION AND IRREVERSIBILITY
% -----------------------------------------------------------------------------

\subsection{Cancellation Is Weaker Than Invertibility}
\label{app:cancellation}

Some categories permit cancellation without general invertibility.

For example, a morphism \(\gamma:x\to y\) may be left-cancellative when

\[
  \gamma\circ\alpha
  =
  \gamma\circ\beta
\]

implies

\[
  \alpha=\beta.
\]

Similarly, it may be right-cancellative.

These properties preserve some historical distinguishability.

They do not imply that \(\gamma\) possesses an inverse.

Thus a process can retain enough structure to distinguish prior paths while
remaining irreversible.

This intermediate case is useful for provenance systems.

One need not choose between total reversibility and total loss.

% -----------------------------------------------------------------------------
%  E.18  INFORMATION-LOSING MORPHISMS
% -----------------------------------------------------------------------------

\subsection{Non-Cancellative Transformations}
\label{app:noncancellative}

Suppose a transformation

\[
  C:x\to y
\]

acts as a lossy compression.

If two distinct preceding histories

\[
  \alpha,\beta:o\to x
\]

satisfy

\[
  C\circ\alpha
  =
  C\circ\beta,
\]

then \(C\) is not left-cancellative with respect to those histories.

The compression has merged distinct incoming derivations.

After the merge, observing the composite does not determine whether the prior
history was \(\alpha\) or \(\beta\).

This categorical pattern is exactly the one formalized by Originist loss:

\[
  \alpha\neq\beta,
  \qquad
  C\circ\alpha=C\circ\beta.
\]

A later repair morphism

\[
  R:y\to z
\]

cannot restore left-cancellativity merely by composition, because

\[
  R\circ C\circ\alpha
  =
  R\circ C\circ\beta.
\]

This recovers \cref{thm:perfect-repair} categorically.

% -----------------------------------------------------------------------------
%  E.19  THE REPAIR COROLLARY
% -----------------------------------------------------------------------------

\subsection{Postcomposition Cannot Separate an Existing Merge}
\label{app:postcomposition-merge}

\begin{theorem}[Categorical repair corollary]
\label{thm:categorical-repair}
Let

\[
  \alpha,\beta:o\to x
\]

be distinct morphisms and let

\[
  C:x\to y
\]

satisfy

\[
  C\circ\alpha=C\circ\beta.
\]

Then for every morphism

\[
  R:y\to z,
\]

one has

\[
  R\circ C\circ\alpha
  =
  R\circ C\circ\beta.
\]
\end{theorem}

\begin{proof}
By assumption,

\[
  C\circ\alpha=C\circ\beta.
\]

Postcompose both sides with \(R\). Associativity gives

\[
  R\circ(C\circ\alpha)
  =
  R\circ(C\circ\beta),
\]

hence

\[
  R\circ C\circ\alpha
  =
  R\circ C\circ\beta.
\]
\end{proof}

The result is trivial algebraically.

Conceptually it is central.

Once two histories have been merged by a representation, deterministic
postprocessing of that representation cannot separate them again.

A distinction must enter from somewhere else.

% -----------------------------------------------------------------------------
%  E.20  SIDE INFORMATION AS A PRODUCT
% -----------------------------------------------------------------------------

\subsection{Restoring Distinguishability by Enrichment}
\label{app:categorical-side-information}

Suppose an archival channel supplies additional information

\[
  A:x\to p.
\]

Instead of representing the history only through \(C\), construct the paired
representation

\[
  \langle C,A\rangle:
  x\to y\times p.
\]

It may happen that

\[
  C\circ\alpha=C\circ\beta
\]

while

\[
  A\circ\alpha\neq A\circ\beta.
\]

Then

\[
  \langle C,A\rangle\circ\alpha
  \neq
  \langle C,A\rangle\circ\beta.
\]

The archive restores distinguishability not by reversing the compression but
by preserving another projection.

This gives a categorical interpretation of the Archivist/Originist split.

The Archivist maintains additional arrows.

The Originist knows when they are required.

% -----------------------------------------------------------------------------
%  E.21  CARD AS AN INVERSE PROBLEM ON MORPHISMS
% -----------------------------------------------------------------------------

\subsection{Hidden Sources}
\label{app:card-categorical}

Card's originism can now be represented as an inference problem over incoming
morphisms \cite{card1989}.

Given a present state \(x\), one seeks candidate pairs

\[
  (O,\gamma)
\]

such that

\[
  \gamma:O\to x.
\]

The inverse problem is therefore not merely

\[
  x\rightsquigarrow O?
\]

but

\[
  x
  \rightsquigarrow
  \{
    (O,\gamma)
    :
    \gamma\in\Hom(O,x)
  \}.
\]

Even if the origin object \(O\) were uniquely identified, the morphism
\(\gamma\) might remain underdetermined.

Thus origin recovery and history recovery are distinct inverse problems.

% -----------------------------------------------------------------------------
%  E.22  ZEBROWSKI AS AN OUTGOING-MORPHISM PROBLEM
% -----------------------------------------------------------------------------

\subsection{Admissible Futures}
\label{app:zebrowski-categorical}

The \emph{Macrolife} problem points in the opposite direction
\cite{zebrowski1979}.

Given an origin \(O\), consider the collection

\[
  \coprod_{X}
  \Hom(O,X).
\]

This is the family of represented outgoing histories from \(O\).

Continuation introduces an admissibility predicate

\[
  A:
  \Mor(\mathbf{Path}_{\mathcal A}(X))
  \to
  \{0,1\}.
\]

The relevant future space becomes

\[
  \{
    \gamma:O\to X
    :
    A(\gamma)=1
  \}.
\]

The problem is therefore not to identify one predetermined future.

It is to characterize the admissible cone of descent.

% -----------------------------------------------------------------------------
%  E.23  LE GUIN AS A LOOP PROBLEM
% -----------------------------------------------------------------------------

\subsection{Return and Endomorphisms}
\label{app:leguin-categorical}

Le Guin's recurrence problem is naturally expressed through endomorphisms
\cite{leguin1974}.

An endomorphism of \(O\) is a morphism

\[
  \gamma:O\to O.
\]

More generally, if return is only to an equivalent origin \(O'\), one has

\[
  \gamma:O\to O'
\]

with

\[
  O'\sim_{\mathcal C}O.
\]

The crucial distinction is between

\[
  1_O
\]

and a nontrivial historical loop

\[
  \gamma.
\]

The former says nothing happened.

The latter says departure and return occurred.

Endpoint recurrence identifies what the two share.

Provenance records what they do not.

% -----------------------------------------------------------------------------
%  E.24  THE THREE PROBES AS HOM-QUESTIONS
% -----------------------------------------------------------------------------

\subsection{One Category, Three Directions}
\label{app:three-hom-questions}

The three literary probes can now be expressed within one structure.

Card begins at \(X\) and asks which incoming histories are warranted:

\[
  \coprod_O\Hom(O,X).
\]

Zebrowski begins at \(O\) and asks which outgoing histories remain admissible:

\[
  \coprod_X\Hom(O,X).
\]

Le Guin asks about histories that return:

\[
  \End(O)
  =
  \Hom(O,O),
\]

or, more generally,

\[
  \coprod_{O'\sim_{\mathcal C}O}
  \Hom(O,O').
\]

The three questions are therefore neither identical nor accidentally similar.

They interrogate different regions of the same path structure.

% -----------------------------------------------------------------------------
%  E.25  THE CATEGORY OF ADMISSIBLE DESCENT
% -----------------------------------------------------------------------------

\subsection{Closure Under Composition}
\label{app:admissible-category}

For admissible histories to form a category, admissibility must satisfy at
least two closure conditions.

Identity histories must be admissible:

\[
  A(1_x)=1.
\]

And whenever

\[
  A(\gamma)=1
\]

and

\[
  A(\eta)=1
\]

for composable histories, one must have

\[
  A(\eta\circ\gamma)=1.
\]

If the second condition fails, admissible segments do not automatically
compose into admissible histories.

This reproduces the local-to-global warning of
Appendix~\ref{app:admissibility} in categorical language.

One should therefore distinguish a category of globally admissible histories
from a graph whose edges are merely locally admissible.

\begin{proposition}[Composition requires a global admissibility principle]
\label{prop:admissibility-composition}
Local admissibility of elementary transformations defines a category of
admissible histories only if admissibility is closed under finite
composition.
\end{proposition}

This condition should never be presumed merely because each step passes an
individual test.

% -----------------------------------------------------------------------------
%  E.26  HISTORY-DEPENDENT ADMISSIBILITY
% -----------------------------------------------------------------------------

\subsection{When Composition Changes the Judgment}
\label{app:history-dependent-admissibility}

Suppose

\[
  A(\gamma)=1
\]

and

\[
  A(\eta)=1
\]

but

\[
  A(\eta\circ\gamma)=0.
\]

Then admissibility depends upon more than the local morphisms considered in
isolation.

This can occur through accumulated drift, forbidden transformation sequences,
resource exhaustion, provenance depth, or global gluing obstruction.

The correct structure may then require enriched objects carrying sufficient
history to make composition local again.

For example, replace a state \(x\) by

\[
  (x,h),
\]

where \(h\) records a relevant historical summary.

A transformation acts on both:

\[
  (x,h)
  \longrightarrow
  (y,h').
\]

This maneuver does not eliminate history dependence.

It moves the required history into the state representation.

The price of recovering Markovian composition is therefore an enlarged state.

% -----------------------------------------------------------------------------
%  E.27  STATE AUGMENTATION
% -----------------------------------------------------------------------------

\subsection{When History Becomes State}
\label{app:history-becomes-state}

This observation produces an apparent challenge to the essay's thesis.

If every relevant historical quantity can be appended to the state, why insist
that continuity is a property of trajectories rather than states?

Let

\[
  \widehat X
  =
  X\times H
\]

be an augmented state space containing a complete history variable \(H\).

Then a sufficiently rich augmented state

\[
  \widehat x=(x,h)
\]

may indeed determine provenance.

But \(h\) is itself a representation of the trajectory.

The construction therefore does not show that history was unnecessary.

It shows that history can be encoded into a larger state.

\begin{objection}[The State-Augmentation Objection]
If a state is allowed to contain its complete derivational history, then every
trajectory-dependent property can apparently be rewritten as a state property.
Why, then, privilege trajectories?
\end{objection}

The answer is that the distinction in this essay concerns representational
content, not grammatical location.

A state containing its complete history is no longer a state in the thin sense
against which the trajectory argument was directed.

It is an endpoint plus a provenance archive.

The Originist has won by annexation.

% -----------------------------------------------------------------------------
%  E.28  THE MINIMAL-STATE QUESTION
% -----------------------------------------------------------------------------

\subsection{How Much History Must Be Carried Forward?}
\label{app:minimal-state}

The useful question is therefore not whether history can be encoded in state.

It is how much historical information must be retained for a specified future
task.

Let

\[
  F:\Gamma(X)\to Y
\]

be a task.

Seek a historical summary

\[
  h_F:\Gamma(X)\to H_F
\]

such that

\[
  F(\gamma)
\]

is determined by

\[
  \bigl(
    \End(\gamma),
    h_F(\gamma)
  \bigr).
\]

The pair

\[
  \widehat x_F
  =
  \bigl(
    \End(\gamma),
    h_F(\gamma)
  \bigr)
\]

is a task-sufficient augmented state.

This connects the path-category formulation directly to the minimal sufficient
provenance problem of \cref{app:minimal-sufficient-provenance}.

The practical Originist role is therefore not to preserve infinite history
indiscriminately.

It is to determine which historical summaries future judgments cannot safely
do without.

% -----------------------------------------------------------------------------
%  E.29  THE PATH-CATEGORY THESIS
% -----------------------------------------------------------------------------

\subsection{The Formal Result}
\label{app:path-category-thesis}

We can now state the categorical result needed by the essay.

\begin{theorem}[Path-category thesis]
\label{thm:path-category-thesis}
Let \(\mathbf C\) be a category of admissible histories containing distinct
parallel morphisms

\[
  \gamma_1,\gamma_2:x\to y.
\]

Any property \(P\) satisfying

\[
  P(\gamma_1)\neq P(\gamma_2)
\]

cannot be represented as a function solely of the source-target pair
\((x,y)\).
\end{theorem}

\begin{proof}
Assume there exists

\[
  \widetilde P:
  \Ob(\mathbf C)\times\Ob(\mathbf C)
  \to Z
\]

such that for every morphism

\[
  \gamma:a\to b
\]

one has

\[
  P(\gamma)=\widetilde P(a,b).
\]

Since \(\gamma_1\) and \(\gamma_2\) have the same source \(x\) and target \(y\),

\[
  P(\gamma_1)
  =
  \widetilde P(x,y)
  =
  P(\gamma_2),
\]

contradicting

\[
  P(\gamma_1)\neq P(\gamma_2).
\]

Therefore no such state-pair representation exists.
\end{proof}

The theorem is simply the trajectory necessity criterion translated into
categorical language.

Its advantage is that branching, convergence, composition, loops, and
irreversibility now belong to one formal object.

% -----------------------------------------------------------------------------
%  E.30  THE IRREVERSIBILITY COROLLARY
% -----------------------------------------------------------------------------

\subsection{History Survives Recurrence}
\label{app:irreversibility-corollary}

\begin{theorem}[Recurrence does not erase history]
\label{thm:recurrence-history}
Let

\[
  \gamma:x\to x
\]

be a nonidentity endomorphism in a path category. Then equality of source and
target does not imply

\[
  \gamma=1_x.
\]
\end{theorem}

\begin{proof}
By hypothesis, \(\gamma\) is a nonidentity endomorphism. Therefore

\[
  \gamma\neq1_x
\]

despite having the same source and target as \(1_x\).
\end{proof}

The proof is intentionally almost empty.

The philosophical mistake it blocks is not.

Endpoint restoration is not historical erasure.

% -----------------------------------------------------------------------------
%  E.31  THE CATEGORICAL ORIGINIST OBJECTION
% -----------------------------------------------------------------------------

\subsection{Against Premature Thinness}
\label{app:categorical-originist-objection}

The Originist objection can finally be expressed as a warning about functors
and quotient categories.

\begin{objection}[Categorical Originist Objection]
Let

\[
  F:\mathbf C\to\mathbf D
\]

be a representation of a history-bearing category. If \(F\) identifies
morphisms whose distinction is required by a future provenance,
admissibility, or continuation judgment, then \(F\) is too coarse for that
judgment even when it preserves every endpoint relevant to the present task.
\end{objection}

This objection does not oppose abstraction.

Category theory itself proceeds by finding the correct level of abstraction.

The objection concerns the word \enquote{correct}.

A quotient appropriate for one inference can be destructive for another.

The loss becomes visible only when histories are allowed to remain distinct
long enough to ask what the quotient has done to them.

% -----------------------------------------------------------------------------
%  E.32  FROM PATH CATEGORY TO DESCENT
% -----------------------------------------------------------------------------

\subsection{Descent as the Primitive Candidate}
\label{app:descent-primitive}

The essay began by asking whether the primitive object should be the state

\[
  x,
\]

the trajectory

\[
  \gamma,
\]

the admissibility class

\[
  [\gamma]_{\mathcal A},
\]

or the provenance graph

\[
  \graphG.
\]

The categorical analysis does not force a single answer.

It does, however, establish an ordering of expressive capacity.

A bare object records less than an object together with its incoming and
outgoing morphisms.

A single morphism records less than the branching category in which alternative
morphisms remain available.

A quotient category may recover tractability while discarding distinctions
among those morphisms.

A provenance enrichment may restore some of what the quotient forgets.

Thus there is no representation-free primitive waiting to be discovered by
notation alone.

The selection of primitive structure determines which histories can still be
distinguished.

That selection therefore has provenance consequences of its own.

% -----------------------------------------------------------------------------
%  E.33  CLOSURE OF THE FORMAL ARGUMENT
% -----------------------------------------------------------------------------

\subsection{What the Category Adds}
\label{app:categorical-closure}

The path-category formulation adds no new literary thesis to the essay.

It shows that the distinctions already extracted from Card, Zebrowski, and
Le Guin possess a common compositional structure.

Origin is an incoming-morphism problem.

Continuation is an admissible outgoing-morphism problem.

Return is an endomorphism problem.

Branching is multiplicity of outgoing morphisms.

Convergence is multiplicity of incoming morphisms.

Irreversibility is failure of general invertibility.

Compression can appear as non-cancellation.

Reconstruction is postcomposition and therefore cannot separate histories
already identified upstream.

Archival side information enriches the representation rather than reversing
the loss.

World-continuation lifts the same structure to a category whose objects are
candidate world realizations.

Nothing in this formalization requires the endpoint to disappear.

It requires only that the endpoint cease to monopolize the ontology.

\begin{held}[Path-Category Thesis]
\centering
The state tells us where an arrow ends.\\[3pt]
It does not tell us which arrow occurred.
\end{held}

This is enough to connect the formal apparatus back to the sentence around
which the essay has repeatedly circled:

\[
  \boxed{
    \text{continuity is a property of trajectories, not of states}
  }.
\]

Even this sentence now requires one qualification.

A trajectory may itself be too thin if branching alternatives, provenance
relations, admissibility constraints, or world-level gluing are relevant.

The primitive may therefore be not a path but a structured space of possible
paths.

That question is left open deliberately.

Closing it by fiat would reproduce exactly the operation this essay has spent
its length learning to inspect.

% ============================================================================
%  APPENDIX F
%  PROVENANCE TYPES, EPISTEMIC ANNOTATION, AND THE [omega?] WOUND
% ============================================================================

\section{Provenance Types, Epistemic Annotation, and the \texorpdfstring{$[\omega?]$}{[omega?]} Wound}
\label{app:provenance-types}

\movementquestion{What changes when provenance ceases to be metadata attached to a proposition and becomes part of the proposition's epistemic type?}

\begin{orientation}
A citation normally answers a question adjacent to a claim. It tells the reader
where to look for support, attribution, or prior discussion. The claim itself
is ordinarily treated as unchanged by the citation attached to it.

The apparatus used throughout this essay makes a stronger distinction.
Identical linguistic content can arrive through different histories. One copy
may have been inspected directly. Another may have been reconstructed from a
summary. A third may have passed through repeated compression and regeneration.
A fourth may have a derivational history that certainly exists but can no
longer be recovered.

These propositions may say exactly the same thing.

They need not occupy the same epistemic position.

The notation \( [0] \), \( [1] \), \( [2] \), and the longer transformation
words used in the main text are therefore not intended as decorative citation
styles. They encode properties of a derivational trajectory. The exceptional
tag \( [\omega?] \) marks a different failure: not a known large generation
count, but the loss of our ability to establish what the generation count or
transformation history was.

That distinction permits the essay's typographic conceit to be stated
formally. A proposition has content. Its occurrence has provenance. Our
knowledge of that provenance is a third object. Truth, derivation, and
provenance knowledge can therefore change independently.
\end{orientation}

% -----------------------------------------------------------------------------
%  F.1  PROPOSITIONAL CONTENT AND OCCURRENCE
% -----------------------------------------------------------------------------

\subsection{Content Is Not Occurrence}
\label{app:content-occurrence}

Let

\[
  \Phi
\]

be a space of propositional contents.

An element

\[
  \varphi\in\Phi
\]

represents what is asserted.

Now let

\[
  \mathcal O(\varphi)
\]

denote the set of concrete occurrences or realizations of that content.

Two occurrences

\[
  o_1,o_2\in\mathcal O(\varphi)
\]

may therefore satisfy

\[
  \Content(o_1)
  =
  \Content(o_2)
  =
  \varphi
\]

while

\[
  o_1\neq o_2.
\]

The distinction is elementary but indispensable.

Provenance belongs first to occurrences and histories, not to abstract
propositional content considered without realization.

% -----------------------------------------------------------------------------
%  F.2  TYPED CLAIMS
% -----------------------------------------------------------------------------

\subsection{The Provenance-Typed Proposition}
\label{app:typed-proposition}

Let \(D\) be a space of provenance descriptions.

Define a provenance-typed claim as a pair

\[
  \widehat{\varphi}
  =
  (\varphi,\pi),
\]

where

\[
  \varphi\in\Phi
\]

and

\[
  \pi\in D.
\]

The projection

\[
  \Content:
  \Phi\times D
  \to
  \Phi
\]

is given by

\[
  \Content(\varphi,\pi)=\varphi.
\]

The provenance projection is

\[
  \Prov:
  \Phi\times D
  \to
  D
\]

with

\[
  \Prov(\varphi,\pi)=\pi.
\]

Consequently,

\[
  (\varphi,\pi_1)
  \neq
  (\varphi,\pi_2)
\]

whenever

\[
  \pi_1\neq\pi_2,
\]

even though

\[
  \Content(\varphi,\pi_1)
  =
  \Content(\varphi,\pi_2).
\]

The essay's notation

\[
  \varphi^{[0]},
  \qquad
  \varphi^{[1]},
  \qquad
  \varphi^{[3:R\circ C\circ R]}
\]

is shorthand for such typed occurrences.

% -----------------------------------------------------------------------------
%  F.3  GENERATION DEPTH
% -----------------------------------------------------------------------------

\subsection{Generation as a Provenance Observable}
\label{app:generation-depth}

Let

\[
  \gamma
  =
  (
    o_0
    \xrightarrow{T_1}
    o_1
    \xrightarrow{T_2}
    \cdots
    \xrightarrow{T_n}
    o_n
  )
\]

be a derivational history.

Define the generation-depth observable

\[
  g(\gamma)=n.
\]

A directly inspected source occurrence has

\[
  g(\gamma)=0
\]

relative to the selected source boundary.

A first-generation reconstruction has

\[
  g(\gamma)=1.
\]

More generally,

\[
  \varphi^{[n]}
\]

asserts that the occurrence carrying \(\varphi\) lies at known derivational
depth \(n\).

Generation depth is therefore not a measure of semantic error.

One may have

\[
  \Content(o_0)
  =
  \Content(o_n)
\]

while

\[
  g(o_0)\neq g(o_n).
\]

This is the smallest possible formal example of

\[
  \text{state}\neq\text{history}.
\]

% -----------------------------------------------------------------------------
%  F.4  TRANSFORMATION WORDS
% -----------------------------------------------------------------------------

\subsection{Depth Is Too Coarse}
\label{app:transformation-words}

Generation count forgets which transformations occurred.

Let

\[
  \mathcal T
  =
  \{
    R,C,S,E,\ldots
  \}
\]

be an alphabet of transformation classes.

A history determines a transformation word

\[
  w(\gamma)
  =
  T_n\circ\cdots\circ T_2\circ T_1.
\]

Two histories can satisfy

\[
  g(\gamma_1)
  =
  g(\gamma_2)
\]

while

\[
  w(\gamma_1)
  \neq
  w(\gamma_2).
\]

For example,

\[
  w(\gamma_1)
  =
  R\circ C\circ R
\]

and

\[
  w(\gamma_2)
  =
  C\circ R\circ R
\]

both have depth three.

Thus

\[
  [3]
\]

is insufficient whenever transformation order matters.

The fuller annotation

\[
  [3:R\circ C\circ R]
\]

retains strictly more provenance information.

% -----------------------------------------------------------------------------
%  F.5  PROVENANCE RECORDS
% -----------------------------------------------------------------------------

\subsection{A Composite Provenance Type}
\label{app:composite-provenance}

For the purposes of this essay, define a provenance record schematically by

\[
  \pi
  =
  \left(
    g,
    w,
    \iota,
    p
  \right),
\]

where \(g\) is generation depth, \(w\) a transformation word, \(\iota\) an
invariant-preservation descriptor, and \(p\) an epistemic confidence assigned
to the provenance reconstruction.

Thus the notation

\[
  \varphi^{[3:R\circ C\circ R;\ \mathcal I^+;\ p=.71]}
\]

abbreviates a typed proposition whose provenance record contains

\[
  g=3,
\]

\[
  w=R\circ C\circ R,
\]

\[
  \iota=\mathcal I^+,
\]

and

\[
  p=.71.
\]

The confidence coordinate refers to confidence in the provenance
characterization unless explicitly stated otherwise.

It is not automatically confidence in the truth of \(\varphi\).

% -----------------------------------------------------------------------------
%  F.6  TRUTH AND PROVENANCE
% -----------------------------------------------------------------------------

\subsection{Orthogonal Questions}
\label{app:truth-provenance}

Let

\[
  \Truth:\Phi\to\{0,1\}
\]

be a truth valuation in a setting where bivalent notation is appropriate.

Let

\[
  \Prov:\widehat\Phi\to D
\]

be the provenance projection on typed propositions.

Then

\[
  \Truth(\varphi)
\]

and

\[
  \Prov(\widehat\varphi)
\]

answer different questions.

The first asks whether the proposition is true.

The second asks how this occurrence of the proposition was derived.

Consequently,

\[
  \Truth(\varphi)
  \neq
  \Prov(\widehat\varphi)
\]

not merely in value but in type.

More importantly, a change in one need not imply a change in the other.

\begin{proposition}[Truth-provenance independence]
\label{prop:truth-provenance-independence}
Two occurrences may have identical propositional content and therefore the
same truth value while possessing distinct provenance records.
\end{proposition}

\begin{proof}
Let

\[
  \widehat\varphi_1
  =
  (\varphi,\pi_1)
\]

and

\[
  \widehat\varphi_2
  =
  (\varphi,\pi_2)
\]

with

\[
  \pi_1\neq\pi_2.
\]

Since both project to the same content \(\varphi\),

\[
  \Truth(\Content(\widehat\varphi_1))
  =
  \Truth(\varphi)
  =
  \Truth(\Content(\widehat\varphi_2)).
\]

Yet by construction,

\[
  \Prov(\widehat\varphi_1)
  \neq
  \Prov(\widehat\varphi_2).
\]
\end{proof}

% -----------------------------------------------------------------------------
%  F.7  PROVENANCE AND EVIDENTIAL SUPPORT
% -----------------------------------------------------------------------------

\subsection{Provenance Is Not Evidence Strength}
\label{app:provenance-evidence}

A second distinction is required.

A low generation count does not guarantee that a proposition is well
supported.

A directly inspected source can be wrong.

Likewise, a high generation count does not entail falsehood.

A repeatedly copied mathematical identity may remain exactly correct.

Let

\[
  E(\varphi)
\]

denote evidential support for a proposition and let

\[
  g(\widehat\varphi)
\]

denote generation depth.

There is no universal monotone law

\[
  E(\varphi)=f(g(\widehat\varphi)).
\]

Generation depth may affect reliability under a particular transformation
process.

It is not itself reliability.

This prevents the provenance notation from becoming a disguised hierarchy of
truth.

% -----------------------------------------------------------------------------
%  F.8  KNOWN AND UNKNOWN PROVENANCE
% -----------------------------------------------------------------------------

\subsection{The Epistemic Layer}
\label{app:epistemic-provenance}

So far \(\pi\) has represented actual provenance.

An observer does not necessarily know \(\pi\).

Let

\[
  K_t(\pi)
\]

denote the observer's epistemic state concerning provenance at time \(t\).

For example,

\[
  K_t(\pi)
  =
  \{\pi\}
\]

may represent exact provenance knowledge.

A set

\[
  K_t(\pi)
  =
  \{
    \pi_1,\pi_2,\pi_3
  \}
\]

may represent unresolved alternatives.

A probability distribution

\[
  \mu_t(\pi)
\]

may encode graded uncertainty.

The important distinction is therefore among

\[
  \varphi,
  \qquad
  \pi,
  \qquad
  K_t(\pi).
\]

These are respectively content, derivation, and knowledge of derivation.

% -----------------------------------------------------------------------------
%  F.9  THREE KINDS OF CHANGE
% -----------------------------------------------------------------------------

\subsection{Content Change, History Change, Knowledge Change}
\label{app:three-changes}

Consider a typed occurrence

\[
  \widehat\varphi_t
  =
  \bigl(
    \varphi_t,
    \pi_t,
    K_t(\pi_t)
  \bigr).
\]

Three changes must be distinguished.

A content change occurs when

\[
  \varphi_t\neq\varphi_{t+1}.
\]

A derivational change occurs when a new transformation extends the history:

\[
  \pi_t\neq\pi_{t+1}.
\]

An epistemic provenance change occurs when

\[
  K_t(\pi)\neq K_{t+1}(\pi)
\]

while both

\[
  \varphi_t=\varphi_{t+1}
\]

and

\[
  \pi_t=\pi_{t+1}.
\]

The last case is the one enacted by the wound.

Nothing happened to the proposition.

Nothing happened to its past.

Something happened to what could responsibly be said about that past.

% -----------------------------------------------------------------------------
%  F.10  THE ORDINARY UNCERTAINTY TAG
% -----------------------------------------------------------------------------

\subsection{\texorpdfstring{$[?]$}{[?]} as Unresolved Provenance}
\label{app:ordinary-uncertainty}

The ordinary tag

\[
  [?]
\]

means that the provenance class required by the current discussion has not
been established.

Formally, if

\[
  \Pi
\]

is the relevant provenance space, then

\[
  [?]
\]

indicates that the epistemically admissible set

\[
  K_t(\pi)
  \subseteq
  \Pi
\]

contains more than one live possibility.

The notation makes no claim that the provenance is intrinsically
unrecoverable.

Further archival work may reduce

\[
  K_t(\pi)
\]

to a singleton.

Thus

\[
  [?]
\]

means unresolved.

It does not yet mean lost.

% -----------------------------------------------------------------------------
%  F.11  THE OMEGA QUESTION
% -----------------------------------------------------------------------------

\subsection{What \texorpdfstring{$[\omega?]$}{[omega?]} Means}
\label{app:omega-question}

The exceptional tag

\[
  [\omega?]
\]

marks a stronger condition.

It does not mean

\[
  g(\gamma)=\omega.
\]

Nor does it mean that infinitely many transformations literally occurred.

The symbol \(\omega\) is used here as a visual marker for unbounded or
unrecoverably open derivational depth.

Let

\[
  \Pi_{\mathrm{recoverable}}(E)
\]

be the provenance distinctions recoverable from the surviving evidence \(E\).

Then

\[
  [\omega?]
\]

marks a case in which the provenance coordinate required by the argument is
not determined by \(E\), and no currently available derivation establishes a
finite exact generation count or transformation word.

Schematically,

\[
  K_t(g)
  =
  \{
    n\in\mathbb N:
    n\text{ remains compatible with }E
  \}
\]

may be unbounded.

The tag therefore means not \enquote{infinite history} but
\enquote{derivational depth no longer responsibly bounded by the available
record}.

% -----------------------------------------------------------------------------
%  F.12  WHY THE WOUND IS UNIQUE
% -----------------------------------------------------------------------------

\subsection{Scarcity as Semantics}
\label{app:wound-scarcity}

The preamble enforces that

\[
  \provomegadisplay
\]

may occur only once.

This is not merely a typographic joke.

The ordinary uncertainty tag already handles routine provenance ambiguity.

If every unresolved claim received

\[
  [\omega?],
\]

the symbol would cease to distinguish ordinary uncertainty from a specific
failure of provenance recovery.

Its scarcity is therefore part of its semantics.

The compile-time constraint

\[
  N_{\omega?}\leq1
\]

acts as a tiny formal discipline imposed upon the document itself.

The source file refuses a second use unless the author deliberately changes
the rule.

% -----------------------------------------------------------------------------
%  F.13  THE PLANTED PROPOSITION
% -----------------------------------------------------------------------------

\subsection{The Wound as an Epistemic Event}
\label{app:planted-wound}

Let the proposition planted in Movement~I be

\[
  \varphi.
\]

At its first occurrence, suppose the essay represented its provenance as

\[
  \widehat\varphi_{t_0}
  =
  \varphi^{[0]}.
\]

The reader was therefore licensed to infer

\[
  K_{t_0}(g)=\{0\}.
\]

Movement~IV revises that status.

The proposition's wording remains

\[
  \varphi.
\]

Its truth conditions remain those of \(\varphi\).

Its actual historical derivation, whatever it was, has not retroactively
changed.

What changes is the warranted provenance description:

\[
  K_{t_1}(g)
  \neq
  K_{t_0}(g).
\]

If the surviving derivational record is insufficient to establish any bounded
exact depth, the annotation becomes

\[
  \varphi^{[\omega?]}.
\]

Hence

\[
  \Content_{t_0}(\varphi)
  =
  \Content_{t_1}(\varphi),
\]

while

\[
  K_{t_0}(\Prov(\varphi))
  \neq
  K_{t_1}(\Prov(\varphi)).
\]

That is the wound.

% -----------------------------------------------------------------------------
%  F.14  THE WOUND DOES NOT RETROACTIVELY ALTER HISTORY
% -----------------------------------------------------------------------------

\subsection{Discovery Is Not Derivation}
\label{app:discovery-not-derivation}

It would be a mistake to say that the proposition's provenance changed when
the essay discovered its uncertainty.

Let the actual provenance be

\[
  \pi^\ast.
\]

Then before and after the discovery,

\[
  \pi^\ast_{t_0}
  =
  \pi^\ast_{t_1}.
\]

Only the epistemic representation changes:

\[
  K_{t_0}(\pi^\ast)
  \neq
  K_{t_1}(\pi^\ast).
\]

This gives the precise distinction requested by the main argument:

\[
  \text{history}
  \neq
  \text{knowledge of history}.
\]

An Originist must track both.

Otherwise uncertainty about provenance can be mistaken for instability in the
past itself.

% -----------------------------------------------------------------------------
%  F.15  PROVENANCE REVISION
% -----------------------------------------------------------------------------

\subsection{Revision Without Semantic Revision}
\label{app:provenance-revision}

Define a provenance-revision operator

\[
  \mathcal U_E
\]

acting on epistemic provenance states.

Given new evidence \(E'\),

\[
  \mathcal U_{E'}:
  K_t(\pi)
  \mapsto
  K_{t+1}(\pi).
\]

A pure provenance revision satisfies

\[
  \Content_{t+1}
  =
  \Content_t
\]

while

\[
  K_{t+1}(\pi)
  \neq
  K_t(\pi).
\]

Thus a document may require revision even when none of its propositional
sentences requires rewriting.

The revision concerns what the document claims to know about its own
derivation.

This is one reason provenance cannot be reduced to bibliography.

% -----------------------------------------------------------------------------
%  F.16  CONFIDENCE ABOUT CONTENT AND CONFIDENCE ABOUT PROVENANCE
% -----------------------------------------------------------------------------

\subsection{Two Confidence Coordinates}
\label{app:two-confidences}

The fuller notation can be refined by separating two probabilities.

Let

\[
  p_T
\]

represent confidence assigned to the truth or correctness of the content under
a specified epistemic model.

Let

\[
  p_P
\]

represent confidence assigned to the provenance reconstruction.

Then a fully explicit typed claim could be written

\[
  \varphi^{
    [3:R\circ C\circ R;
      \mathcal I^+;
      p_T=.94;
      p_P=.71]
  }.
\]

The main text suppresses this distinction unless needed, because the notation
would otherwise become unreadable.

Conceptually, however,

\[
  p_T\neq p_P
\]

must remain possible.

One may be highly confident that a proposition is correct while uncertain
about where it came from.

Conversely, one may know exactly where a false proposition came from.

% -----------------------------------------------------------------------------
%  F.17  FALSEHOOD WITH PERFECT PROVENANCE
% -----------------------------------------------------------------------------

\subsection{A Useful Counterexample}
\label{app:false-perfect-provenance}

Suppose a source \(S\) states a false proposition

\[
  \psi.
\]

The source is inspected directly.

Then the occurrence may correctly carry

\[
  \psi^{[0]}.
\]

Its provenance status is excellent.

Its truth status is not.

Thus

\[
  \Prov(\psi^{[0]})
\]

can be maximally well established while

\[
  \Truth(\psi)=0.
\]

This counterexample prevents the entire provenance apparatus from being
misread as a surrogate truth calculus.

% -----------------------------------------------------------------------------
%  F.18  TRUE CONTENT WITH DEGRADED PROVENANCE
% -----------------------------------------------------------------------------

\subsection{The Converse Counterexample}
\label{app:true-degraded-provenance}

Now suppose a true proposition

\[
  \chi
\]

passes through a chain

\[
  \chi_0
  \xrightarrow{R}
  \chi_1
  \xrightarrow{C}
  \chi_2
  \xrightarrow{R}
  \chi_3
\]

without semantic alteration:

\[
  \Content(\chi_0)
  =
  \Content(\chi_1)
  =
  \Content(\chi_2)
  =
  \Content(\chi_3)
  =
  \chi.
\]

Then

\[
  \Truth(\chi)=1
\]

throughout, while the final occurrence has provenance

\[
  \chi^{[3:R\circ C\circ R]}.
\]

Truth has remained fixed.

History has grown.

The two coordinates are therefore demonstrably independent.

% -----------------------------------------------------------------------------
%  F.19  PROVENANCE AS A TYPE SYSTEM
% -----------------------------------------------------------------------------

\subsection{Typing Rules}
\label{app:provenance-type-system}

The notation may be understood as a lightweight type system.

Write

\[
  \Gamma
  \vdash
  \varphi:\pi
\]

to mean that under provenance context \(\Gamma\), the occurrence of
\(\varphi\) is assigned provenance type \(\pi\).

A direct-source rule has schematic form

\[
  \frac{
    \text{\(\varphi\) inspected in designated source \(S\)}
  }{
    \Gamma
    \vdash
    \varphi:[0]
  }.
\]

A reconstruction rule may take the form

\[
  \frac{
    \Gamma\vdash\varphi:[n:w]
    \qquad
    R(\varphi)=\varphi'
  }{
    \Gamma
    \vdash
    \varphi':[n+1:R\circ w]
  }.
\]

A compression rule is analogous:

\[
  \frac{
    \Gamma\vdash\varphi:[n:w]
    \qquad
    C(\varphi)=\varphi'
  }{
    \Gamma
    \vdash
    \varphi':[n+1:C\circ w]
  }.
\]

These rules propagate derivational information even when

\[
  \varphi'=\varphi.
\]

% -----------------------------------------------------------------------------
%  F.20  THE UNCERTAINTY RULE
% -----------------------------------------------------------------------------

\subsection{Failure of Derivational Typing}
\label{app:uncertainty-rule}

Suppose the available evidence does not determine whether

\[
  \varphi
\]

was copied directly, reconstructed once, or regenerated through several
intermediate states.

Then no exact rule justifies

\[
  \Gamma\vdash\varphi:[n:w].
\]

The responsible judgment is instead

\[
  \Gamma\vdash\varphi:[?].
\]

If the available record additionally fails to bound the relevant derivational
depth, the exceptional judgment is

\[
  \Gamma\vdash\varphi:[\omega?].
\]

This notation records a failure of derivational knowledge rather than filling
the gap with a guessed history.

% -----------------------------------------------------------------------------
%  F.21  SUBSTITUTION FAILURE
% -----------------------------------------------------------------------------

\subsection{Content Equality Does Not License Typed Substitution}
\label{app:typed-substitution}

Suppose

\[
  \varphi^{[0]}
\]

and

\[
  \varphi^{[3:R\circ C\circ R]}
\]

have identical content.

At the untyped semantic level one may substitute one occurrence for the other
where only content matters.

At the provenance-typed level, unrestricted substitution is invalid.

\begin{proposition}[Typed substitution failure]
\label{prop:typed-substitution}
If

\[
  \pi_1\neq\pi_2,
\]

then content equality

\[
  \Content(\varphi,\pi_1)
  =
  \Content(\varphi,\pi_2)
\]

does not imply equality of the typed claims.
\end{proposition}

This is the proposition-level analogue of the failure of state equality to
determine trajectory equality.

% -----------------------------------------------------------------------------
%  F.22  THE FOOTNOTE AS A PROVENANCE CARRIER
% -----------------------------------------------------------------------------

\subsection{Why the Apparatus Is in the Footnotes}
\label{app:footnote-carrier}

The footnote is particularly suited to this experiment because it normally
occupies a subordinate bibliographic role.

Here it carries an additional state variable.

A footnote attached to

\[
  \varphi
\]

may identify a source while the superscript attached to the claim specifies
how the present occurrence relates derivationally to that source.

Thus the bibliography answers

\[
  \text{Which source is relevant?}
\]

while the provenance annotation answers

\[
  \text{What happened between that source and this occurrence?}
\]

These questions overlap.

They are not identical.

% -----------------------------------------------------------------------------
%  F.23  TYPOGRAPHY AS A LOSSY DISPLAY CHANNEL
% -----------------------------------------------------------------------------

\subsection{The Page Does Not Display the Whole Type}
\label{app:typography-lossy}

Even the elaborate notation of this essay does not display complete provenance.

Let

\[
  \pi
\]

be the full provenance record and let

\[
  D_{\mathrm{page}}(\pi)
\]

be the displayed annotation.

Then

\[
  D_{\mathrm{page}}
\]

is itself a quotient.

Many complete provenance histories may produce the same visible tag.

For example,

\[
  D_{\mathrm{page}}(\pi_1)
  =
  D_{\mathrm{page}}(\pi_2)
  =
  [2]
\]

even though

\[
  \pi_1\neq\pi_2.
\]

The essay therefore cannot escape its own Originist objection.

It can only make its chosen losses explicit.

% -----------------------------------------------------------------------------
%  F.24  THE DOCUMENT AS A DERIVED OBJECT
% -----------------------------------------------------------------------------

\subsection{This PDF Has a Trajectory}
\label{app:document-trajectory}

Let

\[
  D_0
\]

be an initial source file and let successive editing, compilation, correction,
and reconstruction operations produce

\[
  D_0
  \xrightarrow{E_1}
  D_1
  \xrightarrow{E_2}
  D_2
  \xrightarrow{K}
  P_1
  \xrightarrow{K}
  P_2,
\]

where \(K\) denotes compilation.

The final PDF \(P_2\) is an endpoint.

It does not contain, merely by existing, the complete derivational chain that
produced it.

Two byte-identical PDFs could have different source histories.

Thus even

\[
  P_A\eqbyte P_B
\]

does not imply

\[
  \Prov(P_A)=\Prov(P_B).
\]

The essay's own artifact therefore instantiates the distinction it argues for.

% -----------------------------------------------------------------------------
%  F.25  TWO-PASS COMPILATION AS A SMALL MODEL
% -----------------------------------------------------------------------------

\subsection{Compilation and Self-Knowledge}
\label{app:two-pass}

The preamble's insistence on multiple compilation passes is intentionally
minor but structurally apt.

On an early pass, the document may not yet possess resolved information about
cross-references or counters whose values depend upon the compiled whole.

After another pass, the visible document can incorporate information generated
by its prior compilation state.

Schematically,

\[
  S_0
  \xrightarrow{K}
  S_1
  \xrightarrow{K}
  S_2.
\]

The final state

\[
  S_2
\]

contains information that depends upon the trajectory

\[
  S_0\to S_1\to S_2.
\]

This does not make LaTeX compilation metaphysically profound.

It makes it an unusually literal local example of a document whose present
display depends upon its own prior states.

% -----------------------------------------------------------------------------
%  F.26  SELF-APPLICATION WITHOUT SELF-ABSORPTION
% -----------------------------------------------------------------------------

\subsection{The Limit of the Device}
\label{app:self-application-limit}

The self-application of provenance notation must remain subordinate to the
argument.

The essay is not primarily about its own footnotes.

Its apparatus serves as a controlled demonstration of a more general claim.

Card's archives concern historical descent \cite{card1989}.

Zebrowski's transformations concern future continuation
\cite{zebrowski1979}.

Le Guin's departures, returns, and dissemination concern recurrence and
branching \cite{leguin1974}.

The document's provenance notation supplies a fourth case at a deliberately
smaller scale.

Its function is evidential and performative, not foundational.

% -----------------------------------------------------------------------------
%  F.27  THE EPISTEMIC WOUND THEOREM
% -----------------------------------------------------------------------------

\subsection{The Formal Payoff}
\label{app:wound-theorem}

\begin{theorem}[Epistemic wound theorem]
\label{thm:epistemic-wound}
Let a proposition occurrence have fixed content \(\varphi\) and fixed actual
provenance \(\pi^\ast\). If new evidence changes the warranted epistemic state
from

\[
  K_0(\pi^\ast)
\]

to

\[
  K_1(\pi^\ast)
\]

with

\[
  K_0(\pi^\ast)\neq K_1(\pi^\ast),
\]

then the epistemic type of the occurrence changes even though neither its
content nor its actual history changes.
\end{theorem}

\begin{proof}
The enriched occurrence is represented by

\[
  \widehat\varphi_t
  =
  \bigl(
    \varphi,
    \pi^\ast,
    K_t(\pi^\ast)
  \bigr).
\]

By hypothesis,

\[
  \varphi_0=\varphi_1
\]

and

\[
  \pi^\ast_0=\pi^\ast_1,
\]

while

\[
  K_0(\pi^\ast)\neq K_1(\pi^\ast).
\]

Therefore

\[
  \widehat\varphi_0
  \neq
  \widehat\varphi_1
\]

at the epistemic-type level despite invariance of content and actual
derivation.
\end{proof}

This is exactly the distinction enacted when a proposition formerly marked

\[
  [0]
\]

must later be marked

\[
  [\omega?].
\]

% -----------------------------------------------------------------------------
%  F.28  THE ORIGINIST RULE OF RESPONSIBLE ANNOTATION
% -----------------------------------------------------------------------------

\subsection{Do Not Invent the Missing Path}
\label{app:responsible-annotation}

The formal apparatus yields a simple methodological rule.

\begin{objection}[Originist Rule of Responsible Annotation]
When the surviving evidence fails to determine a derivational path, a
representation should preserve the uncertainty rather than silently replace
the missing path with the most convenient reconstruction.
\end{objection}

This rule is the epistemic counterpart of preserving branching in the path
category.

If several histories remain admissible, the representation should not collapse
them merely because one endpoint description is easier to print.

% -----------------------------------------------------------------------------
%  F.29  PROVENANCE MONOTONICITY DOES NOT HOLD
% -----------------------------------------------------------------------------

\subsection{Later Knowledge Can Be Worse}
\label{app:nonmonotone-provenance}

One might expect provenance knowledge to improve monotonically as inquiry
continues.

That expectation is false.

New evidence can reveal that an apparently secure derivation was less secure
than previously believed.

Thus an epistemic sequence may move from

\[
  [0]
\]

to

\[
  [?]
\]

or even

\[
  [\omega?].
\]

This is not a contradiction.

The earlier tag represented an earlier epistemic state.

The later tag represents a revision of what was warranted.

Provenance inquiry is therefore itself historical.

Its judgments have trajectories.

% -----------------------------------------------------------------------------
%  F.30  SECOND-ORDER PROVENANCE
% -----------------------------------------------------------------------------

\subsection{The Provenance of Provenance Claims}
\label{app:second-order-provenance}

Once provenance claims are treated as claims, they too possess provenance.

Let

\[
  \rho
  =
  \enquote{the occurrence of \(\varphi\) is generation zero}.
\]

Then \(\rho\) has its own derivational history.

One may therefore write schematically

\[
  \rho^{[1]}
  :
  \varphi^{[0]}.
\]

The outer annotation concerns the provenance of the provenance judgment.

The inner annotation concerns the provenance attributed to \(\varphi\).

This recursion can continue.

There is no practical reason to display it indefinitely.

The important fact is that provenance annotation does not create an
epistemically privileged metalanguage immune to the theory.

% -----------------------------------------------------------------------------
%  F.31  STOPPING THE REGRESS
% -----------------------------------------------------------------------------

\subsection{Operational Closure}
\label{app:provenance-regress}

The regress of provenance-about-provenance is stopped operationally rather
than metaphysically.

For a task \(\tau\), retain provenance information only until additional
provenance distinctions cease to alter the admissible judgments required by
\(\tau\).

If

\[
  P_k
\]

denotes provenance retained to order \(k\), choose \(k\) such that

\[
  J_\tau(P_k)
  =
  J_\tau(P_{k+1})
\]

for all distinctions relevant to the task.

This is the same closure principle governing the essay as a whole.

One stops not because every possible distinction has been exhausted, but
because further extension no longer changes the structure needed by the
argument.

% -----------------------------------------------------------------------------
%  F.32  THE TYPOGRAPHIC THESIS
% -----------------------------------------------------------------------------

\subsection{What the Notation Actually Claims}
\label{app:typographic-thesis}

The provenance typography can now be stated without metaphor.

The expressions

\[
  \varphi^{[0]}
\]

and

\[
  \varphi^{[3:R\circ C\circ R]}
\]

may have identical propositional content.

They are nevertheless different typed occurrences because their derivational
coordinates differ.

Likewise,

\[
  \varphi^{[?]}
\]

and

\[
  \varphi^{[\omega?]}
\]

distinguish ordinary unresolved provenance from a stronger failure to bound or
recover the derivational depth relevant to the inquiry.

Typography is therefore functioning as a visible projection of a provenance
type system.

It does not make the type system complete.

It makes one otherwise invisible coordinate legible.

% -----------------------------------------------------------------------------
%  F.33  CLOSURE
% -----------------------------------------------------------------------------

\subsection{The Wound and the Thesis}
\label{app:wound-closure}

The planted wound now returns the essay to its simplest distinction.

Suppose the sentence is unchanged.

Suppose its truth value is unchanged.

Suppose even its actual derivational history is unchanged.

The reader can nevertheless learn that the earlier representation of that
history was unjustified.

Nothing about the terminal string reveals this.

Only the document's trajectory does.

\begin{woundbox}
\centering
{\color{woundred}\bfseries
The proposition did not change.\\[3pt]
Its history did not change.\\[3pt]
What changed was whether the history could still be known.}
\end{woundbox}

The distinction among

\[
  \Content,
  \qquad
  \Prov,
  \qquad
  K(\Prov)
\]

therefore completes the provenance apparatus required by the essay.

It also prevents the slogan

\[
  \text{state}\neq\text{history}
\]

from becoming too simple.

There are at least three layers:

\[
  \text{state},
  \qquad
  \text{history},
  \qquad
  \text{knowledge of history}.
\]

A representation may preserve the first while losing the second.

An archive may preserve the second while an observer loses access to it.

A reconstruction may recover enough of the second to support one judgment but
not another.

And a proposition may remain word-for-word unchanged while moving from
apparent certainty to an explicit wound.

\begin{held}[Provenance-Type Thesis]
\centering
What a statement says, how this occurrence came to say it,\\
and what we know about how it came to say it\\
are three different objects.
\end{held}

This appendix completes the formal apparatus.

What remains is not another extension of the stack but its bibliographic
grounding: the primary literary works and the secondary sources required by
the argument. The references that follow should therefore remain deliberately
ordinary. After a document devoted to the difference between content and
derivation, the bibliography need not perform another trick.

\newpage

% ============================================================================
%  REFERENCES
%  Traditional thebibliography environment.
%
%  The bibliography is intentionally ordinary. Provenance annotations belong
%  to occurrences in the argument; bibliographic entries identify works.
%  Those functions should not be collapsed at the final stage.
% ============================================================================

\clearpage
\section*{References}
\addcontentsline{toc}{section}{References}
\markboth{References}{References}

\begin{thebibliography}{99}

\bibitem{card1989}
Card, Orson Scott.
``The Originist.''
In \emph{Foundation's Friends: Stories in Honor of Isaac Asimov},
edited by Martin H. Greenberg.
New York: Tor Books, 1989.

\bibitem{zebrowski1979}
Zebrowski, George.
\emph{Macrolife}.
New York: Harper \& Row, 1979.

\bibitem{leguin1974}
Le Guin, Ursula K.
\emph{The Dispossessed: An Ambiguous Utopia}.
New York: Harper \& Row, 1974.

\bibitem{asimov1951}
Asimov, Isaac.
\emph{Foundation}.
New York: Gnome Press, 1951.

\bibitem{asimov1952}
Asimov, Isaac.
\emph{Foundation and Empire}.
New York: Gnome Press, 1952.

\bibitem{asimov1953}
Asimov, Isaac.
\emph{Second Foundation}.
New York: Gnome Press, 1953.

\bibitem{parfit1984}
Parfit, Derek.
\emph{Reasons and Persons}.
Oxford: Clarendon Press, 1984.

\bibitem{williams1970}
Williams, Bernard.
``The Self and the Future.''
\emph{The Philosophical Review}
79, no.~2 (1970): 161--180.

\bibitem{shoemaker1963}
Shoemaker, Sydney.
\emph{Self-Knowledge and Self-Identity}.
Ithaca, NY: Cornell University Press, 1963.

\bibitem{locke1975}
Locke, John.
\emph{An Essay Concerning Human Understanding}.
Edited by Peter H. Nidditch.
Oxford: Clarendon Press, 1975.
Originally published 1689.

\bibitem{butler1736}
Butler, Joseph.
``Of Personal Identity.''
In \emph{The Analogy of Religion, Natural and Revealed, to the Constitution
and Course of Nature}.
London: James, John and Paul Knapton, 1736.

\bibitem{reid1785}
Reid, Thomas.
\emph{Essays on the Intellectual Powers of Man}.
Edinburgh: John Bell, 1785.

\bibitem{wiener1948}
Wiener, Norbert.
\emph{Cybernetics: Or Control and Communication in the Animal and the
Machine}.
Cambridge, MA: MIT Press, 1948.

\bibitem{shannon1948}
Shannon, Claude E.
``A Mathematical Theory of Communication.''
\emph{Bell System Technical Journal}
27 (1948): 379--423, 623--656.

\bibitem{ashby1956}
Ashby, W. Ross.
\emph{An Introduction to Cybernetics}.
London: Chapman \& Hall, 1956.

\bibitem{bateson1972}
Bateson, Gregory.
\emph{Steps to an Ecology of Mind}.
San Francisco: Chandler Publishing Company, 1972.

\bibitem{vonfoerster2003}
von Foerster, Heinz.
\emph{Understanding Understanding: Essays on Cybernetics and Cognition}.
New York: Springer, 2003.

\bibitem{maturana1980}
Maturana, Humberto R., and Francisco J. Varela.
\emph{Autopoiesis and Cognition: The Realization of the Living}.
Dordrecht: D. Reidel, 1980.

\bibitem{varela1974}
Varela, Francisco G., Humberto R. Maturana, and Ricardo Uribe.
``Autopoiesis: The Organization of Living Systems, Its Characterization and
a Model.''
\emph{BioSystems}
5, no.~4 (1974): 187--196.

\bibitem{rosen1991}
Rosen, Robert.
\emph{Life Itself: A Comprehensive Inquiry into the Nature, Origin, and
Fabrication of Life}.
New York: Columbia University Press, 1991.

\bibitem{juarrero1999}
Juarrero, Alicia.
\emph{Dynamics in Action: Intentional Behavior as a Complex System}.
Cambridge, MA: MIT Press, 1999.

\bibitem{juarrero2023}
Juarrero, Alicia.
\emph{Context Changes Everything: How Constraints Create Coherence}.
Cambridge, MA: MIT Press, 2023.

\bibitem{prigogine1984}
Prigogine, Ilya, and Isabelle Stengers.
\emph{Order Out of Chaos: Man's New Dialogue with Nature}.
New York: Bantam Books, 1984.

\bibitem{schrodinger1944}
Schr\"odinger, Erwin.
\emph{What Is Life? The Physical Aspect of the Living Cell}.
Cambridge: Cambridge University Press, 1944.

\bibitem{landauer1961}
Landauer, Rolf.
``Irreversibility and Heat Generation in the Computing Process.''
\emph{IBM Journal of Research and Development}
5, no.~3 (1961): 183--191.

\bibitem{bennett1973}
Bennett, Charles H.
``Logical Reversibility of Computation.''
\emph{IBM Journal of Research and Development}
17, no.~6 (1973): 525--532.

\bibitem{kolmogorov1965}
Kolmogorov, A. N.
``Three Approaches to the Quantitative Definition of Information.''
\emph{Problems of Information Transmission}
1, no.~1 (1965): 1--7.

\bibitem{cover2006}
Cover, Thomas M., and Joy A. Thomas.
\emph{Elements of Information Theory}.
2nd ed.
Hoboken, NJ: Wiley-Interscience, 2006.

\bibitem{mackay2003}
MacKay, David J. C.
\emph{Information Theory, Inference, and Learning Algorithms}.
Cambridge: Cambridge University Press, 2003.

\bibitem{jaynes2003}
Jaynes, E. T.
\emph{Probability Theory: The Logic of Science}.
Edited by G. Larry Bretthorst.
Cambridge: Cambridge University Press, 2003.

\bibitem{fisher1922}
Fisher, R. A.
``On the Mathematical Foundations of Theoretical Statistics.''
\emph{Philosophical Transactions of the Royal Society of London.
Series A}
222 (1922): 309--368.

\bibitem{amari2016}
Amari, Shun-ichi.
\emph{Information Geometry and Its Applications}.
Tokyo: Springer, 2016.

\bibitem{maclane1998}
Mac Lane, Saunders.
\emph{Categories for the Working Mathematician}.
2nd ed.
New York: Springer, 1998.

\bibitem{riehl2016}
Riehl, Emily.
\emph{Category Theory in Context}.
Mineola, NY: Dover Publications, 2016.

\bibitem{awodey2010}
Awodey, Steve.
\emph{Category Theory}.
2nd ed.
Oxford: Oxford University Press, 2010.

\bibitem{gelfand1994}
Gelfand, Sergei I., and Yuri I. Manin.
\emph{Homological Algebra}.
Berlin: Springer, 1994.

\bibitem{bott1982}
Bott, Raoul, and Loring W. Tu.
\emph{Differential Forms in Algebraic Topology}.
New York: Springer, 1982.

\bibitem{bredon1997}
Bredon, Glen E.
\emph{Sheaf Theory}.
2nd ed.
New York: Springer, 1997.

\bibitem{kashiwara2006}
Kashiwara, Masaki, and Pierre Schapira.
\emph{Categories and Sheaves}.
Berlin: Springer, 2006.

\bibitem{maclane1992}
Mac Lane, Saunders, and Ieke Moerdijk.
\emph{Sheaves in Geometry and Logic: A First Introduction to Topos Theory}.
New York: Springer, 1992.

\bibitem{hatcher2002}
Hatcher, Allen.
\emph{Algebraic Topology}.
Cambridge: Cambridge University Press, 2002.

\bibitem{nlabgroupoid}
Baez, John C., and James Dolan.
``From Finite Sets to Feynman Diagrams.''
In \emph{Mathematics Unlimited---2001 and Beyond},
edited by Bj\"orn Engquist and Wilfried Schmid, 29--50.
Berlin: Springer, 2001.

\bibitem{vanfraassen1980}
van Fraassen, Bas C.
\emph{The Scientific Image}.
Oxford: Clarendon Press, 1980.

\bibitem{quine1960}
Quine, W. V. O.
\emph{Word and Object}.
Cambridge, MA: MIT Press, 1960.

\bibitem{goodman1978}
Goodman, Nelson.
\emph{Ways of Worldmaking}.
Indianapolis: Hackett Publishing Company, 1978.

\bibitem{lewis1986}
Lewis, David.
\emph{On the Plurality of Worlds}.
Oxford: Basil Blackwell, 1986.

\bibitem{putnam1981}
Putnam, Hilary.
\emph{Reason, Truth and History}.
Cambridge: Cambridge University Press, 1981.

\bibitem{whitehead1929}
Whitehead, Alfred North.
\emph{Process and Reality: An Essay in Cosmology}.
New York: Macmillan, 1929.

\bibitem{simondon2020}
Simondon, Gilbert.
\emph{Individuation in Light of Notions of Form and Information}.
Translated by Taylor Adkins.
Minneapolis: University of Minnesota Press, 2020.

\bibitem{deleuze1994}
Deleuze, Gilles.
\emph{Difference and Repetition}.
Translated by Paul Patton.
New York: Columbia University Press, 1994.

\bibitem{derrida1996}
Derrida, Jacques.
\emph{Archive Fever: A Freudian Impression}.
Translated by Eric Prenowitz.
Chicago: University of Chicago Press, 1996.

\bibitem{ricoeur2004}
Ricoeur, Paul.
\emph{Memory, History, Forgetting}.
Translated by Kathleen Blamey and David Pellauer.
Chicago: University of Chicago Press, 2004.

\bibitem{assmann2011}
Assmann, Jan.
\emph{Cultural Memory and Early Civilization: Writing, Remembrance, and
Political Imagination}.
Cambridge: Cambridge University Press, 2011.

\bibitem{bowker2005}
Bowker, Geoffrey C.
\emph{Memory Practices in the Sciences}.
Cambridge, MA: MIT Press, 2005.

\bibitem{duranti1997}
Duranti, Luciana.
``The Archival Bond.''
\emph{Archives and Museum Informatics}
11 (1997): 213--218.

\bibitem{yates1966}
Yates, Frances A.
\emph{The Art of Memory}.
Chicago: University of Chicago Press, 1966.

\bibitem{lotman1990}
Lotman, Yuri M.
\emph{Universe of the Mind: A Semiotic Theory of Culture}.
Translated by Ann Shukman.
Bloomington: Indiana University Press, 1990.

\bibitem{kripke1980}
Kripke, Saul A.
\emph{Naming and Necessity}.
Cambridge, MA: Harvard University Press, 1980.

\bibitem{putnam1975}
Putnam, Hilary.
``The Meaning of `Meaning'.''
In \emph{Mind, Language and Reality: Philosophical Papers, Volume 2},
215--271.
Cambridge: Cambridge University Press, 1975.

\bibitem{burge1979}
Burge, Tyler.
``Individualism and the Mental.''
\emph{Midwest Studies in Philosophy}
4, no.~1 (1979): 73--121.

\bibitem{nozick1981}
Nozick, Robert.
\emph{Philosophical Explanations}.
Cambridge, MA: Harvard University Press, 1981.

\bibitem{dennett1991}
Dennett, Daniel C.
\emph{Consciousness Explained}.
Boston: Little, Brown and Company, 1991.

\bibitem{clark1998}
Clark, Andy, and David J. Chalmers.
``The Extended Mind.''
\emph{Analysis}
58, no.~1 (1998): 7--19.

\bibitem{husserl1991}
Husserl, Edmund.
\emph{On the Phenomenology of the Consciousness of Internal Time
(1893--1917)}.
Translated by John Barnett Brough.
Dordrecht: Kluwer Academic Publishers, 1991.

\bibitem{bergson1911}
Bergson, Henri.
\emph{Creative Evolution}.
Translated by Arthur Mitchell.
New York: Henry Holt and Company, 1911.

\bibitem{bergson1910}
Bergson, Henri.
\emph{Time and Free Will: An Essay on the Immediate Data of Consciousness}.
Translated by F. L. Pogson.
London: George Allen, 1910.

\bibitem{priest2006}
Priest, Graham.
\emph{In Contradiction: A Study of the Transconsistent}.
2nd ed.
Oxford: Oxford University Press, 2006.

\bibitem{priest2008}
Priest, Graham.
\emph{An Introduction to Non-Classical Logic: From If to Is}.
2nd ed.
Cambridge: Cambridge University Press, 2008.

\bibitem{belnap1977}
Belnap, Nuel D.
``A Useful Four-Valued Logic.''
In \emph{Modern Uses of Multiple-Valued Logic},
edited by J. Michael Dunn and George Epstein, 8--37.
Dordrecht: D. Reidel, 1977.

\bibitem{barwise1983}
Barwise, Jon, and John Perry.
\emph{Situations and Attitudes}.
Cambridge, MA: MIT Press, 1983.

\bibitem{pearl2009}
Pearl, Judea.
\emph{Causality: Models, Reasoning, and Inference}.
2nd ed.
Cambridge: Cambridge University Press, 2009.

\bibitem{spirtes2000}
Spirtes, Peter, Clark Glymour, and Richard Scheines.
\emph{Causation, Prediction, and Search}.
2nd ed.
Cambridge, MA: MIT Press, 2000.

\bibitem{friston2010}
Friston, Karl.
``The Free-Energy Principle: A Unified Brain Theory?''
\emph{Nature Reviews Neuroscience}
11 (2010): 127--138.

\bibitem{friston2013}
Friston, Karl.
``Life as We Know It.''
\emph{Journal of the Royal Society Interface}
10 (2013): 20130475.

\bibitem{peirce1931}
Peirce, Charles Sanders.
\emph{Collected Papers of Charles Sanders Peirce}.
Edited by Charles Hartshorne, Paul Weiss, and Arthur W. Burks.
Cambridge, MA: Harvard University Press, 1931--1958.

\bibitem{sellars1956}
Sellars, Wilfrid.
``Empiricism and the Philosophy of Mind.''
In \emph{Minnesota Studies in the Philosophy of Science}, vol.~1,
edited by Herbert Feigl and Michael Scriven, 253--329.
Minneapolis: University of Minnesota Press, 1956.

\bibitem{brandom1994}
Brandom, Robert B.
\emph{Making It Explicit: Reasoning, Representing, and Discursive
Commitment}.
Cambridge, MA: Harvard University Press, 1994.

\bibitem{latour2005}
Latour, Bruno.
\emph{Reassembling the Social: An Introduction to Actor-Network-Theory}.
Oxford: Oxford University Press, 2005.

\bibitem{star1989}
Star, Susan Leigh, and James R. Griesemer.
``Institutional Ecology, `Translations' and Boundary Objects: Amateurs and
Professionals in Berkeley's Museum of Vertebrate Zoology, 1907--39.''
\emph{Social Studies of Science}
19, no.~3 (1989): 387--420.

\bibitem{nelson1965}
Nelson, Theodor H.
``Complex Information Processing: A File Structure for the Complex, the
Changing and the Indeterminate.''
In \emph{Proceedings of the 1965 20th National Conference},
84--100.
New York: Association for Computing Machinery, 1965.

\bibitem{bush1945}
Bush, Vannevar.
``As We May Think.''
\emph{The Atlantic Monthly}
176, no.~1 (1945): 101--108.

\bibitem{engelbart1962}
Engelbart, Douglas C.
\emph{Augmenting Human Intellect: A Conceptual Framework}.
AFOSR-3233 Summary Report.
Menlo Park, CA: Stanford Research Institute, 1962.

\bibitem{cerf1974}
Cerf, Vinton G., and Robert E. Kahn.
``A Protocol for Packet Network Intercommunication.''
\emph{IEEE Transactions on Communications}
22, no.~5 (1974): 637--648.

\bibitem{lamport1978}
Lamport, Leslie.
``Time, Clocks, and the Ordering of Events in a Distributed System.''
\emph{Communications of the ACM}
21, no.~7 (1978): 558--565.

\bibitem{gray1978}
Gray, Jim.
``Notes on Data Base Operating Systems.''
In \emph{Operating Systems: An Advanced Course},
edited by R. Bayer, R. M. Graham, and G. Seegm\"uller, 393--481.
Berlin: Springer, 1978.

\bibitem{git2005}
Torvalds, Linus.
``Git: Distributed Version Control System.''
Software project initiated 2005.

\bibitem{merkle1988}
Merkle, Ralph C.
``A Digital Signature Based on a Conventional Encryption Function.''
In \emph{Advances in Cryptology---CRYPTO '87},
edited by Carl Pomerance, 369--378.
Berlin: Springer, 1988.

\bibitem{nakamoto2008}
Nakamoto, Satoshi.
``Bitcoin: A Peer-to-Peer Electronic Cash System.''
2008.

\bibitem{lynch1996}
Lynch, Nancy A.
\emph{Distributed Algorithms}.
San Francisco: Morgan Kaufmann, 1996.

\bibitem{kleppmann2017}
Kleppmann, Martin.
\emph{Designing Data-Intensive Applications}.
Sebastopol, CA: O'Reilly Media, 2017.

\bibitem{crdt2011}
Shapiro, Marc, Nuno Pregui\c{c}a, Carlos Baquero, and Marek Zawirski.
``Conflict-Free Replicated Data Types.''
In \emph{Stabilization, Safety, and Security of Distributed Systems},
386--400.
Berlin: Springer, 2011.

\bibitem{shiptheseus}
Plutarch.
``Theseus.''
In \emph{Plutarch's Lives}.
The traditional Ship of Theseus problem derives from Plutarch's account of
the Athenians preserving Theseus's ship by replacing decayed timbers.

\end{thebibliography}

% ============================================================================
%  DOCUMENT CLOSURE
% ============================================================================

\clearpage
\thispagestyle{empty}

\vspace*{\fill}

\begin{center}

{\small\scshape terminal state}

\vspace{2em}

\[
  \boxed{
    \text{\Large state}
    \neq
    \text{\Large history}
  }
\]

\vspace{2em}

{\small
The inequality is printed in the final state of the document.\\
Its justification is not.
}

\vspace{3em}

{\footnotesize
\textit{Nothing Original Need Survive}\\
Flyxion --- Independent Researcher\\
compiled with \enginename
}

\end{center}

\vspace*{\fill}

% ============================================================================
%  END OF DOCUMENT
% ============================================================================

\end{document}
